% Consolidation des acquis de la méthodologie de l'exercice sur texte — Philosophie Tle Tech — Cours signé OnBuch+
% Programme officiel MINESEC. Compiler avec : tectonic N21-consolidation-acquis-methodologie-exercice-texte.tex
\newcommand{\DOCMATIERE}{Philosophie}
\newcommand{\DOCNIVEAU}{Tle Tech}
\newcommand{\DOCMODULE}{Philosophie – classe de Terminale technique}
\newcommand{\DOCLECON}{N21}
\newcommand{\DOCTITRE}{Consolidation des acquis de la méthodologie de l'exercice sur texte}
% OnBuch+ — préambule « cours signés » ENRICHI (compilé avec tectonic / XeLaTeX)
% Système visuel « Pop » d'OnBuch+ : fond crème (jamais blanc), bordures encre
% épaisses, ombres « dures » décalées, titres Archivo Black en capitales.
% Version MATHÉMATIQUES : graphes (pgfplots/tikz), boîtes pédagogiques
% (définition, propriété, méthode, attention, le savais-tu, exercice résolu,
% exercice, TP, bilan), page de garde et signature OnBuch+ ancrée en bas.
%
% Macros attendues dans main.tex AVANT \input{preamble} :
%   \DOCMATIERE  \DOCNIVEAU  \DOCMODULE  \DOCLECON (numéro)  \DOCTITRE
\documentclass[11pt]{article}

\usepackage{fontspec}
\usepackage[french]{babel}
\usepackage[a4paper,margin=2cm,top=3.4cm,bottom=2.8cm,headheight=1.4cm,footskip=1.3cm]{geometry}
\usepackage{amsmath,amssymb,amsfonts}
\usepackage{mathtools} % \xmapsto, \coloneqq... souvent utilisés par les modèles
\usepackage{cancel} % simplifications barrées (bilans de forces)
% \abs{z} (module, valeur absolue) et \norm{u} (norme), utilisés par les modèles
\providecommand{\abs}{}\let\abs\relax\DeclarePairedDelimiter{\abs}{\lvert}{\rvert}
\providecommand{\norm}{}\let\norm\relax\DeclarePairedDelimiter{\norm}{\lVert}{\rVert}
\usepackage[table]{xcolor} % [table] : fournit aussi \rowcolors (alternance de lignes)
\usepackage{pagecolor}
\usepackage{tikz}
\usetikzlibrary{arrows.meta,positioning,calc,shapes.geometric,shapes.misc,shapes.arrows,decorations.pathmorphing,decorations.pathreplacing,decorations.markings,patterns,shadows,mindmap,trees,fit,backgrounds,matrix,intersections,angles,quotes}
\usepackage{pgfplots}
\usepackage[version=4]{mhchem} % équations nucléaires \ce{^{235}_{92}U}
\usepackage[european,siunitx]{circuitikz} % schémas électriques
\pgfplotsset{compat=1.18}
\usepgfplotslibrary{fillbetween,groupplots} % aires entre courbes (calcul intégral)
\usepackage{enumitem}
\usepackage{fancyhdr}
% [most] charge toutes les bibliothèques tcolorbox (listings, minted, raster,
% documentation...) et gonfle inutilement la mémoire TeX sur un document long
% (~300 boîtes) : on ne charge que ce qui est réellement utilisé.
\usepackage[skins,breakable]{tcolorbox}
\usepackage{graphicx}
\usepackage{colortbl}
\usepackage{booktabs}
\usepackage{tabularx}
\usepackage{multirow}
\usepackage{diagbox} % case d'en-tête diagonale des tableaux à double entrée
% Colonnes extensibles centrée / gauche / droite (C, L, R), souvent utilisées par les modèles
\newcolumntype{C}{>{\centering\arraybackslash}X}
\newcolumntype{L}{>{\raggedright\arraybackslash}X}
\newcolumntype{R}{>{\raggedleft\arraybackslash}X}
\usepackage{array}
\usepackage{multicol}
\usepackage{siunitx}
% parse-numbers=false : accepte aussi \SI{2,5 \times 10^{-3}}{...}
\sisetup{locale=FR,output-decimal-marker={,},group-minimum-digits=4,parse-numbers=false}
\DeclareSIUnit\ohm{\ensuremath{\symup{\Omega}}} % U+2126 absent de la police
\DeclareSIUnit\division{div} % réglages d'oscilloscope (ms/div, V/div)
\DeclareSIUnit\tr{tr} % tours (vitesse de rotation, ex. \SI{1500}{\tr\per\minute})
\DeclareSIUnit\molar{M} % concentration molaire (ex. \SI{3}{\molar}, \SI{1}{\micro\molar})
\DeclareSIUnit\an{an} % année (le modèle confond parfois avec \year, un compteur TeX) — ex. \SI{5}{\kilo\meter\per\an}
\DeclareSIUnit\FCFA{FCFA} % franc CFA (ex. \SI{500}{\FCFA\per\kilo\gram})
\DeclareSIUnit\degreeBrix{\textdegree Brix} % teneur en sucre d'un sirop/jus (réfractométrie)
\DeclareSIUnit\month{mois} % le modèle utilise parfois \month, un compteur TeX, comme unité de durée
\DeclareSIUnit\microliter{\micro\liter} % le modèle écrit parfois \microliter en un seul token
\DeclareSIUnit\million{millions} % dénombrement (ex. \SI{15}{\million\per\mL} spermatozoïdes)
\usepackage{qrcode}
% \degree : commande fréquemment utilisée par les modèles pour un angle ou une
% température en dehors de \SI{}{} (non fournie par les paquets chargés).
\providecommand{\degree}{^{\circ}}
% \qed (fin de démonstration), utilisé par les modèles sans amsthm
\providecommand{\qed}{\hfill$\square$}
% \barre : double barre des valeurs interdites dans un tableau de variations (tkz-tab)
\providecommand{\barre}{\ensuremath{\|}}
% environnement proof (amsthm non chargé) utilisé par les modèles
\ifdefined\proof\else\newenvironment{proof}[1][Démonstration]{\par\noindent\textit{#1.}\ }{\qed\par}\fi
\usepackage{lastpage}
\usepackage{hyperref}
\hypersetup{hidelinks}

% ============================================================
% Palette « Pop » OnBuch+ — mêmes hex que la classe P (plus_ui.dart)
% ============================================================
\definecolor{poporange}{HTML}{F59321}
\definecolor{popdark}{HTML}{E07A0C}
\definecolor{popcream}{HTML}{FFF6E8}
\definecolor{popink}{HTML}{1C1714}
\definecolor{poppurple}{HTML}{7A5AE0}
\definecolor{popgreen}{HTML}{2E9E6A}
\definecolor{poppink}{HTML}{E0457B}
\definecolor{popblue}{HTML}{2F7DE1}
\definecolor{popgold}{HTML}{C98A12}
\definecolor{popmuted}{HTML}{8A7F76}
\definecolor{popline}{HTML}{EADFD2}
\definecolor{popgray}{HTML}{8A7F76}
\colorlet{popgrey}{popgray}
% Alias tolérants (le modèle nomme parfois les couleurs différemment)
\colorlet{popwhite}{white}
\colorlet{popblack}{popink}
\colorlet{popred}{poppink}
\colorlet{popyellow}{popgold}
% Teintes claires (fonds de tableaux/encadrés) — le modèle nomme parfois
% les couleurs avec un suffixe L (ex. popblueL) sans les avoir définies.
\colorlet{poporangeL}{poporange!20}
\colorlet{popdarkL}{popdark!20}
\colorlet{poppurpleL}{poppurple!20}
\colorlet{popgreenL}{popgreen!20}
\colorlet{poppinkL}{poppink!20}
\colorlet{popblueL}{popblue!20}
\colorlet{popgoldL}{popgold!20}
\colorlet{popredL}{popred!20}
\colorlet{popgrayL}{popgray!20}
% Teintes claires pour fonds de tableaux / zones
\colorlet{popblueL}{popblue!10!white}
\colorlet{poporangeL}{poporange!14!white}
\colorlet{popgreenL}{popgreen!12!white}
\colorlet{poppurpleL}{poppurple!12!white}
\colorlet{poppinkL}{poppink!12!white}
\colorlet{popgoldL}{popgold!14!white}

\pagecolor{popcream}

% ============================================================
% Typographie
% ============================================================
\defaultfontfeatures{Ligatures=TeX}
\setmainfont{PlusJakartaSans-Medium.ttf}[Path=../../../fonts/,BoldFont=PlusJakartaSans-Bold.ttf]
% Maths en sans-serif (Fira Math) avec chiffres et lettres droites de Plus
% Jakarta, pour que les formules se fondent dans le texte.
\usepackage[math-style=ISO,bold-style=ISO]{unicode-math}
\setmathfont{FiraMath-Regular.otf}[Path=../../../fonts/]
\setmathfont{PlusJakartaSans-Medium.ttf}[Path=../../../fonts/,range={up/num,up/{Latin,latin},\mathup}]
\usepackage{icomma} % virgule décimale sans espace en mode math
% Caractères mathématiques Unicode absents des polices de texte (Plus Jakarta,
% Archivo Black) : rendus par leur équivalent mathématique, en texte comme en
% formule — sinon ils s'affichent en carré vide (ex. « Arithmétique dans ℤ »).
\usepackage{newunicodechar}
\newunicodechar{ℝ}{\ensuremath{\mathbb{R}}}
\newunicodechar{ℤ}{\ensuremath{\mathbb{Z}}}
\newunicodechar{ℂ}{\ensuremath{\mathbb{C}}}
\newunicodechar{ℕ}{\ensuremath{\mathbb{N}}}
\newunicodechar{ℚ}{\ensuremath{\mathbb{Q}}}
\newunicodechar{𝔻}{\ensuremath{\mathbb{D}}}
\newunicodechar{α}{\ensuremath{\alpha}}
\newunicodechar{β}{\ensuremath{\beta}}
\newunicodechar{γ}{\ensuremath{\gamma}}
\newunicodechar{δ}{\ensuremath{\delta}}
\newunicodechar{ε}{\ensuremath{\varepsilon}}
\newunicodechar{θ}{\ensuremath{\theta}}
\newunicodechar{λ}{\ensuremath{\lambda}}
\newunicodechar{μ}{\ensuremath{\mu}}
\newunicodechar{σ}{\ensuremath{\sigma}}
\newunicodechar{τ}{\ensuremath{\tau}}
\newunicodechar{φ}{\ensuremath{\varphi}}
\newunicodechar{ψ}{\ensuremath{\psi}}
\newunicodechar{ω}{\ensuremath{\omega}}
\newunicodechar{Σ}{\ensuremath{\Sigma}}
% majuscules produites par \MakeUppercase dans les titres (α-aminés -> Α)
\newunicodechar{Α}{\ensuremath{\alpha}}
\newunicodechar{Β}{\ensuremath{\beta}}
\newunicodechar{Γ}{\ensuremath{\Gamma}}
\newunicodechar{Δ}{\ensuremath{\Delta}}
\newunicodechar{Ε}{\ensuremath{\varepsilon}}
\newunicodechar{Θ}{\ensuremath{\Theta}}
\newunicodechar{Λ}{\ensuremath{\Lambda}}
\newunicodechar{Μ}{\ensuremath{\mu}}
\newunicodechar{Τ}{\ensuremath{\tau}}
\newunicodechar{Φ}{\ensuremath{\Phi}}
\newunicodechar{Ψ}{\ensuremath{\Psi}}
\newunicodechar{Ω}{\ensuremath{\Omega}}
\newunicodechar{↦}{\ensuremath{\mapsto}}
\newunicodechar{∈}{\ensuremath{\in}}
\newunicodechar{∉}{\ensuremath{\notin}}
\newunicodechar{∧}{\ensuremath{\wedge}}
\newunicodechar{⇔}{\ensuremath{\Leftrightarrow}}
\newunicodechar{⟺}{\ensuremath{\Longleftrightarrow}}
\newunicodechar{⇒}{\ensuremath{\Rightarrow}}
\newunicodechar{⟹}{\ensuremath{\Longrightarrow}}
\newunicodechar{∘}{\ensuremath{\circ}}
\newunicodechar{∪}{\ensuremath{\cup}}
\newunicodechar{∩}{\ensuremath{\cap}}
\newunicodechar{≡}{\ensuremath{\equiv}}
\newunicodechar{⊂}{\ensuremath{\subset}}
\newunicodechar{⊥}{\ensuremath{\perp}}
\newunicodechar{∃}{\ensuremath{\exists}}
\newunicodechar{∀}{\ensuremath{\forall}}
\newunicodechar{∛}{\ensuremath{\sqrt[3]{\,}}}
\newunicodechar{∜}{\ensuremath{\sqrt[4]{\,}}}
\newunicodechar{⃗}{\ensuremath{{}^{\to}}}
\newunicodechar{ⁿ}{\ifmmode^{n}\else\textsuperscript{\ensuremath{n}}\fi}
\newunicodechar{ˣ}{\ifmmode^{x}\else\textsuperscript{\ensuremath{x}}\fi}
\newunicodechar{ᵇ}{\ifmmode^{b}\else\textsuperscript{\ensuremath{b}}\fi}
\newunicodechar{ʳ}{\ifmmode^{r}\else\textsuperscript{\ensuremath{r}}\fi}
\newunicodechar{ᵘ}{\ifmmode^{u}\else\textsuperscript{\ensuremath{u}}\fi}
\newunicodechar{ʸ}{\ifmmode^{y}\else\textsuperscript{\ensuremath{y}}\fi}
\newunicodechar{ᵃ}{\ifmmode^{a}\else\textsuperscript{\ensuremath{a}}\fi}
\newunicodechar{ᵉ}{\ifmmode^{e}\else\textsuperscript{\ensuremath{e}}\fi}
\newunicodechar{ᵏ}{\ifmmode^{k}\else\textsuperscript{\ensuremath{k}}\fi}
\newunicodechar{ᵖ}{\ifmmode^{p}\else\textsuperscript{\ensuremath{p}}\fi}
\newunicodechar{ᴺ}{\ifmmode^{N}\else\textsuperscript{\ensuremath{N}}\fi}
\newunicodechar{ᶜ}{\ifmmode^{c}\else\textsuperscript{\ensuremath{c}}\fi}
\newunicodechar{ʰ}{\ifmmode^{h}\else\textsuperscript{\ensuremath{h}}\fi}
\newunicodechar{ᵠ}{\ifmmode^{q}\else\textsuperscript{\ensuremath{q}}\fi}
\newunicodechar{ₙ}{\ifmmode_{n}\else\textsubscript{\ensuremath{n}}\fi}
\newunicodechar{ₐ}{\ifmmode_{a}\else\textsubscript{\ensuremath{a}}\fi}
\newunicodechar{ᵢ}{\ifmmode_{i}\else\textsubscript{\ensuremath{i}}\fi}
\newunicodechar{ᵣ}{\ifmmode_{r}\else\textsubscript{\ensuremath{r}}\fi}
\newunicodechar{ₖ}{\ifmmode_{k}\else\textsubscript{\ensuremath{k}}\fi}
\newunicodechar{ᵦ}{\ifmmode_{\beta}\else\textsubscript{\ensuremath{\beta}}\fi}
\newunicodechar{ₓ}{\ifmmode_{x}\else\textsubscript{\ensuremath{x}}\fi}
\newfontfamily\popdisplay{ArchivoBlack-Regular.ttf}[Path=../../../fonts/]
\newfontfamily\poplogo{SpaceGrotesk-Bold.ttf}[Path=../../../fonts/]
\newfontfamily\popsemi{PlusJakartaSans-SemiBold.ttf}[Path=../../../fonts/]
\newfontfamily\popxbold{PlusJakartaSans-ExtraBold.ttf}[Path=../../../fonts/]
\newfontfamily\popxboldls{PlusJakartaSans-ExtraBold.ttf}[Path=../../../fonts/,LetterSpace=3.0]

\setlist[enumerate]{leftmargin=1.7em,itemsep=5pt,topsep=4pt}
\setlist[itemize]{leftmargin=1.7em,itemsep=4pt,topsep=4pt,label={\color{poporange}\textbullet}}
\setlength{\parindent}{0pt}
\setlength{\parskip}{5pt}
\renewcommand{\arraystretch}{1.25}

% pgfplots : style Pop par défaut
\pgfplotsset{
  every axis/.append style={axis line style={popink,line width=1pt},tick style={popink},
    grid style={popline,line width=0.5pt},label style={font=\small},tick label style={font=\footnotesize}},
  popcurve/.style={poporange,line width=1.8pt,no markers,smooth},
}
% Même style disponible dans un \draw TikZ simple (hors pgfplots)
\tikzset{popcurve/.style={poporange,line width=1.8pt}}
\tikzset{popgrid/.style={popline,very thin}}
\colorlet{popcurve}{poporange} % le nom du style est parfois utilisé comme couleur

% ============================================================
% Pill / squiggle
% ============================================================
\newcommand{\poppill}[3]{%
  \tikz[baseline=-0.55ex]\node[draw=popink,line width=1.2pt,rounded corners=2.6pt,%
    fill=#2,inner xsep=6.5pt,inner ysep=3pt]{%
    \popxboldls\fontsize{7}{7}\selectfont\color{#3}\MakeUppercase{#1}};%
}
\newcommand{\popsquiggle}[1]{%
  \tikz[baseline=-0.4ex]{
    \draw[#1,line width=2.6pt,line cap=round]
      plot [smooth] coordinates {(0,0.09) (0.34,0.01) (0.68,0.21) (1.02,0.01) (1.36,0.10)};
  }%
}
% Mot-clé surligné (marqueur orange sous le texte)
\newcommand{\cle}[1]{{\popxbold\color{popdark}#1}}

% ============================================================
% En-tête / pied
% ============================================================
\pagestyle{fancy}
\fancyhf{}
\renewcommand{\headrulewidth}{0pt}
\renewcommand{\footrulewidth}{0pt}
\lhead{\raisebox{-0.15ex}{\poplogo\fontsize{13}{13}\selectfont\color{popink}OnBuch\color{poporange}+}}
\rhead{\poppill{\DOCMATIERE\ \textbullet\ \DOCNIVEAU\ \textbullet\ Leçon \DOCLECON}{white}{popink}}
\lfoot{\popxbold\fontsize{7.6}{7.6}\selectfont\color{popmuted}\MakeUppercase{Cours signé OnBuch+}\ \textbullet\ {\popsemi\DOCTITRE}}
\rfoot{\tikz[baseline=-0.6ex]\node[rounded corners=8.5pt,fill=popink,minimum height=17pt,minimum width=17pt,inner xsep=5pt,inner sep=0]{\popxbold\fontsize{8}{8}\selectfont\color{popcream}\thepage\,/\,\pageref*{LastPage}};}

% ============================================================
% Titres
% ============================================================
\newcommand{\chaptitre}[2]{%
  \begin{tcolorbox}[enhanced,colback=poporange,colframe=popink,boxrule=2.6pt,arc=11pt,
      drop shadow=popink,left=15pt,right=15pt,top=12pt,bottom=11pt,before skip=3pt,after skip=18pt]
    \poppill{#2}{popink}{popcream}\par\vspace{9pt}
    {\raggedright\hyphenpenalty=10000\exhyphenpenalty=10000\popdisplay\fontsize{22}{24}\selectfont\color{popcream}\MakeUppercase{#1}\par}
  \end{tcolorbox}
}
% Section numérotée : pastille orange avec le numéro + titre Archivo
\newcounter{popsec}
\newcommand{\coursec}[1]{%
  \stepcounter{popsec}\setcounter{popsub}{0}%
  \par\vspace{16pt}\noindent
  \begin{minipage}{\linewidth}
  \tikz[baseline=-0.7ex]\node[draw=popink,line width=1.6pt,fill=poporange,rounded corners=4pt,minimum size=22pt,inner sep=0]{\popdisplay\fontsize{12}{12}\selectfont\color{popcream}\thepopsec};%
  \hspace{8pt}{\popdisplay\fontsize{15}{17}\selectfont\color{popink}\MakeUppercase{#1}}\par
  \vspace{4pt}\noindent{\color{poporange}\rule{36pt}{3.6pt}}
  \end{minipage}\par\nopagebreak\vspace{8pt}
}
\newcounter{popsub}
\newcommand{\courssub}[1]{%
  \stepcounter{popsub}%
  \par\vspace{9pt}\noindent{\popxbold\fontsize{11.5}{13}\selectfont\color{popink}{\color{poporange}\thepopsec.\thepopsub}\hspace{6pt}#1}\par\nopagebreak\vspace{4pt}
}

% ============================================================
% Boîtes pédagogiques — même recette : fond blanc, bordure épaisse colorée,
% ombre dure encre, badge plein. Titre optionnel [#1].
% ============================================================
\tcbset{popbase/.style={breakable,enhanced,colback=white,boxrule=2.3pt,arc=9pt,
  shadow={3pt}{-3pt}{0pt}{popink},left=14pt,right=14pt,top=11pt,bottom=11pt,before skip=11pt,after skip=15pt,
  fontupper=\popsemi\fontsize{10.6}{14.5}\selectfont\color{popink},
  fontlower=\popsemi\fontsize{10.6}{14.5}\selectfont\color{popink}}}
% Coche dessinée (le glyphe ✓ n'existe pas dans les polices du document)
\newcommand{\popcheck}[1]{\tikz[baseline=-0.5ex]\draw[#1,line width=1.6pt,line cap=round,line join=round](-0.13,0.0)--(-0.04,-0.09)--(0.13,0.1);}
\providecommand{\coche}{\popcheck{popgreen}}
\providecommand{\resultat}[1]{\par\smallskip\noindent\fcolorbox{popgreen}{popgreenL}{\parbox{\dimexpr\linewidth-2\fboxsep-2\fboxrule\relax}{\textbf{Résultat.} #1}}\par\smallskip}
\newcommand{\popboxhead}[3]{\poppill{#1}{#2}{white}\ifx\relax#3\relax\else\hspace{6pt}{\popxbold\fontsize{10.6}{12}\selectfont\color{popink}#3}\fi\par\vspace{7pt}}

\newtcolorbox{aretenir}[1][]{popbase,colframe=popgreen,before upper={\popboxhead{À retenir}{popgreen}{#1}}}
\newtcolorbox{exemplebox}[1][]{popbase,colframe=poporange,before upper={\popboxhead{Exemple}{poporange}{#1}}}
\newtcolorbox{definition}[1][]{popbase,colframe=popblue,before upper={\popboxhead{Définition}{popblue}{#1}}}
\newtcolorbox{propriete}[1][]{popbase,colframe=poppurple,before upper={\popboxhead{Propriété}{poppurple}{#1}}}
\newtcolorbox{methode}[1][]{popbase,colframe=popgold,before upper={\popboxhead{Méthode}{popgold}{#1}}}
\newtcolorbox{attention}[1][]{popbase,colframe=poppink,before upper={\popboxhead{Attention}{poppink}{#1}}}
\newtcolorbox{experience}[1][]{popbase,colframe=popblue,colback=popblueL,before upper={\popboxhead{Expérience}{popblue}{#1}}}
% « Le savais-tu ? » : carte violette pleine (contexte, culture, Cameroun)
\newtcolorbox{savaistu}[1][]{popbase,colback=poppurple,colframe=popink,
  fontupper=\popsemi\fontsize{10.4}{14.2}\selectfont\color{white},
  before upper={\poppill{Le savais-tu\,?}{white}{poppurple}\ifx\relax#1\relax\else\hspace{6pt}{\popxbold\color{white}#1}\fi\par\vspace{7pt}}}
% Exercice résolu : énoncé en haut, \tcblower puis la solution sur fond crème
\newcounter{popexo}\newcounter{popexr}
\newtcolorbox{exoresolu}[1][]{popbase,colframe=poporange,colbacklower=popcream,
  segmentation style={popink,line width=1.2pt,dashed},
  before upper={\stepcounter{popexr}\popboxhead{Exercice résolu \thepopexr}{poporange}{#1}},
  before lower={\poppill{Solution}{popink}{popcream}\par\vspace{6pt}}}
% Exercice (non résolu en place) — la correction est dans la partie Corrigés
\newtcolorbox{exercice}[1][]{popbase,colframe=popink,
  before upper={\stepcounter{popexo}\popboxhead{Exercice \thepopexo}{popink}{#1}}}
% Corrigé
\newtcolorbox{corrige}[1][]{popbase,colframe=popgreen,colback=popgreenL,
  before upper={\popboxhead{Corrigé}{popgreen}{#1}}}
% Objectifs / prérequis / bilan
\newtcolorbox{objectifs}{popbase,colframe=popink,colback=poporangeL,
  before upper={\popboxhead{Objectifs de la leçon}{popink}{}}}
\newtcolorbox{prerequis}{popbase,colframe=popink,colback=popblueL,
  before upper={\popboxhead{Prérequis}{popblue}{}}}
\newtcolorbox{bilan}{popbase,colframe=popink,colback=white,boxrule=2.8pt,
  before upper={\popboxhead{Fiche bilan}{poporange}{L'essentiel à connaître pour l'examen}}}
% Figure encadrée avec légende
\newtcolorbox{popfigure}[1][]{enhanced,breakable=false,colback=white,colframe=popline,boxrule=1.4pt,arc=8pt,
  left=10pt,right=10pt,top=10pt,bottom=8pt,before skip=10pt,after skip=12pt,halign=center,
  lower separated=false,halign lower=center,
  fontlower=\popsemi\fontsize{9}{11.5}\selectfont\color{popmuted},
  #1}
\newcommand{\legende}[1]{\par\vspace{6pt}{\popsemi\fontsize{9}{11.5}\selectfont\color{popmuted}#1}}

% ============================================================
% Signature OnBuch+
% ============================================================
\newcommand{\poplogoinline}[1]{{\poplogo\fontsize{#1}{#1}\selectfont\color{popink}OnBuch\color{poporange}+}}
\newcommand{\signatureonbuch}{%
  \begin{tcolorbox}[enhanced,colback=popink,colframe=popink,boxrule=2.2pt,arc=10pt,
      drop shadow=poporange,left=14pt,right=14pt,top=10pt,bottom=10pt,before skip=0pt,after skip=0pt]
    \tikz[baseline=-0.7ex]\node[circle,fill=popgreen,draw=popcream,line width=1.3pt,minimum size=22pt,inner sep=0]{\popcheck{white}};%
    \hspace{9pt}\begin{minipage}[c]{0.62\linewidth}
      {\popxbold\fontsize{12}{13}\selectfont\color{popcream}Signé\ \textbullet\ {\poplogo OnBuch\color{poporange}+}}\par\vspace{2pt}
      {\raggedright\popsemi\fontsize{8}{10.5}\selectfont\color{popline}\MakeUppercase{Équipe pédagogique OnBuch+}\\\MakeUppercase{Conforme au programme officiel MINESEC}\par}
    \end{minipage}\hfill
    {\poplogo\fontsize{22}{22}\selectfont\color{popcream}OnBuch\color{poporange}+}
  \end{tcolorbox}%
}

% ============================================================
% Page de garde
% #1 titre · #2 matière · #3 niveau · #4 module · #5 n° leçon · #6 durée ·
% #7 description / objectifs (texte ou itemize)
% ============================================================
\newcommand{\pagegarde}[7]{%
  \thispagestyle{empty}
  \newgeometry{a4paper,margin=1.7cm,top=1.6cm,bottom=1.4cm}%
  \begin{tikzpicture}[remember picture,overlay]
    \fill[poporange!18] ([xshift=-3.2cm,yshift=-3.6cm]current page.north east) circle (4.6cm);
    \fill[poppurple!14] ([xshift=2.4cm,yshift=7.6cm]current page.south west) circle (3.8cm);
  \end{tikzpicture}%
  {\poplogo\fontsize{30}{30}\selectfont\color{popink}OnBuch\color{poporange}+}\hfill
  \raisebox{0.6ex}{\poppill{Cours signé \textbullet\ Premium}{popink}{popcream}}\par
  \vspace{1pt}\popsquiggle{poppurple}\par
  \vspace{8pt}
  \begin{tcolorbox}[enhanced,colback=poporange,colframe=popink,boxrule=2.8pt,arc=13pt,
      drop shadow=popink,left=18pt,right=18pt,top=12pt,bottom=13pt,after skip=10pt]
    \poppill{#2 \textbullet\ #3}{popink}{popcream}\hspace{5pt}\poppill{#4}{white}{popink}\par\vspace{8pt}
    {\popdisplay\fontsize{14}{14}\selectfont\color{popink}LEÇON #5}\par\vspace{4pt}
    {\raggedright\hyphenpenalty=10000\exhyphenpenalty=10000\popdisplay\fontsize{25}{27}\selectfont\color{popcream}\MakeUppercase{#1}\par}
    \vspace{8pt}
    {\popxbold\fontsize{9.5}{9.5}\selectfont\color{popink}\MakeUppercase{Durée conseillée : #6}}
  \end{tcolorbox}
  \begin{tcolorbox}[enhanced,colback=white,colframe=popink,boxrule=2.2pt,arc=10pt,
      drop shadow=popink,left=16pt,right=16pt,top=10pt,bottom=10pt,after skip=8pt]
    \poppill{Dans cette leçon}{poppurple}{white}\par\vspace{5pt}
    \popsemi\fontsize{9.3}{12.6}\selectfont\color{popink}#7
  \end{tcolorbox}
  \noindent\begin{minipage}[t]{0.487\linewidth}
    \begin{tcolorbox}[enhanced,colback=popgreen,colframe=popink,boxrule=2.2pt,arc=10pt,
        drop shadow=popink,left=11pt,right=11pt,top=10pt,bottom=10pt,height=3.1cm,valign=center]
      \tcbox[colback=white,colframe=popink,boxrule=1.3pt,arc=5pt,boxsep=0pt,nobeforeafter,
        left=3pt,right=3pt,top=3pt,bottom=3pt,valign=center]{\qrcode[height=1.9cm]{https://play.google.com/store/apps/details?id=com.luvvix.onbuch&hl=fr}}%
      \hspace{8pt}\begin{minipage}[c]{0.5\linewidth}
        {\popdisplay\fontsize{11}{12.5}\selectfont\color{white}\MakeUppercase{Télécharge OnBuch}}\par\vspace{4pt}
        {\raggedright\popsemi\fontsize{8.4}{11}\selectfont\color{white}Gratuit \textbullet\ Google Play\\Tuteur IA, annales, concours, résultats\par}
      \end{minipage}
    \end{tcolorbox}
  \end{minipage}\hfill
  \begin{minipage}[t]{0.487\linewidth}
    \begin{tcolorbox}[enhanced,colback=poppurple,colframe=popink,boxrule=2.2pt,arc=10pt,
        drop shadow=popink,left=11pt,right=11pt,top=10pt,bottom=10pt,height=3.1cm,valign=center]
      \tikz[baseline=-0.7ex]\node[circle,fill=white,draw=popink,line width=1.3pt,minimum size=1.6cm,inner sep=0]{\popdisplay\fontsize{24}{24}\selectfont\color{poppurple}+};%
      \hspace{8pt}\begin{minipage}[c]{0.6\linewidth}
        {\popdisplay\fontsize{11}{12.5}\selectfont\color{white}\MakeUppercase{Passe à OnBuch+}}\par\vspace{4pt}
        {\raggedright\popsemi\fontsize{8.4}{11}\selectfont\color{white}Tous les cours signés, Tuteur IA illimité, fiches PDF — onglet OnBuch+ de l'app\par}
      \end{minipage}
    \end{tcolorbox}
  \end{minipage}
  \vspace{8pt}
  \popcontenu
  \vfill
  \signatureonbuch
  \restoregeometry
  \newpage
}

% Rangée « Ce cours contient » (page de garde)
\newcommand{\poptile}[3]{% #1 couleur #2 gros texte #3 légende
  \begin{tcolorbox}[enhanced,colback=#1,colframe=popink,boxrule=2pt,arc=9pt,shadow={2.5pt}{-2.5pt}{0pt}{popink},
      left=6pt,right=6pt,top=9pt,bottom=9pt,halign=center,valign=center,height=1.75cm,width=\linewidth,nobeforeafter]
    {\popdisplay\fontsize{12.5}{13}\selectfont\color{white}\MakeUppercase{#2}}\par\vspace{3pt}
    {\centering\popsemi\fontsize{7.8}{9.5}\selectfont\color{white}#3\par}
  \end{tcolorbox}}
\newcommand{\popcontenu}{%
  \par\noindent{\popxboldls\fontsize{7.5}{7.5}\selectfont\color{popmuted}CE COURS SIGNÉ CONTIENT}\par\vspace{7pt}
  \noindent\begin{minipage}[t]{0.235\linewidth}\poptile{popblue}{Cours}{complet, illustré, pas à pas}\end{minipage}\hfill
  \begin{minipage}[t]{0.235\linewidth}\poptile{poporange}{Méthodes}{et exercices résolus}\end{minipage}\hfill
  \begin{minipage}[t]{0.235\linewidth}\poptile{poppink}{Exercices}{type examen + corrigés}\end{minipage}\hfill
  \begin{minipage}[t]{0.235\linewidth}\poptile{popgreen}{Fiche bilan}{l'essentiel à retenir}\end{minipage}\par}

% Sommaire
\newtcolorbox{sommaire}{enhanced,colback=white,colframe=popink,boxrule=2.2pt,arc=10pt,
  drop shadow=popink,left=18pt,right=18pt,top=13pt,bottom=13pt,before skip=6pt,after skip=20pt,
  before upper={\poppill{Sommaire}{poppurple}{white}\par\vspace{9pt}\popsemi\fontsize{10.6}{14.5}\selectfont\color{popink}}}

% Signature de fin de cours — poussée tout en bas de la dernière page
\newcommand{\findecours}{%
  \par\vfill\nopagebreak
  \begin{center}\popsquiggle{poporange}\end{center}\vspace{4pt}
  \signatureonbuch
}
\AtBeginDocument{\def\llbracket{\mathopen{[\mkern-3.5mu[}}\def\rrbracket{\mathclose{]\mkern-3.5mu]}}} % intervalles d'entiers (glyphe ⟦ absent de la police)
% Glyphes absents de Fira Math : redéfinis à partir de glyphes disponibles
\AtBeginDocument{%
  \def\setminus{\mathbin{\backslash}}\let\smallsetminus\setminus
  \def\vdots{\mathord{\vbox{\baselineskip3.2pt\lineskiplimit0pt\kern4pt\hbox{.}\hbox{.}\hbox{.}}}}%
  \def\ddots{\mathinner{\mkern1mu\raise6.5pt\hbox{.}\mkern2mu\raise3.6pt\hbox{.}\mkern2mu\raise0.7pt\hbox{.}\mkern1mu}}%
  \def\triangle{\mathord{\scalebox{0.9}{$\symup{\Delta}$}}}%
  \def\nabla{\mathord{\raisebox{\depth}{\scalebox{1}[-1]{$\symup{\Delta}$}}}}%
  \def\checkmark{\popcheck{popdark}}%
  \def\star{\mathbin{\ast}}\def\bigstar{\mathord{\ast}}%
  \def\diamond{\mathbin{\scalebox{0.6}{$\Diamond$}}}%
  \def\bigcup{\mathop{\mathchoice{\raisebox{-0.3em}{\scalebox{1.7}{$\cup$}}}{\scalebox{1.25}{$\cup$}}{\cup}{\cup}}\displaylimits}%
  \def\bigcap{\mathop{\mathchoice{\raisebox{-0.3em}{\scalebox{1.7}{$\cap$}}}{\scalebox{1.25}{$\cap$}}{\cap}{\cap}}\displaylimits}%
  \def\lesseqqgtr{\mathrel{\lessgtr}}\def\bigoplus{\mathop{\oplus}\displaylimits}%
}
\providecommand{\ssi}{si et seulement si} % abréviation fréquente
\AtBeginDocument{\providecommand{\wideparen}{\overparen}} % arc de cercle

%% FIGFIX-BEGIN v3 — correctifs globaux des figures (docs/onbuch-generation/tools/figfix)
\usepackage{adjustbox}\usepackage{etoolbox}
% toute figure TikZ/pgfplots plus large que la ligne est réduite à la largeur disponible
\BeforeBeginEnvironment{tikzpicture}{\begin{adjustbox}{max width=\linewidth}}
\AfterEndEnvironment{tikzpicture}{\end{adjustbox}}
% légendes pgfplots lisibles par-dessus les courbes
\pgfplotsset{every axis/.append style={legend style={fill=white,fill opacity=0.92,text opacity=1,draw=black!15,inner sep=3pt}}}
% étiquettes d'arêtes (syntaxe edge["..."]) avec fond blanc
\tikzset{every edge quotes/.append style={fill=white,inner sep=1.2pt}}
% bulles de cartes mentales : texte à la ligne dans la bulle, bulle un peu plus grande
\tikzset{every concept/.append style={text width=2.0cm,align=center,minimum size=2.3cm,inner sep=2pt}}
%% FIGFIX-END
\begin{document}

\pagegarde{Consolidation des acquis de la méthodologie de l'exercice sur texte}{Philosophie}{Tle Tech}{Philosophie – classe de Terminale technique}{N21}{2 séances (fiche de progression harmonisée)}{Cette leçon de consolidation revient sur l'ensemble des principes et exigences méthodologiques de l'exercice sur texte philosophique au programme de Terminale technique. À travers des travaux pratiques guidés sur des textes d'auteurs et des situations de la vie courante camerounaise, l'élève s'entraîne à répondre aux questions, à justifier sa démarche et à rédiger des réponses conformes aux attentes du Baccalauréat.
\vspace{4pt}
\begin{multicols}{2}\small
\begin{itemize}
\item À la fin de cette leçon, je sais rappeler les principes fondamentaux de l'exercice sur texte philosophique (fidélité au texte, réponse aux questions posées, pas de paraphrase)
\item À la fin de cette leçon, je sais identifier et dégager la thèse, l'argumentation et les concepts clés d'un texte philosophique
\item À la fin de cette leçon, je sais expliquer un passage ou une expression du texte en m'appuyant sur le contexte de l'œuvre
\item À la fin de cette leçon, je sais rédiger une réponse structurée à une question de compréhension en évitant la paraphrase et le hors-sujet
\item À la fin de cette leçon, je sais discuter une thèse en mobilisant des arguments personnels et des références philosophiques
\item À la fin de cette leçon, je sais appliquer les notions philosophiques étudiées à des situations concrètes de la vie courante au Cameroun
\end{itemize}\end{multicols}}

\begin{sommaire}
\begin{multicols}{2}
\begin{enumerate}
\item Rappel des principes fondamentaux de l'exercice sur texte
\item La lecture méthodique du texte
\item Techniques de réponse aux questions de compréhension et d'analyse
\item Techniques de réponse aux questions de discussion
\item Application des notions philosophiques aux situations concrètes
\item Rédaction, justification de la démarche et gestion de l'épreuve
\item Travaux pratiques de consolidation guidés
\item Activité d'intégration
\item Méthodes et exercices résolus
\item Exercices
\item Corrigés des exercices
\item Fiche bilan
\end{enumerate}
\end{multicols}
\end{sommaire}

\begin{objectifs}
\begin{itemize}
\item À la fin de cette leçon, je sais rappeler les principes fondamentaux de l'exercice sur texte philosophique (fidélité au texte, réponse aux questions posées, pas de paraphrase)
\item À la fin de cette leçon, je sais identifier et dégager la thèse, l'argumentation et les concepts clés d'un texte philosophique
\item À la fin de cette leçon, je sais expliquer un passage ou une expression du texte en m'appuyant sur le contexte de l'œuvre
\item À la fin de cette leçon, je sais rédiger une réponse structurée à une question de compréhension en évitant la paraphrase et le hors-sujet
\item À la fin de cette leçon, je sais discuter une thèse en mobilisant des arguments personnels et des références philosophiques
\item À la fin de cette leçon, je sais appliquer les notions philosophiques étudiées à des situations concrètes de la vie courante au Cameroun
\item À la fin de cette leçon, je sais justifier ma démarche, mes réponses et mes conclusions de manière cohérente
\item À la fin de cette leçon, je sais gérer mon temps et respecter les exigences formelles de l'épreuve du Baccalauréat technique
\end{itemize}
\end{objectifs}

\begin{prerequis}
\begin{itemize}
\item Notion de texte philosophique et distinction texte/dissertation (début d'année de Terminale)
\item Notions de thèse, d'argument et de concept vues dans les unités précédentes
\item Techniques de lecture analytique et de repérage des articulations logiques d'un texte
\item Notions philosophiques des unités de l'année (vérité, devoir, liberté, État, art, technique...)
\item Règles générales de rédaction d'une copie philosophique (introduction, développement, conclusion)
\end{itemize}
\end{prerequis}

\newpage

\coursec{Rappel des principes fondamentaux de l'exercice sur texte}

\courssub{Situation d'accroche et problématique}

Imaginez Amina, élève de Terminale G2 au lycée de Bafoussam, à deux semaines de l'épreuve du Baccalauréat technique. Elle ouvre un sujet zéro : un extrait de \emph{Fondements de la métaphysique des mœurs} de Kant sur l'autonomie de la volonté. La panique la saisit. Face à la première question (« Expliquez la distinction entre autonomie et hétéronomie »), elle recopie machinalement la phrase du texte : « La volonté est une espèce de causalité... ». À la question « Dégagez la thèse du texte », elle écrit : « Kant dit que l'autonomie est le fondement de la morale. » Point final. Pas d'argumentation, pas de mise en relation des concepts, pas d'explication du vocabulaire technique. Son professeur, M. Fouda, corrige sa copie et lui dit : « Amina, l'exercice sur texte n'est pas un résumé, ni une récitation. C'est une \textbf{explication méthodique} guidée par des questions précises. Tu dois \emph{expliquer}, pas \emph{recopier} ; \emph{analyser}, pas \emph{paraphraser} ; \emph{discuter}, pas \emph{opiner} sans arguments. »

\medskip
\noindent \textbf{Problématique :} Comment transformer la lecture d'un texte philosophique en réponses rigoureuses, argumentées et conformes aux exigences du jury du Baccalauréat technique ? Cette leçon de consolidation permet de revoir, pas à pas, toute la méthodologie acquise durant l'année pour aborder l'épreuve avec sérénité et efficacité.

\courssub{Définition et nature de l'exercice sur texte}

\begin{definition}[Exercice sur texte philosophique]
L'exercice sur texte (souvent appelé « commentaire dirigé » ou « explication de texte ») est une épreuve écrite du Baccalauréat technique consistant à \textbf{expliquer un texte philosophique} en répondant à une série de questions posées par l'examinateur. Le texte est imposé (généralement un auteur au programme : Platon, Aristote, Descartes, Kant, Bergson, Sartre, etc., ou un auteur contemporain prolongeant les thèmes du programme). L'élève ne choisit ni le sujet, ni la thèse, ni le plan : il suit le fil conducteur imposé par les questions.
\end{definition}

\begin{aretenir}[Ce que l'exercice sur texte N'EST PAS]
\begin{itemize}
    \item Ce n'est \textbf{pas un résumé} : on ne raconte pas le texte, on en explicite le sens.
    \item Ce n'est \textbf{pas une paraphrase} : remplacer les mots de l'auteur par des synonymes ne prouve pas la compréhension.
    \item Ce n'est \textbf{pas une dissertation} : on ne construit pas son propre plan, on ne choisit pas sa thèse (voir comparatif ci-dessous).
    \item Ce n'est \textbf{pas un commentaire libre} : on répond \emph{uniquement} aux questions posées, dans l'ordre.
\end{itemize}
\end{aretenir}

\begin{popfigure}
\begin{tikzpicture}[
    level 1/.style={sibling distance=4.5cm, level distance=1.8cm},
    level 2/.style={sibling distance=2.2cm, level distance=1.5cm},
    every node/.style={font=\small, align=center},
    box/.style={draw=popblue, thick, rounded corners, fill=popblueL, text width=3.8cm, minimum height=1.2cm},
    boxtwo/.style={draw=poporange, thick, rounded corners, fill=poporangeL, text width=3.5cm, minimum height=1.1cm},
    boxthree/.style={draw=popgreen, thick, rounded corners, fill=popgreenL, text width=3.5cm, minimum height=1.1cm},
    edge from parent/.style={draw=popdark, thick, -{Stealth[length=3mm]}}
]
\node[box, align=center]{\textbf{EXERCICE SUR TEXTE}\\(Commentaire dirigé)}
    child { node[boxtwo, align=center]{\textbf{1. QUESTIONS DE\\COMPRÉHENSION}\\(\og Expliquez \fg, \og Définissez \fg)\\\textit{Expliciter le sens,\\le vocabulaire, les concepts}} }
    child { node[boxtwo, align=center]{\textbf{2. QUESTIONS D'ANALYSE}\\(\og Dégagez la thèse \fg, \og Montrez comment \fg)\\\textit{Reconstituer l'argumentation,\\la structure, le raisonnement}} }
    child { node[boxthree, align=center]{\textbf{3. QUESTIONS DE\\DISCUSSION}\\(\og Discutez \fg, \og Appréciez \fg)\\\textit{Évaluer la portée,\\les limites, confronter}} };
\end{tikzpicture}
\legende{Les trois types de questions classiques de l'exercice sur texte et leurs exigences respectives.}
\end{popfigure}

\courssub{Différence fondamentale : Exercice sur texte vs Dissertation}

C'est la confusion la plus fréquente et la plus pénalisante. Le tableau ci-dessous synthétise les distinctions cruciales pour ne plus jamais les mélanger.

\begin{popfigure}
\begin{center}
\begin{tikzpicture}
\node[inner sep=0pt, align=center] (table){
\begin{tabularx}{\linewidth}{|>{\bfseries}l|X|X|}
\hline
\rowcolor{popblueL}
\textbf{Critère de distinction} & \textbf{Exercice sur texte} & \textbf{Dissertation} \\ \hline
\textbf{Objet} & Un texte philosophique \textbf{imposé} (extrait d'auteur). & Une \textbf{question} ou un \textbf{sujet} philosophique formulé par le jury. \\ \hline
\textbf{Point de départ} & La \textbf{lecture} du texte. Tout part de l'auteur. & La \textbf{problématisation} du sujet. Tout part de l'élève. \\ \hline
\textbf{Liberté du candidat} & \textbf{Nulle} sur le choix du texte, de la thèse, du plan. \textbf{Totale} sur l'explication, l'analyse, la discussion. & \textbf{Totale} : choix de la thèse, du plan, des arguments, des exemples, des auteurs. \\ \hline
\textbf{Structure imposée} & Le \textbf{questionnaire} (généralement 3 questions : compréhension, analyse, discussion). & \textbf{Aucune} : l'élève construit son plan (dialectique, analytique, thématique...). \\ \hline
\textbf{Exigence centrale} & \textbf{Fidélité au texte} : toute affirmation doit s'y appuyer (citation, référence précise). & \textbf{Cohérence et pertinence} : arguments logiques, exemples pertinents, culture philosophique. \\ \hline
\textbf{Rôle de la culture} & \textbf{Secondaire} : elle éclaire le texte (contexte, auteur, concepts), ne le remplace pas. & \textbf{Principale} : elle nourrit l'argumentation (auteurs, concepts, références précises). \\ \hline
\textbf{Piège majeur} & La \textbf{paraphrase}, le \textbf{hors-texte}, la \textbf{recopie}. & Le \textbf{hors-sujet}, le \textbf{catalogage d'auteurs}, le \textbf{plan incohérent}. \\ \hline
\end{tabularx}
};
\draw[popblue, thick] (table.north west) rectangle (table.south east);
\end{tikzpicture}
\end{center}
\legende{Tableau comparatif : Exercice sur texte vs Dissertation.}
\end{popfigure}

\begin{attention}[Piège n°1 : Traiter l'exercice sur texte comme une dissertation]
Ne faites \textbf{jamais} de plan dialectique (Thèse/Antithèse/Synthèse) pour répondre aux questions. Répondez \textbf{question par question}, dans l'ordre, en numérotant vos réponses exactement comme les questions (1, 2, 3...). Une réponse « hors questionnaire » (ex: une introduction générale, une conclusion globale) est sanctionnée.
\end{attention}

\courssub{Le principe de fidélité au texte : la règle d'or}

\begin{definition}[Principe de fidélité]
Toute réponse dans l'exercice sur texte doit \textbf{s'appuyer explicitement sur le texte}. Cela signifie :
\begin{enumerate}
    \item \textbf{Citer} : utiliser des guillemets pour reprendre les termes exacts de l'auteur (mots-clés, phrases courtes).
    \item \textbf{Référencer} : indiquer la ligne ou le paragraphe (ex: « l. 5-6 », « §2 »).
    \item \textbf{Expliquer} : montrer \emph{comment} le passage cité répond à la question.
\end{enumerate}
Une affirmation sans ancrage textuel est une \emph{opinion personnelle} hors sujet, notée 0 ou très faiblement.
\end{definition}

\begin{exemplebox}[Fidélité vs Hors-texte]
\textbf{Texte (Kant) :} « L'autonomie de la volonté est la propriété qu'a la volonté d'être une loi à elle-même... » (l. 3-4).
\medskip
\noindent \textbf{Question :} « Expliquez ce qu'est l'autonomie chez Kant. »
\medskip
\noindent \textbf{Mauvaise réponse (Hors-texte / Opinion) :} « L'autonomie, c'est quand on est libre de choisir sans contrainte extérieure, comme quand je décide de réviser mon bac tout seul. » $\rightarrow$ \textbf{0 point.} (Aucune référence au texte, confusion avec la liberté empirique).
\medskip
\noindent \textbf{Bonne réponse (Fidélité) :} « Chez Kant, l'autonomie est la propriété de la volonté de se donner \textbf{elle-même sa loi} (l. 3-4 : « une loi à elle-même »), indépendamment de tout désir ou inclination extérieure. Ce n'est pas l'indépendance factuelle, mais la législation rationnelle de soi par soi. » $\rightarrow$ \textbf{Point plein.}
\end{exemplebox}

\courssub{L'interdiction de la paraphrase : expliquer $\neq$ reformuler}

\begin{attention}[Erreur fatale : La paraphrase (recopie déguisée)]
Recopier des passages entiers du texte en changeant quelques mots (\emph{« Kant dit que la volonté est une causalité... »}) ne vaut \textbf{zéro point} à la question concernée. L'examinateur cherche la \textbf{compréhension intellectuelle} : la définition des concepts, l'articulation logique, l'implicite rendu explicite.
\end{attention}

\begin{methode}[Comment éviter la paraphrase ? (Technique du \og Déplier \fg)]
\begin{enumerate}
    \item \textbf{Identifiez} le concept-clé ou la phrase cible dans le texte.
    \item \textbf{Citez-le} brièvement (mots-clés entre guillemets).
    \item \textbf{Définissez-le} avec \textbf{VOS mots} (vocabulaire philosophique précis : \cle{maxim}, \cle{impératif catégorique}, \cle{raison pratique}...).
    \item \textbf{Mettez en relation} avec le contexte immédiat (phrase d'avant, phrase d'après) ou la thèse générale.
    \item \textbf{Illustrez} si possible par un exemple concret (situation courante, contre-exemple) pour montrer que vous avez \emph{compris}, pas seulement \emph{lu}.
\end{enumerate}
\end{methode}

\begin{exoresolu}[Paraphrase vs Explication]
\textbf{Texte :} « Agis uniquement d'après la maxime qui fait que tu puisses vouloir en même temps qu'elle devienne une loi universelle. » (Kant, \emph{Fondements}, l. 10-11).
\medskip
\noindent \textbf{Question :} « Expliquez la formule de la loi universelle. »
\medskip
\noindent \tcblower
\noindent \textbf{Paraphrase (0 pt) :} « Kant dit qu'il faut agir seulement selon la maxime qui permet de vouloir qu'elle devienne une loi pour tous. » $\rightarrow$ \textit{Simple reformulation syntaxique.}
\medskip
\noindent \textbf{Explication (Plein pts) :} « La \textbf{maxime} (règle subjective d'action, ex: \og je mens pour sortir d'embarras \fg) doit passer le test de l'\textbf{universalisation} : je dois pouvoir \textbf{vouloir sans contradiction} qu'elle devienne une \textbf{loi universelle} (valable pour tous les êtres rationnels). Si mentir devenait loi universelle, la promesse n'aurait plus de sens (contradiction dans la conception), donc je ne peux pas vouloir cette maxime comme loi. C'est le critère de moralité \emph{a priori}. » $\rightarrow$ \textit{Concepts définis, logique exposée, exemple concret.}
\end{exoresolu}

\courssub{Les trois types de questions : identification et exigences}

Au Baccalauréat technique, le questionnaire suit presque invariablement cette progression logique : \textbf{Comprendre $\rightarrow$ Analyser $\rightarrow$ Discuter}. Savoir identifier le type dès la lecture de l'énoncé permet d'adopter la bonne stratégie immédiatement.

\begin{methode}[Identifier le type de question à sa formulation]
\begin{center}
\begin{tabularx}{\linewidth}{|l|X|X|}
\hline
\rowcolor{popblueL}
\textbf{Mots-clés / Verbes} & \textbf{Type} & \textbf{Exigence centrale} \\ \hline
\og \textbf{Expliquez} \fg, \og \textbf{Définissez} \fg, \og \textbf{Que signifie} \fg, \og \textbf{Précisez le sens de} \fg & \textbf{1. COMPRÉHENSION} & \textbf{Expliciter le sens} : vocabulaire technique, concepts, phrases complexes, distinctions conceptuelles. \textit{Ne pas donner son avis.} \\ \hline
\og \textbf{Dégagez la thèse} \fg, \og \textbf{Quelle est la thèse} \fg, \og \textbf{Montrez comment l'auteur démontre} \fg, \og \textbf{Analysez l'argumentation} \fg, \og \textbf{Relevez les étapes du raisonnement} \fg & \textbf{2. ANALYSE} & \textbf{Reconstituer la logique} : Thèse principale + Arguments (premisses) + Enchaînement (connecteurs logiques : \og car \fg, \og donc \fg, \og or \fg, \og si... alors \fg). Structure du texte. \\ \hline
\og \textbf{Discutez} \fg, \og \textbf{Appréciez} \fg, \og \textbf{Qu'en pensez-vous ?} \fg, \og \textbf{Cette thèse est-elle acceptable ?} \fg, \og \textbf{Les arguments sont-ils convaincants ?} \fg & \textbf{3. DISCUSSION} & \textbf{Évaluer critique} : Portée / Limites / Objections / Contre-exemples / Nuances. \textbf{Obligation :} s'appuyer sur le texte ET la culture philosophique (auteurs, concepts du programme). \textit{Pas d'opinion gratuite.} \\ \hline
\end{tabularx}
\end{center}
\end{methode}

\begin{savaistu}[Le Bac Tech au Cameroun]
Au Baccalauréat technique camerounais, l'épreuve de Philosophie dure \textbf{3 heures}. Elle comporte généralement \textbf{deux parties} :
\begin{itemize}
    \item \textbf{Exercice sur texte} (coefficient souvent 1, sur 10 ou 12 points) : texte d'un auteur au programme (Kant, Bergson, Sartre, Aristote, etc. selon l'année).
    \item \textbf{Dissertation} (coefficient souvent 1, sur 10 ou 12 points) : deux sujets au choix.
\end{itemize}
L'exercice sur texte compte donc pour \textbf{environ 50\% de la note finale}. Le négliger, c'est perdre la moitié des points avant d'avoir commencé la dissertation. La maîtrise de cette méthodologie est donc \textbf{stratégique}.
\end{savaistu}

\courssub{Exigences formelles et barème indicatif}

\begin{propriete}[Exigences formelles obligatoires]
Pour que la copie soit lisible et notée correctement, le candidat \textbf{doit} :
\begin{enumerate}
    \item \textbf{Numérotez} chaque réponse exactement comme la question (1, 2, 3... a, b, c...). Pas de numérotation personnelle.
    \item \textbf{Sautez une ligne} entre chaque réponse numérotée.
    \item \textbf{Rédigez en phrases complètes}, en français correct (orthographe, grammaire, syntaxe). Le style télégraphique ou les listes à puces \emph{à l'intérieur} d'une réponse sont tolérés seulement pour l'énumération d'arguments (analyse), mais le développement doit être phrasé.
    \item \textbf{Citez le texte} avec guillemets et référence de ligne (l. ...) ou paragraphe (§...).
    \item \textbf{Ne pas écrire d'introduction globale} ni de \textbf{conclusion globale}. L'exercice s'arrête à la fin de la dernière question.
\end{enumerate}
\end{propriete}

\begin{aretenir}[Barème indicatif type (sur 12 points) — Variable selon sessions]
\begin{center}
\begin{tabularx}{\linewidth}{|>{\bfseries}X|>{\centering}X|X|}
\hline
\rowcolor{popblueL}
\textbf{Type de question} & \textbf{Points approx.} & \textbf{Critères de réussite} \\ \hline
Compréhension (Q1) & 3--4 pts & Précision des définitions, clarté des distinctions, ancrage textuel rigoureux. \\ \hline
Analyse (Q2) & 4--5 pts & Thèse bien formulée, arguments identifiés et ordonnés, connecteurs logiques explicités, structure du texte respectée. \\ \hline
Discussion (Q3) & 4--5 pts & Objections pertinentes, nuances apportées, exemples philosophiques/concrets, retour au texte pour juger sa solidité, conclusion mesurée. \\ \hline
\textbf{Total} & \textbf{12 pts} & \textbf{+ Points de présentation / langue (souvent 1 à 2 pts bonus ou malus).} \\ \hline
\end{tabularx}
\end{center}
\end{aretenir}

\courssub{Application immédiate : Entraînement à l'identification}

\begin{exercice}[Identification des types de questions]
Pour chacune des questions ci-dessous, indiquez le type (Compréhension / Analyse / Discussion) et justifiez en un mot-clé le verbe opératoire.
\begin{enumerate}
    \item « Expliquez la distinction entre \og nature \fg et \og culture \fg dans le texte. »
    \item « Dégagez la thèse principale de l'auteur et montrez comment il l'établit. »
    \item « L'argument de l'auteur vous semble-t-il convaincant ? Discutez. »
    \item « Que signifie l'expression \og condition de possibilité \fg (l. 4) ? »
    \item « Analysez les étapes du raisonnement qui mènent à la conclusion. »
    \item « Appréciez la portée de cette conception de la liberté. »
\end{enumerate}
\end{exercice}

\begin{corrige}[Exercice n°1]
\begin{enumerate}
    \item \textbf{Compréhension} — Verbe : \og Expliquez \fg (expliciter sens/distinction).
    \item \textbf{Analyse} — Verbes : \og Dégagez \fg (thèse) + \og Montrez comment \fg (argumentation/structure).
    \item \textbf{Discussion} — Verbe : \og Discutez \fg (évaluation critique, avis argumenté).
    \item \textbf{Compréhension} — Verbe : \og Que signifie \fg (définition concept vocabulaire précis).
    \item \textbf{Analyse} — Verbe : \og Analysez \fg (étapes, enchaînement logique, structure).
    \item \textbf{Discussion} — Verbe : \og Appréciez \fg (portée, limites, valeur philosophique).
\end{enumerate}
\end{corrige}

\begin{attention}[Dernier rappel avant la suite]
La \textbf{lecture méthodique} (Section 2) est la condition \emph{sine qua non} de la réussite. On ne répond pas aux questions \emph{pendant} la première lecture. On lit, on annote, on structure, \emph{puis} on rédige. La section suivante détaille ce protocole de lecture active (repérage, vocabulaire, structure, thèse, arguments) qui transforme un texte opaque en matériau exploitable.
\end{attention}

\coursec{La lecture méthodique du texte}

Dans la section précédente, nous avons rappelé les principes fondamentaux de l'exercice sur texte : fidélité à l'auteur, réponse aux questions posées et non à celles qu'on aurait aim\'ees recevoir, distinction entre compr\'ehension et discussion. Mais aucun de ces principes ne peut \^etre appliqu\'e sans une condition pr\'ealable : \cle{bien lire le texte}. Or lire un texte philosophique ne ressemble ni \`a la lecture d'un roman ni \`a celle d'un message sur un t\'el\'ephone. Le texte philosophique est un \emph{raisonnement} : il avance une id\'ee, la justifie, r\'epond \`a des objections. La lecture m\'ethodique est donc l'ensemble des op\'erations ordonn\'ees qui permettent de reconstituer ce raisonnement. C'est l'objet de cette section.

\courssub{Pourquoi une lecture m\'ethodique ?}

Imaginons un \'el\`eve qui, \`a l'\'epreuve, lit le texte une seule fois, rapidement, puis se pr\'ecipite sur les questions. Que risque-t-il ? De r\'epondre \`a c\^ot\'e, d'attribuer \`a l'auteur une id\'ee qu'il combat, de citer un exemple en le prenant pour un argument, ou de paraphraser le texte sans en comprendre la structure. La lecture m\'ethodique \'evite ces \'ecueils. Elle proc\`ede en \cle{\'etapes successives}, de la plus globale \`a la plus fine : d'abord saisir l'ensemble, ensuite analyser le d\'etail, enfin reconstruire le plan du raisonnement.

\begin{popfigure}
\begin{center}
\begin{tikzpicture}[font=\small]
\foreach \i/\txt/\col in {0/{1. Lecture globale : auteur, th\`eme, th\`ese}/popblueL,
                          1/{2. Lecture analytique : concepts, connecteurs}/popgreenL,
                          2/{3. D\'ecoupage : moments de l'argumentation}/poporangeL,
                          3/{4. Plan sch\'ematique au brouillon}/poppinkL}{
  \fill[\col,draw=popdark] (0,-\i*1.1) rectangle (8,-\i*1.1-0.9);
  \node[anchor=west] at (0.2,-\i*1.1-0.45) {\txt};
  \draw[-{Stealth},thick,popdark] (8.3,-\i*1.1-0.45) -- (8.3,-\i*1.1-1.1-0.45);
}
\node[right,popmuted] at (8.5,-2.2) {de l'ensemble vers le d\'etail};
\end{tikzpicture}
\legende{Les quatre \'etapes de la lecture m\'ethodique : on descend marche par marche, du plus g\'en\'eral (qui parle ? de quoi ? pour dire quoi ?) au plus pr\'ecis (comment l'auteur le d\'emontre-t-il ?).}
\end{center}
\end{popfigure}

\courssub{La premi\`ere lecture : lecture globale}

La premi\`ere lecture se fait \emph{stylo pos\'e}, sans s'arr\^eter sur les difficult\'es. Elle vise \`a r\'epondre \`a quatre questions simples.

\textbf{1. Qui parle ?} On rep\`ere l'\cle{auteur} et, si possible, l'\cle{\oe uvre} dont le texte est extrait. Ces indications figurent g\'en\'eralement sous le texte ou dans l'\'enonc\'e. Conna\^itre l'auteur (Platon, Kant, Descartes, Nietzsche, Eboussi Boulaga...) donne des rep\`eres : \'epoque, probl\'ematiques philosophiques habituelles, style. Mais attention : ce rep\'erage sert \`a \emph{orienter} la lecture, non \`a remplacer l'analyse du texte par un cours sur l'auteur.

\textbf{2. De quoi parle-t-on ?} On identifie le \cle{th\`eme g\'en\'eral} : le sujet dont traite le texte (la libert\'e, l'\'Etat, la technique, l'\'education, le bonheur...). Le th\`eme se formule en un ou deux mots, souvent par un nom abstrait.

\textbf{3. Que dit l'auteur sur ce th\`eme ?} On d\'egage la \cle{th\`ese}, c'est-\`a-dire la position pr\'cise que l'auteur d\'efend.

\begin{definition}[La th\`ese]
La \cle{th\`ese} est la position ou l'id\'ee principale que l'auteur affirme et cherche \`a \'etablir dans le texte. Elle r\'epond \`a la question : \og Que veut d\'emontrer l'auteur ? \fg Elle s'exprime toujours par une phrase compl\`ete, jamais par un simple mot.
\end{definition}

\textbf{4. Comment le texte est-il construit ?} On per\c{c}oit globalement le mouvement du texte : l'auteur part-il d'une opinion commune pour la critiquer ? Pose-t-il une question avant d'y r\'epondre ? Cette perception restera provisoire ; la deuxi\`eme lecture la v\'erifiera.

\begin{attention}[Th\`eme ou th\`ese ? Ne pas confondre]
L'erreur la plus fr\'equente au Baccalaur\'eat technique consiste \`a confondre le \cle{th\`eme} (le sujet g\'en\'eral) avec la \cle{th\`ese} (la position pr\'cise de l'auteur sur ce sujet). Le th\`eme tient en un mot ; la th\`ese exige une phrase. Dire \og le texte parle de la libert\'e \fg n'est pas d\'egager une th\`ese : c'est seulement nommer le th\`eme.
\end{attention}

\begin{exemplebox}[Th\`eme et th\`ese dans un texte sur l'\'education]
Soit un texte dont l'auteur critique l'\'education fond\'ee sur la soumission et l'ob\'eissance. Le \textbf{th\`eme} est : \og l'\'education \fg{}. La \textbf{th\`ese} est : \og l'\'education v\'eritable lib\`ere l'homme de la tutelle et le rend capable de penser par lui-m\^eme \fg{}. On voit la diff\'erence : le th\`eme nomme le terrain, la th\`ese indique ce que l'auteur y plante et y d\'efend.
\end{exemplebox}

\courssub{La deuxi\`eme lecture : lecture analytique}

La deuxi\`eme lecture se fait \emph{au crayon}, lentement, phrase par phrase. Trois op\'erations s'y succ\`edent.

\textbf{1. Souligner les concepts cl\'es.} Un \cle{concept} est une id\'ee g\'en\'erale exprim\'ee par un terme pr\'cis (libert\'e, devoir, conscience, \'Etat...). On souligne les termes qui portent la charge philosophique du texte, ainsi que leurs synonymes ou leurs contraires \'eventuels. Si un terme revient souvent, c'est un indice fort : il est sans doute au c\oe ur de la th\`ese.

\textbf{2. Rep\'erer les articulations logiques.} Le raisonnement philosophique est tiss\'e de \cle{connecteurs logiques} qui indiquent le r\^ole de chaque phrase. Les rep\'erer, c'est voir la charpente du texte.

\begin{center}
\begin{tabularx}{\linewidth}{|l|X|}\hline
\rowcolor{popblueL}\textbf{Connecteur} & \textbf{Fonction logique} \\ \hline
car, parce que, en effet, puisque & introduit une \emph{raison}, une justification (un argument) \\ \hline
donc, ainsi, c'est pourquoi, par cons\'equent & introduit une \emph{cons\'equence}, souvent la th\`ese ou une conclusion partielle \\ \hline
cependant, mais, or, n\'eanmoins & introduit une \emph{opposition} : objection, nuance, critique d'une autre position \\ \hline
d'abord, ensuite, enfin & marque les \emph{\'etapes} du raisonnement \\ \hline
par exemple, ainsi le montre & introduit une \emph{illustration} concr\`ete \\ \hline
si... alors, supposons que & introduit une \emph{hypoth\`ese} ou une exp\'erience de pens\'ee \\ \hline
\end{tabularx}
\end{center}

\textbf{3. V\'erifier la th\`ese rep\'er\'ee.} La phrase de la th\`ese se trouve souvent au d\'ebut du texte (annonc\'ee puis d\'emontr\'ee) ou \`a la fin (conclue apr\`es d\'emonstration), parfois signal\'ee par \og donc \fg ou \og c'est pourquoi \fg{}. Mais il arrive qu'elle soit implicite : il faut alors la formuler soi-m\^eme.

\begin{methode}[D\'egager la th\`ese en une seule phrase]
\begin{enumerate}
\item Identifier le th\`eme (le sujet g\'en\'eral du texte).
\item Se demander : \og Que soutient l'auteur \emph{\`a propos} de ce sujet ? Quelle position prend-il ? \fg{}
\item R\'ediger la r\'eponse en \textbf{une seule phrase compl\`ete}, commen\c{c}ant par : \og L'auteur soutient que... \fg ou \og Selon l'auteur, ... \fg{}
\item V\'erifier que cette phrase r\'esume l'ensemble du texte et non un seul passage : si une partie du texte \'echappe \`a la formule, la th\`ese est mal d\'egag\'ee.
\end{enumerate}
\end{methode}

\courssub{Le d\'ecoupage du texte en moments de l'argumentation}

Un texte philosophique n'est pas un bloc compact : il avance par \cle{\'etapes} ou \og moments \fg{}. Le d\'ecoupage consiste \`a diviser le texte en deux, trois ou quatre parties, chacune correspondant \`a une fonction dans le raisonnement. Les sch\'emas les plus fr\'equents sont :

\begin{itemize}
\item \textbf{Th\`ese adverse $\rightarrow$ critique $\rightarrow$ th\`ese de l'auteur} : l'auteur expose d'abord une opinion qu'il r\'efute, puis d\'efend la sienne.
\item \textbf{Question $\rightarrow$ analyse $\rightarrow$ r\'eponse} : l'auteur part d'un probl\`eme et le r\'esout progressivement.
\item \textbf{Th\`ese $\rightarrow$ arguments $\rightarrow$ exemples ou cons\'equences} : l'auteur annonce sa position puis la justifie.
\end{itemize}

Pour d\'ecouper, on s'appuie sur les connecteurs rep\'er\'es (chaque \og cependant \fg ou \og donc \fg marque souvent un changement de moment), sur les changements de sujet de phrase, et sur les alin\'eas \'eventuels. On indique le d\'ecoupage par les num\'eros de lignes : \og de la ligne 1 \`a la ligne 5 : ... \fg{}.

\courssub{Identifier les arguments, les exemples, les objections}

Une fois la th\`ese d\'egag\'ee, il faut identifier ce qui la soutient.

\begin{definition}[L'argument]
Un \cle{argument} est une raison logique avanc\'ee pour \'etablir ou justifier la th\`ese. Il r\'epond \`a la question : \og Pourquoi l'auteur affirme-t-il cela ? \fg Un texte philosophique contient en g\'en\'eral plusieurs arguments, chacun pouvant \^etre reformul\'e en une phrase.
\end{definition}

Il faut distinguer soigneusement trois \'el\'ements que le correcteur attend de voir s\'epar\'es :

\begin{itemize}
\item L'\textbf{argument} : une raison g\'en\'erale (ex. : \og l'\'education d\'eveloppe l'usage de la raison \fg{}).
\item L'\textbf{exemple} : un cas particulier qui illustre l'argument sans le prouver (ex. : \og ainsi, un enfant qui apprend \`a lire acc\`ede seul aux savoirs \fg{}). L'exemple \emph{illustre}, il ne \emph{d\'emontre} pas.
\item L'\textbf{objection} et la \textbf{r\'eponse \`a l'objection} : l'auteur anticipe parfois une critique (\og on objectera que... \fg{}, \og dira-t-on... \fg{}) puis y r\'epond. Rep\'erer ce couple objection/r\'eponse est un signe de lecture fine, tr\`es valoris\'e.
\end{itemize}

\begin{exemplebox}[Texte annot\'e : rep\'erage des connecteurs et de la th\`ese]
\begin{center}
\begin{tikzpicture}[font=\small]
\node[anchor=north west,text width=9.2cm,align=justify] (t) at (0,0) {%
\og On croit souvent que la technique n'est qu'un ensemble d'outils neutres.
\colorbox{poporangeL}{Cependant}~%
tout instrument transforme celui qui s'en sert, \colorbox{poporangeL}{car}~%
il modifie ses gestes, ses habitudes et finalement sa mani\`ere de penser.
\colorbox{poporangeL}{Donc}~%
\colorbox{popgreenL}{la technique n'est jamais neutre : elle façonne l'homme autant que l'homme la façonne}. \fg{}};
\draw[rounded corners,popgreen,thick] (0.05,-4.35) rectangle (9.3,-3.15);
\node[anchor=west,poporange,font=\footnotesize] at (9.6,-0.9) {opposition};
\node[anchor=west,poporange,font=\footnotesize] at (9.6,-2.0) {raison (argument)};
\node[anchor=west,poporange,font=\footnotesize] at (9.6,-3.0) {conclusion};
\node[anchor=west,popgreen,font=\footnotesize\bfseries] at (9.6,-3.8) {THÈSE};
\end{tikzpicture}
\par\smallskip
\textit{Annotation d'un texte court : en orange, les connecteurs logiques ; en vert, la th\`ese. Le d\'ecoupage est visible : opinion commune critiqu\'ee (\og cependant \fg{}) $\rightarrow$ argument (\og car \fg{}) $\rightarrow$ th\`ese conclue (\og donc \fg{}).}
\end{center}
\end{exemplebox}

\courssub{Le plan sch\'ematique au brouillon}

Derni\`ere op\'eration avant de r\'epondre aux questions : r\'ediger au brouillon un \cle{plan sch\'ematique} du texte. Il ne s'agit pas de recopier le texte, mais de r\'esumer chaque moment en une ligne, en indiquant sa fonction logique. Ce plan est la carte qui guidera toutes les r\'eponses : il emp\^eche de paraphraser et permet de situer chaque citation.

\begin{exoresolu}[Construire le plan sch\'ematique d'un texte]
\textbf{\'Enonc\'e.} \`A partir du texte annot\'e ci-dessus sur la technique, r\'edigez le plan sch\'ematique du texte et d\'egagez sa th\`ese en une phrase.
\tcblower
\textbf{Solution d\'etaill\'ee.}
\begin{itemize}
\item \textbf{Th\`eme} : la technique (sa neutralit\'e suppos\'ee).
\item \textbf{Moment 1} (phrase 1) : exposition de l'opinion commune --- la technique serait un outil neutre. Fonction : th\`ese adverse.
\item \textbf{Moment 2} (\og Cependant... car...\fg{}) : critique de cette opinion et argument --- l'instrument transforme les gestes, les habitudes et la pens\'ee de l'utilisateur.
\item \textbf{Moment 3} (\og Donc...\fg{}) : conclusion --- la th\`ese de l'auteur.
\item \textbf{Th\`ese} : \og L'auteur soutient que la technique n'est jamais neutre, car elle façonne l'homme autant que l'homme la façonne. \fg{}
\end{itemize}
On remarque que le plan tient en trois lignes, alors que le texte en compte davantage : c'est bien un \emph{sch\'ema}, non une paraphrase.
\end{exoresolu}

\begin{attention}[Pièges de la lecture m\'ethodique]
\begin{itemize}
\item \textbf{Paraphraser} au lieu d'analyser : r\'ep\'eter le texte avec d'autres mots ne montre pas qu'on l'a compris. Il faut d\'egager la \emph{fonction} de chaque passage (th\`ese, argument, exemple, objection).
\item \textbf{Attribuer \`a l'auteur la th\`ese adverse} : quand un texte commence par \og on croit que... \fg{}, \og certains pensent que... \fg{}, l'id\'ee exprim\'ee est celle que l'auteur \emph{combat}, non celle qu'il d\'efend.
\item \textbf{Prendre l'exemple pour l'argument} : l'exemple illustre, il ne prouve pas.
\item \textbf{S'arr\^eter \`a la premi\`ere difficult\'e} : un mot incompris ne doit pas bloquer la lecture globale ; on y revient \`a la deuxi\`eme lecture, souvent \'eclair\'e par le contexte.
\end{itemize}
\end{attention}

\begin{savaistu}[Philosophie et m\'ethodologie au Cameroun]
L'enseignement de la philosophie dans le secondaire camerounais (programmes MINESEC) accorde une place centrale \`a la ma\^itrise de l'exercice sur texte \`a l'\'ecrit (Probatoire et Baccalaur\'eat technique). Les correcteurs sanctionnent s\'ev\`erement la confusion th\`eme/th\`ese, la paraphrase et l'absence de plan sch\'ematique. La m\'ethodologie pr\'esent\'ee ici est le standard exig\'e pour r\'eussir l'\'epreuve.
\end{savaistu}

\begin{aretenir}[L'essentiel de la lecture m\'ethodique]
\begin{enumerate}
\item \textbf{Lecture globale} : auteur, \oe uvre, th\`eme, premi\`ere id\'ee de la th\`ese.
\item \textbf{Lecture analytique} : concepts cl\'es soulign\'es, connecteurs logiques rep\'er\'es, th\`ese v\'erifi\'ee.
\item \textbf{D\'ecoupage} : le texte se divise en moments (th\`ese adverse, argumentation, conclusion...).
\item \textbf{Distinction} : th\`ese (ce que l'auteur affirme) / arguments (pourquoi) / exemples (illustrations) / objections et r\'eponses.
\item \textbf{Plan sch\'ematique au brouillon} : une ligne par moment, avec sa fonction logique, avant toute r\'eponse aux questions.
\end{enumerate}
\end{aretenir}

La lecture m\'ethodique est ainsi un investissement rentable : dix minutes de lecture rigoureuse \'economisent trente minutes d'h\'esitation devant les questions. C'est pr\'ecis\'ement \`a ces questions --- d'abord celles de compr\'ehension et d'analyse --- que nous consacrerons la section suivante, en nous appuyant sur le plan sch\'ematique que nous venons d'apprendre \`a construire.

\coursec{Techniques de réponse aux questions de compréhension et d'analyse}

La lecture méthodique du texte, étudiée dans la section précédente, nous a permis d'identifier la thèse de l'auteur, ses arguments et la structure de son raisonnement. Mais ce travail préparatoire reste invisible pour le correcteur tant qu'il n'est pas \emph{mis en forme}. C'est précisément l'objet de cette section : transformer notre compréhension en réponses rédigées, conformes aux consignes et dignes du maximum de points. Car à l'examen, ce n'est pas ce que l'on a compris qui est noté, mais ce que l'on a écrit. Nous distinguerons deux grandes familles de questions : les \cle{questions de compréhension} (expliquer) et les \cle{questions d'analyse} (dégager la thèse, l'argumentation, la structure).

\courssub{Les questions de compréhension : expliquer un passage}

\textbf{Intuition.} Imaginons un technicien qui reçoit une machine en panne. Dire « la machine ne marche pas » ne suffit pas : il doit dire \emph{pourquoi} elle ne marche pas, en nommant les pièces défaillantes et leur rôle dans l'ensemble. Expliquer un passage philosophique, c'est exactement cela : ouvrir le texte, nommer ses concepts, montrer comment ils fonctionnent ensemble.

\begin{definition}[Question de compréhension]
Une \cle{question de compréhension} invite le candidat à \emph{expliquer} une expression, une phrase ou un passage du texte, c'est-à-dire à en dégager le sens précis, à en définir les termes et à montrer le rôle que ce passage joue dans l'argumentation de l'auteur. Elle se reconnaît à des consignes comme : « Expliquez... », « Que veut dire l'auteur lorsqu'il affirme... ? », « Dégagez le sens de... ».
\end{definition}

\begin{methode}[Expliquer un passage en quatre étapes]
\begin{enumerate}
\item \textbf{Dégager le sens littéral} : reformuler le passage avec ses propres mots, sans le recopier. On répond à la question : « Que dit l'auteur, simplement ? »
\item \textbf{Définir les concepts} : identifier les termes philosophiques importants du passage et les définir brièvement (liberté, devoir, État, conscience, bonheur...).
\item \textbf{Situer le passage dans l'argumentation} : montrer si le passage énonce la thèse, présente un argument, réfute une objection ou illustre une idée. On répond à la question : « À quoi sert ce passage dans le texte ? »
\item \textbf{Illustrer par un exemple} : donner un exemple concret, de préférence tiré de la vie courante ou du contexte camerounais, qui rend le sens du passage évident.
\end{enumerate}
\end{methode}

\textbf{Mobiliser le contexte de l'œuvre et de la pensée de l'auteur.} Un passage ne se comprend jamais isolément. Si le texte est extrait de la \emph{Lettre sur la tolérance} de Locke, savoir que l'auteur défend la séparation des affaires religieuses et politiques éclaire immédiatement une phrase sur la conscience. De même, une phrase de Nietzsche sur la « mort de Dieu » ne peut être expliquée sans rappeler sa critique de la morale traditionnelle. Le contexte ne doit pas envahir la réponse (une ou deux phrases suffisent), mais il montre au correcteur que le candidat situe le texte dans une pensée d'ensemble.

\begin{attention}[Erreur fréquente : répondre à toutes les questions de la même manière]
Beaucoup de candidats rédigent des réponses identiques quelle que soit la consigne. Or « expliquez », « dégagez la thèse », « dégagez l'argumentation » et « discutez » n'attendent pas le même travail. Une question « expliquez » porte sur un passage précis ; une question « dégagez la thèse » porte sur le texte entier ; une question « discutez » exige une prise de position argumentée. \textbf{Lisez la consigne deux fois} et soulignez le verbe d'action avant de rédiger.
\end{attention}

\begin{exemplebox}[Un exemple chiffré pour calibrer sa réponse]
À l'examen, une question « expliquez » vaut en général \textbf{2 à 3 points}. Le barème répartit ces points entre le sens dégagé, la définition des termes et le rattachement à l'argumentation. Une réponse d'une seule ligne ne peut donc pas prétendre au maximum : elle ne contient matériellement pas assez d'éléments notables. Comptez en moyenne \textbf{une phrase par demi-point} : pour 3 points, visez 5 à 7 phrases construites.
\end{exemplebox}

\courssub{Les questions d'analyse : thèse, argumentation, structure}

\textbf{Intuition.} Si la compréhension zoome sur un passage, l'analyse prend du recul sur le texte entier, comme un géomètre qui, après avoir examiné chaque pierre, décrit le plan de l'édifice.

\begin{definition}[Les trois questions d'analyse]
\begin{itemize}
\item \textbf{Dégager la thèse} : formuler en une ou deux phrases l'idée principale que l'auteur défend, avec ses propres mots, sans paraphrase.
\item \textbf{Dégager l'argumentation} : énumérer les arguments dans l'ordre où ils apparaissent et montrer leur enchaînement logique (comment chaque argument prépare le suivant et conduit à la thèse).
\item \textbf{Dégager la structure} : découper le texte en moments (parties) et indiquer la fonction de chacun (exposition du problème, réfutation, démonstration, conclusion).
\end{itemize}
\end{definition}

\begin{methode}[Répondre à une question d'argumentation]
\begin{enumerate}
\item Relever dans le texte les \textbf{articulations logiques} : « d'abord », « en effet », « cependant », « donc », « par conséquent ». Ce sont les charnières du raisonnement.
\item Numéroter les arguments dans l'\textbf{ordre du texte} (jamais dans le désordre) : « L'auteur soutient d'abord que..., ensuite il montre que..., enfin il conclut que... »
\item Pour chaque argument, indiquer sa \textbf{nature} : argument d'autorité, exemple, raisonnement par l'absurde, analogie, réfutation d'une objection.
\item Montrer l'\textbf{enchaînement} : expliquer pourquoi l'argument 2 suppose l'argument 1 et prépare la conclusion.
\end{enumerate}
\end{methode}

\begin{popfigure}
\begin{center}
\begin{tikzpicture}[font=\small]
% Entonnoir : citation -> sens -> argumentation
\fill[popblue!25] (0,4.4) -- (8,4.4) -- (6.6,3.0) -- (1.4,3.0) -- cycle;
\fill[poporange!30] (1.4,2.8) -- (6.6,2.8) -- (5.6,1.4) -- (2.4,1.4) -- cycle;
\fill[popgreen!30] (2.4,1.2) -- (5.6,1.2) -- (4.7,0) -- (3.3,0) -- cycle;
\node[align=center] at (4,3.7) {\textbf{1. Citation du texte}\\ \footnotesize « ... » (guillemets, bref extrait)};
\node[align=center] at (4,2.1) {\textbf{2. Explication du sens}\\ \footnotesize reformulation + définition des concepts};
\node[align=center] at (4,0.6) {\textbf{3. Rôle dans l'argumentation}\\ \footnotesize thèse, argument, objection ?};
\draw[->,thick,popdark] (4,3.0) -- (4,2.8);
\draw[->,thick,popdark] (4,1.4) -- (4,1.2);
\end{tikzpicture}
\end{center}
\legende{Schéma en entonnoir de la technique d'explication : on part de la citation exacte (1), on en dégage le sens (2), puis on le rattache à l'argumentation globale du texte (3). Chaque étape rétrécit le propos : du texte vers son idée.}
\end{popfigure}

\courssub{La rédaction de la réponse : forme et citation}

Une bonne réponse à une question de compréhension ou d'analyse possède une \textbf{architecture en trois temps} :

\begin{enumerate}
\item \textbf{Une phrase d'introduction de la réponse} qui reprend la consigne et annonce le sens dégagé : « Dans ce passage, l'auteur affirme que... » ou « Par cette expression, l'auteur signifie que... ».
\item \textbf{Un développement appuyé sur le texte} : définition des concepts, explication du raisonnement, rattachement à l'argumentation, exemple.
\item \textbf{Une courte conclusion} : une phrase qui résume le sens et ouvre éventuellement sur la suite du texte : « Ce passage prépare ainsi la réfutation qui suit. »
\end{enumerate}

\begin{propriete}[Règle de la citation]
Pour justifier son interprétation, le candidat \textbf{cite brièvement le texte entre guillemets}, sans jamais le recopier en entier. Une citation efficace compte de 3 à 10 mots et est immédiatement suivie de son explication : « Lorsque l'auteur écrit que "la conscience est inviolable", il entend par là que... ». La citation est une \emph{preuve}, non un remplissage.
\end{propriete}

\begin{exemplebox}[Réponse paraphrasée contre réponse réussie]
Texte (extrait librement adapté) : « L'homme n'est point fait pour la servitude ; obéir à la peur, c'est déjà cesser d'être libre. »

\textbf{Réponse paraphrasée (note faible) :} « L'auteur dit que l'homme n'est pas fait pour la servitude et que quand on obéit par peur on n'est plus libre. » — Cette réponse ne fait que répéter la phrase avec d'autres mots : aucun concept défini, aucun rôle dans l'argumentation, aucun exemple.

\textbf{Réponse réussie (note maximale) :} « Dans ce passage, l'auteur affirme que la liberté est constitutive de la nature humaine et incompatible avec la soumission fondée sur la crainte. La \emph{servitude} désigne ici l'état de celui qui obéit non par adhésion rationnelle mais par \emph{peur} de la sanction. Lorsque l'auteur écrit que "obéir à la peur, c'est déjà cesser d'être libre", il signifie que la liberté réside dans le \emph{motif} de l'obéissance, non dans l'acte lui-même. Ce passage énonce la thèse du texte, que l'auteur défendra ensuite contre l'objection du maintien de l'ordre. Ainsi, un citoyen qui respecte le code de la route par conviction civique reste libre, tandis que celui qui s'y plie uniquement par crainte du gendarme agit déjà en esclave. »
\end{exemplebox}

\begin{exoresolu}[Expliquer un passage : entraînement guidé]
\textbf{Énoncé.} Texte : « La technique, loin d'être une simple servante de la science, est devenue la maîtresse qui façonne notre monde et nos désirs. » Question : Expliquez ce passage (3 points).
\tcblower
\textbf{Solution détaillée.} \emph{Introduction :} Dans ce passage, l'auteur renverse la conception courante selon laquelle la technique ne ferait qu'appliquer les découvertes de la science. \emph{Développement :} La \emph{science} désigne la connaissance désintéressée du réel ; la \emph{technique} désigne l'ensemble des moyens et procédés d'action sur le monde. L'image de la « servante » devenue « maîtresse » signifie un renversement hiérarchique : ce n'est plus la connaissance qui guide l'outil, mais l'outil qui oriente ce que nous cherchons à connaître et même ce que nous désirons. Lorsque l'auteur écrit que la technique « façonne nos désirs », il montre qu'elle ne répond pas seulement à nos besoins : elle les crée. \emph{Exemple :} le téléphone portable, au Cameroun comme ailleurs, n'a pas seulement répondu au besoin de communiquer ; il a fait naître des désirs nouveaux (réseaux sociaux, mobile money, connexion permanente). \emph{Conclusion :} Ce passage énonce la thèse du texte, que l'auteur illustrera par la suite.
\end{exoresolu}

\begin{savaistu}[Le corrigé-type des correcteurs]
Au Baccalauréat technique comme au Probatoire, les correcteurs disposent d'un \emph{corrigé-type} établi lors de réunions d'harmonisation. Ce corrigé liste les \textbf{éléments essentiels attendus} (sens du passage, concepts définis, arguments relevés), mais n'impose pas les mots exacts. Deux réponses rédigées différemment peuvent donc obtenir toutes deux le maximum, pourvu qu'elles contiennent les mêmes éléments de fond. Inutile d'apprendre des formules par cœur : entraînez-vous à dégager les \emph{idées} attendues.
\end{savaistu}

\begin{exercice}[À vous de jouer]
Texte : « On ne gouverne bien les hommes qu'en gouvernant d'abord leurs esprits : la force enchaîne les corps, l'éducation libère les volontés. » 1) Expliquez ce passage (3 points). 2) Dégagez la thèse de l'auteur en une phrase (1 point).
\end{exercice}

\begin{corrige}[Exercice]
1) \emph{Introduction :} L'auteur affirme que le véritable gouvernement politique repose sur la formation des esprits plutôt que sur la contrainte physique. \emph{Développement :} « Gouverner les esprits » signifie éduquer, c'est-à-dire former le jugement des citoyens ; la \emph{force} ne produit qu'une obéissance extérieure (« enchaîne les corps »), tandis que l'\emph{éducation} rend l'adhésion intérieure possible (« libère les volontés »). Le contraste force/éducation structure l'argument : l'obéissance par la contrainte est servitude, l'obéissance par conviction est liberté. \emph{Exemple :} un élève qui respecte le règlement parce qu'il en comprend le sens est libre ; celui qui le respecte par peur de l'exclusion ne fait que subir. \emph{Conclusion :} ce passage énonce la thèse du texte. 2) \emph{Thèse :} le bon gouvernement repose sur l'éducation des citoyens et non sur la force.
\end{corrige}

\begin{aretenir}[L'essentiel de la section]
\begin{itemize}
\item Question de \textbf{compréhension} : expliquer = sens littéral + définition des concepts + rôle dans l'argumentation + exemple.
\item Question d'\textbf{analyse} : thèse (une phrase), argumentation (arguments dans l'ordre, enchaînement logique), structure (moments du texte).
\item Toute réponse s'organise en trois temps : introduction de la réponse, développement appuyé sur le texte, courte conclusion.
\item On \textbf{cite brièvement entre guillemets} pour prouver, jamais pour recopier.
\item Adapter sa réponse à la consigne exacte et calibrer sa longueur sur le barème (environ une phrase par demi-point).
\end{itemize}
\end{aretenir}

\coursec{Techniques de réponse aux questions de discussion}

Dans la section précédente, nous avons appris à répondre aux questions de \cle{compréhension} et d'\cle{analyse} : il s'agissait de rester fidèle au texte, d'expliquer ce que l'auteur dit, comment il le dit et pourquoi il le dit. Mais l'exercice sur texte ne s'arrête pas là. L'examinateur veut aussi savoir ce que \emph{vous}, vous pensez. C'est l'objet des questions de \cle{discussion} : après avoir compris l'auteur, il faut maintenant dialoguer avec lui, l'apprécier, parfois le contredire. C'est le moment où le candidat cesse d'être un commentateur pour devenir un philosophe.

\courssub{Qu'est-ce qu'une question de discussion ?}

\begin{definition}[Question de discussion]
Une \cle{question de discussion} est une question qui invite le candidat à prendre position personnellement sur la thèse de l'auteur : l'\emph{apprécier}, la \emph{discuter}, la \emph{critiquer}, dire si l'on \emph{partage} ou non le point de vue défendu dans le texte.
\end{definition}

On reconnaît facilement ces questions à leurs formulations typiques :
\begin{itemize}
\item « Appréciez la thèse de l'auteur. »
\item « Discutez le point de vue de l'auteur. »
\item « Partagez-vous l'avis de l'auteur ? Justifiez votre réponse. »
\item « La position de l'auteur vous paraît-elle acceptable ? »
\item « Critiquez la conception que l'auteur se fait de la liberté (du devoir, de la technique, de l'État\ldots). »
\end{itemize}

\begin{attention}[Compréhension et discussion : ne pas confondre]
À la question de compréhension, on répond : « l'auteur pense que\ldots ». À la question de discussion, on répond : « \emph{je} pense que\ldots », « à mon avis\ldots », « nous soutiendrons que\ldots ». Employer « l'auteur a raison de dire que\ldots » suivi d'une paraphrase du texte, c'est répondre à côté : on a compris, mais on n'a pas discuté.
\end{attention}

\courssub{Prendre clairement position}

La première exigence d'une discussion réussie est de \cle{prendre position} dès la première phrase. L'examinateur ne doit pas attendre la fin de votre paragraphe pour savoir ce que vous pensez. Trois positions sont possibles :

\begin{enumerate}
\item \textbf{L'accord} : « Nous partageons le point de vue de l'auteur, car\ldots »
\item \textbf{Le désaccord} : « Nous ne pouvons souscrire à la thèse de l'auteur, et ce pour deux raisons\ldots »
\item \textbf{La position nuancée} : « La thèse de l'auteur est dans une large mesure acceptable, car\ldots ; cependant, elle doit être nuancée sur un point\ldots »
\end{enumerate}

\begin{propriete}[Règle d'or de la prise de position]
Toute position est admissible à condition d'être \emph{annoncée clairement}, \emph{justifiée par des arguments personnels} et \emph{maintenue jusqu'à la conclusion}. Il n'existe pas de « bonne » réponse imposée : ce qui est évalué, c'est la qualité de la justification, non le contenu de l'opinion.
\end{propriete}

\begin{attention}[L'erreur la plus fréquente]
Écrire « je suis d'accord avec l'auteur » puis \emph{répéter les arguments du texte} n'est pas discuter : c'est réciter. La discussion exige des arguments \emph{personnels}, c'est-à-dire des raisons que \emph{vous} apportez et qui ne figuraient pas dans le texte. Si l'auteur dit « la technique libère l'homme du travail pénible », vous ne pouvez pas vous contenter de redire cela : vous devez apporter \emph{vos} raisons, \emph{vos} références, \emph{vos} exemples.
\end{attention}

\courssub{Construire des arguments personnels}

Un \cle{argument personnel} est une raison que vous construisez vous-même pour défendre votre position. Pour être valable, il doit être :
\begin{itemize}
\item \textbf{distinct} des arguments de l'auteur (sinon vous paraphrasez) ;
\item \textbf{cohérent} avec votre position annoncée (pas de contradiction interne) ;
\item \textbf{étayé} par une référence philosophique ou un exemple concret.
\end{itemize}

\textbf{Mobiliser des références philosophiques.} Votre culture du programme de Terminale technique est votre premier réservoir d'arguments. Face à une thèse sur la liberté, pensez à Sartre (« l'homme est condamné à être libre ») ou à Spinoza (la liberté comme nécessité comprise) ; face à une thèse sur le devoir, pensez à Kant (le devoir pour le devoir) ; face à une thèse sur la technique, pensez à Bergson ou aux critiques contemporaines de la technique. Citer brièvement un autre auteur montre que vous savez mettre le texte en perspective.

\textbf{Mobiliser des exemples concrets.} Un argument sans exemple reste abstrait. Puisez dans la vie courante et dans le contexte camerounais : la téléphonie mobile qui a transformé le commerce à Mokolo ou à Mvog-Mbi, le mobile money qui dispense d'aller à la banque, la mécanisation agricole dans la plaine des Mbo, l'automatisation dans les entreprises de Douala, l'usage des réseaux sociaux par les élèves. L'exemple doit \emph{illustrer} l'argument, non le remplacer.

\begin{exemplebox}[Argument personnel bien construit]
Thèse de l'auteur : « la technique libère l'homme ».\\
Mauvaise réponse : « Je suis d'accord car la technique libère l'homme du travail pénible, comme le dit l'auteur. » (simple répétition)\\
Bonne réponse : « Nous partageons partiellement ce point de vue. Certes, la technique allège les efforts : le broyeur de manioc épargne aux femmes de nos villages des heures de travail épuisant. Mais cette libération a un revers : l'automatisation détruit aussi des emplois. Notre argument est donc que la technique libère le corps tout en menaçant parfois le pain de l'ouvrier. »
\end{exemplebox}

\courssub{Envisager un contre-argument et y répondre}

\begin{definition}[Contre-argument ou objection]
Un \cle{contre-argument} (ou objection) est une raison sérieuse qu'un adversaire pourrait opposer à votre position. L'envisager soi-même, puis y \emph{répondre}, est la marque d'une discussion aboutie : cela prouve que vous avez examiné le problème des deux côtés avant de trancher.
\end{definition}

La démarche se fait en deux temps :
\begin{enumerate}
\item \textbf{Formuler l'objection} honnêtement : « On pourrait cependant objecter que\ldots », « Un adversaire de notre thèse dirait que\ldots »
\item \textbf{Y répondre} par un raisonnement : « Mais cette objection ne tient pas, car\ldots » ou « Cette objection est partiellement fondée ; elle nous oblige seulement à préciser que\ldots »
\end{enumerate}

\begin{exemplebox}[Discussion de la thèse « la technique libère l'homme »]
Un élève de Terminale technique discute ainsi : « Nous soutenons que la technique libère l'homme de la fatigue et lui ouvre du temps pour penser et créer. On pourrait toutefois objecter le \emph{chômage technologique} : dans certaines entreprises de Douala, l'introduction de machines automatisées a conduit à des licenciements, et l'ouvrier remplacé n'est pas libéré, il est mis au chômage. Mais cette objection, si elle est réelle, ne renverse pas notre thèse : elle montre seulement que la libération technique exige une organisation sociale juste — reconversion, formation professionnelle — pour profiter à tous. La technique libère donc l'homme \emph{à condition} que la société en organise équitablement les fruits. »
\end{exemplebox}

\begin{popfigure}
\begin{center}
\begin{tikzpicture}[scale=1, every node/.style={font=\small}]
% fléau de la balance
\draw[line width=1.2pt, popdark] (-4.2,0) -- (4.2,0);
\draw[line width=1.2pt, popdark] (0,0) -- (0,-0.6);
\node[below, popdark] at (0,-0.6) {\textbf{Position personnelle justifiée}};
% plateau gauche : pour
\draw[line width=1pt, popgreen] (-4.2,0) -- (-5.0,-1.6) (-4.2,0) -- (-3.4,-1.6);
\draw[line width=1.5pt, popgreen] (-5.3,-1.6) arc[start angle=180, end angle=360, radius=0.95];
\node[align=center, popgreen!60!black] at (-4.2,-2.6) {\textbf{Arguments POUR}\\ allègement du travail\\ gain de temps\\ progrès médical};
% plateau droit : contre
\draw[line width=1pt, poporange] (4.2,0) -- (3.4,-1.6) (4.2,0) -- (5.0,-1.6);
\draw[line width=1.5pt, poporange] (3.1,-1.6) arc[start angle=180, end angle=360, radius=0.95];
\node[align=center, poporange!70!black] at (4.2,-2.6) {\textbf{Arguments CONTRE}\\ chômage technologique\\ dépendance aux machines\\ inégalités d'accès};
% verdict
\node[draw=popblue, rounded corners, fill=popblueL, align=center, font=\small\bfseries] at (0,1.1) {Discussion = peser, puis trancher};
\end{tikzpicture}
\end{center}
\legende{La discussion comme balance : on pèse les arguments pour et contre, puis on fait pencher le plateau par une position personnelle \emph{justifiée}.}
\end{popfigure}

\courssub{Conclure la discussion en tranchant}

La discussion doit se terminer par une \cle{conclusion} qui tranche clairement. Après avoir pesé le pour et le contre, il faut dire nettement où l'on se situe. Une conclusion de discussion efficace :
\begin{itemize}
\item rappelle la position adoptée (« En définitive, nous estimons que\ldots ») ;
\item résume en une phrase l'argument décisif ;
\item éventuellement, précise la portée ou la limite de la position (« \ldots à condition que\ldots », « \ldots dans la mesure où\ldots »).
\end{itemize}

\begin{methode}[Discuter une thèse en 5 étapes]
\begin{enumerate}
\item \textbf{Annoncer sa position} dès la première phrase : accord, désaccord ou position nuancée.
\item \textbf{Argumenter} : formuler un ou deux arguments \emph{personnels}, distincts de ceux de l'auteur.
\item \textbf{Illustrer} : étayer chaque argument par une référence philosophique (auteur, notion du programme) et un exemple concret (vie courante, contexte camerounais).
\item \textbf{Répondre à une objection} : formuler le meilleur contre-argument possible, puis le réfuter ou l'intégrer comme nuance.
\item \textbf{Conclure} : trancher clairement en rappelant sa position et son argument décisif.
\end{enumerate}
\end{methode}

\begin{popfigure}
\begin{center}
\begin{tikzpicture}[
  grow=right,
  level distance=3.1cm,
  level 1/.style={sibling distance=2.6cm},
  level 2/.style={sibling distance=1.5cm, level distance=2.6cm},
  every node/.style={font=\scriptsize, align=center},
  edge from parent/.style={draw=popmuted, -{Stealth[length=2mm]}}
]
\node[draw=popdark, fill=popdark, text=white, rounded corners, font=\small\bfseries] {Discussion}
  child { node[draw=popblue, fill=popblueL, rounded corners, align=center]{6. Conclusion\\ qui tranche} }
  child { node[draw=poppurple, fill=poppurpleL, rounded corners, align=center]{5. Réponse à\\ l'objection} }
  child { node[draw=popgreen, fill=popgreenL, rounded corners, align=center]{4. Exemples\\ concrets}
    child { node[draw=popgreen!50, fill=popgreenL!50, rounded corners, align=center]{contexte\\ camerounais} }
  }
  child { node[draw=poporange, fill=poporangeL, rounded corners, align=center]{3. Références\\ philosophiques}
    child { node[draw=poporange!50, fill=poporangeL!50, rounded corners, align=center]{auteurs, notions\\ du programme} }
  }
  child { node[draw=poppink, fill=poppinkL, rounded corners, align=center]{2. Arguments\\ personnels} }
  child { node[draw=popgold, fill=popgoldL, rounded corners, align=center]{1. Position\\ annoncée} };
\end{tikzpicture}
\end{center}
\legende{Arbre de construction d'une discussion : de la position annoncée jusqu'à la conclusion, en passant par les arguments, les références, les exemples et la réponse à l'objection.}
\end{popfigure}

\begin{exoresolu}[Discuter une thèse : sujet d'entraînement]
\textbf{Énoncé.} Un texte soutient la thèse suivante : « Le devoir est une contrainte qui empêche l'homme d'être libre. » Partagez-vous ce point de vue ? Justifiez votre réponse.
\tcblower
\textbf{Solution détaillée.}\\
\emph{Étape 1 — Position.} Nous ne partageons pas entièrement ce point de vue : le devoir n'est pas seulement une contrainte, il peut être une voie vers la véritable liberté.\\
\emph{Étape 2 — Argument personnel.} Celui qui accomplit son devoir par choix réfléchi, et non par peur de la sanction, agit librement : il se donne à lui-même sa propre loi. La contrainte extérieure devient alors discipline intérieure.\\
\emph{Étape 3 — Référence et exemple.} Kant l'a montré : agir par devoir, c'est obéir à la loi morale que notre raison se prescrit à elle-même, donc c'est être autonome. Concrètement, l'élève de Terminale technique qui respecte ses heures de révision avant le Baccalauréat technique ne subit pas une prison : il se libère de l'échec et de la dépendance future. De même, le conducteur de taxi à Yaoundé qui respecte le code de la route accomplit un devoir qui protège sa liberté et celle des autres.\\
\emph{Étape 4 — Objection et réponse.} On objectera que certains devoirs sont imposés de l'extérieur et vécus comme de pures servitudes : corvées, obligations familiales pesantes. Cette objection est fondée pour les devoirs subis sans consentement. Mais elle confirme notre thèse plutôt qu'elle ne la renverse : c'est précisément le devoir \emph{accepté par la raison} qui libère, tandis que le devoir aveuglément subi aliène.\\
\emph{Étape 5 — Conclusion.} En définitive, le devoir n'empêche la liberté que lorsqu'il est subi ; compris et assumé, il en est la condition. Nous rejetons donc la thèse de l'auteur, en reconnaissant qu'elle décrit le cas réel mais partiel du devoir-contrainte.
\end{exoresolu}

\begin{savaistu}[La nuance tranchée, stratégie gagnante]
Les correcteurs du Baccalauréat technique le répètent dans leurs rapports de jury : une discussion \emph{nuancée mais tranchée} — « dans une large mesure, car\ldots cependant\ldots donc\ldots » — est souvent mieux notée qu'un accord total sans recul critique. Pourquoi ? Parce qu'elle montre que le candidat a examiné le problème sous plusieurs angles, qu'il a du discernement, et qu'il sait néanmoins conclure. Attention toutefois : nuancée ne veut pas dire indécise. « Peut-être oui, peut-être non » n'est pas une position philosophique.
\end{savaistu}

\begin{aretenir}[L'essentiel de la section]
\begin{itemize}
\item La question de discussion demande votre \emph{position personnelle} : apprécier, critiquer, dire si l'on partage le point de vue de l'auteur.
\item Annoncez votre position dès la première phrase : accord, désaccord ou nuance.
\item Construisez des arguments \emph{personnels}, distincts de ceux du texte, étayés par des références philosophiques et des exemples concrets (contexte camerounais).
\item Envisagez au moins un contre-argument et répondez-y : c'est la marque d'une discussion aboutie.
\item Concluez en tranchant clairement.
\end{itemize}
\end{aretenir}

\begin{exercice}[À vous de discuter]
Un texte affirme : « L'homme n'est libre que lorsqu'il obéit à ses désirs. » Discutez cette thèse en suivant les cinq étapes de la méthode : position annoncée, deux arguments personnels, une référence philosophique du programme, un exemple tiré de la vie courante au Cameroun, une objection suivie de votre réponse, et une conclusion tranchée. (Rédaction attendue : 15 à 20 lignes.)
\end{exercice}

\coursec{Application des notions philosophiques aux situations concrètes}

\courssub{Transition : de l'analyse du texte à l'ancrage dans le réel}

Après avoir maîtrisé les techniques de réponse aux questions de compréhension, d'analyse et de discussion (sections 3 et 4), l'élève doit franchir une étape décisive : \textbf{sortir du texte pour confronter les notions philosophiques à la réalité vécue}. L'exercice sur texte ne se limite pas à une exégèse interne ; il exige, notamment dans la question de discussion, que le candidat \emph{problématise} la thèse de l'auteur en l'éprouvant à l'aune d'exemples concrets, historiques, sociaux ou professionnels. C'est ici que se joue la différence entre une copie scolaire et une copie de bachelier : la capacité à \cle{philosophier en situation}, c'est-à-dire à faire opérer les concepts (liberté, devoir, État, technique, art, etc.) sur le terrain concret de l'expérience camerounaise et de la future vie professionnelle du technicien.

\courssub{Pourquoi illustrer : l'exemple, preuve de compréhension réelle}

\cle{Illustrer}, ce n'est pas « décorer » sa copie. C'est \textbf{démontrer que l'on a saisi la portée opératoire de la notion}. Une notion philosophique (ex. : \emph{la liberté}) est une construction abstraite ; elle ne devient intelligible que lorsqu'on montre comment elle structure une situation réelle (ex. : le débat sur la loi relative à la liberté de communication sociale au Cameroun).

\begin{definition}[L'exemple philosophique]
Un exemple philosophique est un \textbf{cas particulier, singulier et daté} (un fait historique, une situation juridique, un dilemme professionnel, une œuvre d'art, une innovation technique) mobilisé pour \textbf{manifester la pertinence ou les limites d'une thèse générale}. Il ne prouve pas la thèse (seul l'argument logique le fait), mais il la \textbf{rend sensible, tangible et vérifiable}.
\end{definition}

\begin{propriete}[Fonctions de l'exemple dans la dissertation et le commentaire]
\begin{itemize}
    \item \textbf{Épistémologique :} Il ancre l'abstraction dans l'expérience (phénoménologie).
    \item \textbf{Argumentatif :} Il sert de \emph{test} pour la thèse (confirmation, objection, nuance).
    \item \textbf{Pédagogique :} Il signale au correcteur la maîtrise du programme et l'ouverture d'esprit du candidat.
\end{itemize}
\end{propriete}

\courssub{Choisir des exemples pertinents, précis et en lien direct}

\cle{La pertinence} est le critère suprême. Un exemple « qui fait joli » mais qui ne cible pas la notion précise discutée est un hors-sujet déguisé.

\begin{methode}[Sélectionner et articuler un exemple pertinent]
\begin{enumerate}
    \item \textbf{Identifier la notion-clé} au cœur de la question de discussion (ex. : « L'État a-t-il pour seul rôle d'assurer la sécurité ? » $\rightarrow$ notions : \emph{État}, \emph{Sécurité}, \emph{Justice}, \emph{Légitimité}).
    \item \textbf{Interroger son stock d'exemples} (cours, actualité, stage, vie locale) : « Quel fait camerounais récent ou structurel met en jeu \emph{cette} notion sous \emph{cet} angle précis ? »
    \item \textbf{Qualifier l'exemple} : Date, lieu, acteurs, enjeu précis. Éviter le vague (« Dans notre société... », « Autrefois... »).
    \item \textbf{Articuler explicitement} : « Cet exemple illustre la thèse de Hobbes selon laquelle... car... » ou « Cet exemple contredit l'idée que... car... ».
\end{enumerate}
\end{methode}

\begin{attention}[Piège : l'anecdote personnelle non articulée]
Raconter longuement « Mon oncle a eu un problème avec la police... » ou « Lors de mon stage à l'ENEO... » sans \textbf{lien explicite et conceptuel} avec la notion (ex. : \emph{légalité}, \emph{arbitraire}, \emph{responsabilité technique}) est une \textbf{perte de temps et de points}. Le correcteur n'évalue pas votre vie, mais votre pensée philosophique. \textbf{Une phrase d'anecdote, trois phrases d'analyse conceptuelle.}
\end{attention}

\courssub{Tableau de mobilisation : Notions du programme $\times$ Situations camerounaises}

Le tableau ci-dessous synthétise les entrées principales du programme de Terminale Technique et propose des ancrages locaux précis, mobilisables en examen.

\begin{popfigure}
\legende{Tableau de correspondance : Notions philosophiques (programme Terminale Tech) et exemples de situations camerounaises exploitables.}
\begin{center}
\begin{tabularx}{\linewidth}{|>{\bfseries}l|X|X|}\hline
\rowcolor{popblueL}
\textbf{Notion / Chapitre} & \textbf{Exemples de situations camerounaises concrètes} & \textbf{Pistes d'articulation conceptuelle} \\ \hline
\textbf{Liberté} (Libre arbitre, conditionnel, politique) & \begin{itemize}\item Débats sur la \emph{Loi n°90/052} (liberté de communication) et la cybercriminalité (loi 2010/012).\item Couvre-feu / État d'urgence (régions NOSO, Covid-19) : sécurité vs liberté de circulation.\item Choix d'orientation technique (industriel vs tertiaire) : déterminisme social vs autonomie.\end{itemize} & \emph{Spinoza} : liberté = nécessité comprise. \emph{Kant} : autonomie vs hétéronomie. \emph{Rousseau} : « Obéir à la loi qu'on s'est prescrite, c'est être libre. » \\ \hline
\textbf{Devoir / Morale} (Kant, utilitarisme, devoirs civiques) & \begin{itemize}\item Paiement de l'impôt (patente, IRPP, TVA) : contrainte légale vs impératif catégorique.\item Respect du code du travail / normes de sécurité (chantier BTP, usine) : devoir professionnel vs rentabilité.\item Don de sang bénévole / entraide communautaire (njangi, tontines) : devoir imparfait vs inclination.\end{itemize} & \emph{Kant} : Devoir par devoir (bonne volonté). \emph{Durkheim} : Devoir comme fait social contraignant. Devoirs parfaits/imparfaits. \\ \hline
\textbf{État / Politique} (Légitimité, Souveraineté, Justice, Droit) & \begin{itemize}\item Décentralisation (lois 2004/017, 2004/018, 2019/024) : transfert de compétences aux CTD (Communes, Régions).\item Justice militaire vs justice civile (procès civils devant tribunaux militaires).\item Carte d'identité biométrique / Recensement : identification administrative vs reconnaissance citoyenne.\end{itemize} & \emph{Hobbes/Locke/Rousseau} : Contrat social. \emph{Weber} : Monopole de la violence physique légitime. \emph{Rawls} : Justice comme équité. \\ \hline
\textbf{Technique / Progrès} (Outil, machine, automatisation, éthique) & \begin{itemize}\item Projet \emph{Barrage de Nachtigal} / \emph{Mekin} : énergie, déplacement populations, impact éco.\item Agriculture : semences améliorées (IRAD) vs semences paysannes / OGM.\item Numérique : \emph{Mobile Money} (Orange Money, MTN MoMo) : inclusion financière vs traçabilité/surveillance.\item Formation technique : acquisition de \emph{compétences} vs \emph{aliénation} à la machine.\end{itemize} & \emph{Aristote} : \emph{Technè} / \emph{Phronèsis}. \emph{Heidegger} : Engeancement (\emph{Gestell}). \emph{Jonas} : Principe responsabilité. \emph{Marx} : Travail / Aliénation. \\ \hline
\textbf{Art / Esthétique} (Beau, Goût, Génie, Industrie culturelle) & \begin{itemize}\item \emph{Makossa}, \emph{Bikutsi}, \emph{Afrobeats} : tradition/oraliture vs production industrielle/export.\item Architecture : cases traditionnelles (Musgum, Bamileke) vs bâtiments administratifs modernistes.\item Artisanat d'art (poterie Baham, pipes Bamoun, tissus Ndop) : utilitaire vs esthétique, brevet/IGP.\end{itemize} & \emph{Kant} : Jugement de goût (désintéressement, universalité subjective). \emph{Benjamin} : Œuvre d'art à l'ère de sa reproductibilité technique. \emph{Heidegger} : L'origine de l'œuvre d'art. \\ \hline
\textbf{Science / Vérité} (Méthode, Démonstration, Modèle) & \begin{itemize}\item Recherche agronomique (IRAD, CIRAD) : protocole expérimental, variétés de manioc/maïs.\item Métrologie / Normalisation (ANOR) : étalonnage, certificats de conformité (normes ISO).\item Gestion de la pandémie Covid-19 : modèles épidémiologiques, décisions politiques sous incertitude.\end{itemize} & \emph{Descartes} : Règle de l'évidence / méthode. \emph{Popper} : Falsifiabilité. \emph{Bachelard} : Obstacle épistémologique / rupture. \\ \hline
\end{tabularx}
\end{center}
\end{popfigure}

\courssub{Lien entre les notions de l'année et la vie professionnelle du technicien}

L'enseignement technique vise une \cle{insertion professionnelle}. La philosophie n'en est pas décorative ; elle outille le futur technicien supérieur pour assumer ses \textbf{responsabilités éthiques, techniques et citoyennes}.

\begin{aretenir}[Trois registres de professionnalisation philosophique]
\begin{enumerate}
    \item \textbf{L'éthique professionnelle (Déontologie) :} Le technicien n'est pas un simple exécutant. Il engage sa \emph{responsabilité} (juridique et morale) : sécurité des usagers (bâtiment, électricité, transport), secret professionnel (informatique, santé), respect de l'environnement (rejets industriels, BTP). \emph{Ex. :} Un technicien en génie civil qui valide un béton non conforme par pression hiérarchique viole son \emph{devoir} (Kant) et sa \emph{responsabilité} (Jonas).
    \item \textbf{Rapport critique à la technique :} Maîtriser la machine, ne pas en être l'esclave. Comprendre l'\emph{aliénation} (Marx) ou l'\emph{engeancement} (Heidegger) pour concevoir des solutions \emph{appropriées} (technique intermédiaire, maintenance locale, low-tech) plutôt que de subir l'importation clé en main.
    \item \textbf{Citoyenneté technique :} Participer aux choix de société (énergie, aménagement, numérique). L'expertise technique doit éclairer le débat public (ex. : avis d'ingénieurs sur le barrage de Lom Pangar, sur la 5G, sur la gestion des déchets plastiques).
\end{enumerate}
\end{aretenir}

\begin{exemplebox}[Illustration : Le devoir professionnel et l'impératif catégorique]
\textbf{Situation :} Un technicien en maintenance industrielle découvre une fissure critique sur une cuve sous pression dans une usine agroalimentaire. Son chef lui demande de « colmater à la soudure rapide » pour ne pas arrêter la chaîne de production (perte financière), promettant une réparation « plus tard ».
\textbf{Analyse philosophique :}
\begin{itemize}
    \item \textbf{Maxime de l'action :} « Je peux mentir sur l'état de sécurité / contourner la procédure pour sauver la production. »
    \item \textbf{Universalisation (Kant) :} Si tout technicien faisait cela, la confiance dans la certification technique s'effondrerait, les accidents se multiplieraient $\rightarrow$ \textbf{Contradiction dans la conception}. L'action est \textbf{interdite}.
    \item \textbf{Humanité comme fin :} Exposer les ouvriers/riverains à un risque mortel pour un profit, c'est les traiter comme \emph{un moyen} seulement.
    \item \textbf{Responsabilité (Jonas) :} L'effet de l'action (explosion potentielle) est irréversible et touche des innocents $\rightarrow$ Principe de précaution éthique.
\end{itemize}
\textbf{Conclusion :} Le technicien a un \textbf{devoir de désobéissance éthique} (alerter, consigner par écrit, refuser l'intervention non conforme), fondé sur la primauté de la \emph{loi morale} (sécurité, vie) sur l'ordre hiérarchique illégitime.
\end{exemplebox}

\courssub{Distinguer l'exemple (cas particulier) de l'argument (raison générale)}

C'est l'erreur méthodologique n°1 en discussion : \textbf{confondre illustration et démonstration}.

\begin{propriete}[Distinction fondamentale]
\begin{center}
\begin{tabularx}{\linewidth}{|>{\bfseries}l|X|X|}\hline
\rowcolor{popblueL}
& \textbf{L'Argument (Raison générale)} & \textbf{L'Exemple (Cas particulier)} \\ \hline
\textbf{Nature} & Proposition logique, universelle, abstraite & Fait singulier, daté, concret \\ \hline
\textbf{Rôle} & \textbf{Démontre} / Justifie la thèse & \textbf{Illustre} / Éprouve / Incarnne \\ \hline
\textbf{Forme} & « Puisque... », « Car... », « En effet... » & « Par exemple... », « C'est le cas de... », « On observe que... » \\ \hline
\textbf{Test} & Valide par la logique (déduction, induction) & Valide par la \emph{reconnaissance} (oui, cela existe) \\ \hline
\end{tabularx}
\end{center}
\end{propriete}

\begin{methode}[Structure standard d'un paragraphe de discussion « Argument + Exemple »]
\begin{enumerate}
    \item \textbf{Thèse / Argument} : Énoncer la raison générale (ex. : « L'État ne se réduit pas à la sécurité car la justice sociale est condition de sa légitimité [Rawls]. »).
    \item \textbf{Transition explicite} : « Cela se vérifie concrètement... », « Prenons l'exemple de... ».
    \item \textbf{Exemple précis} : Citer le fait camerounais (ex. : « La politique des \emph{écoles de la seconde chance} ou les programmes \emph{PUDC} (Plan d'Urgence de Développement Communautaire) visent à réduire les inégalités territoriales. »).
    \item \textbf{Retour conceptuel (Indispensable !)} : « Cet exemple montre que l'État assume une fonction \emph{distributive} (justice), non seulement \emph{répressive} (sécurité). Il confirme donc que... ».
\end{enumerate}
\end{methode}

\begin{attention}[Erreur fatale : L'exemple « flottant »]
Ne jamais laisser un exemple « flotter » sans \textbf{retour conceptuel}. \textbf{Mal} : « Au Cameroun, il y a des routes qui se dégradent vite. » (Constat brut, sans notion). \textbf{Bien} : « La dégradation précoce de la route Yaoundé-Douala malgré les péages illustre le problème de la \emph{délégation de service public} : le contractant (privé) cherche le profit, l'État (régulateur) faillit à son devoir de \emph{contrôle technique}, compromettant le \emph{bien commun}. »
\end{attention}

\courssub{Vérifier la pertinence : l'exemple doit éclairer, non remplacer}

\begin{methode}[Check-list de validation d'un exemple (à appliquer en brouillon)]
\begin{enumerate}
    \item \textbf{Ciblage :} L'exemple traite-t-il \emph{exactement} la notion en jeu dans la question (ex. : « Technique » et non « Science » ; « Devoir juridique » et non « Devoir moral ») ?
    \item \textbf{Précision :} Est-il daté, localisé, nommé ? (Éviter : « Dans les entreprises... » $\rightarrow$ Préférer : « Lors de la construction du pont sur le Wouri (2013-2017)... »).
    \item \textbf{Articulation :} La phrase de lien (« Cet exemple illustre... car... ») est-elle rédigée ?
    \item \textbf{Proportion :} L'exemple occupe-t-il \textbf{moins de 30\%} du paragraphe ? (L'argumentation doit dominer).
    \item \textbf{Non-équivoque :} L'exemple ne contredit-il pas involontairement ma thèse ? (Ex. : Citer la répression d'une manifestation pour illustrer « L'État garantit la liberté »).
\end{enumerate}
\end{methode}

\begin{savaistu}[Le regard du correcteur : l'ancrage local, marqueur d'excellence]
Les correcteurs du Probatoire / Baccalauréat technique au Cameroun sont instruits de valoriser les copies qui \textbf{sortent du seul manuel scolaire}. Un élève qui mobilise \emph{le Code du travail camerounais} pour illustrer le contrat de travail (liberté/obligation), \emph{la loi sur la décentralisation} pour l'État, \emph{le Mobile Money} pour la technique/liberté, ou \emph{le Ngondo / Nyem-Nyem} pour l'art/tradition, signale une \textbf{philosophie vivante, ancrée, critique}. C'est souvent la différence entre la mention \emph{Assez Bien} et \emph{Bien}, ou entre l'admis et le recalé de justesse. \textbf{Philosophier, c'est penser \emph{ici et maintenant} avec les outils des grands auteurs.}
\end{savaistu}

\courssub{Exercice résolu : De la notion à l'exemple articulé}

\begin{exoresolu}[Sujet type : « La technique libère-t-elle l'homme ? »]
\textbf{Énoncé :} À partir de la notion de \emph{technique} (chapitre : Technique et Progrès), proposer deux exemples camerounais distincts : un qui soutient la thèse de la libération, un qui soutient la thèse de l'aliénation. Les articuler chacun à un auteur du programme.

\textbf{Solution détaillée :}

\textbf{Exemple 1 : Libération — Le Mobile Money (Orange Money / MTN MoMo) et l'inclusion financière.}
\begin{itemize}
    \item \textbf{Fait :} Depuis 2011, le paiement mobile permet à des millions de Camerounais non bancarisés (ruraux, informels, femmes au foyer) d'épargner, transférer, payer des factures (Eneo, Camwater), recevoir des salaires sans compte bancaire classique.
    \item \textbf{Articulation (Thèse libératrice) :} Cela illustre \textbf{Marx} (\emph{Le Capital}, Livre I) : la technique comme « développement des forces productives » qui réduit la \emph{peine} (peine de se déplacer, risque de vol d'espèces, exclusion du crédit) et augmente la \emph{puissance} d'agir économique des plus faibles.
    \item \textbf{Nuance (Maîtrise) :} C'est une libération \emph{conditionnelle} : elle suppose l'accès au téléphone, au réseau, à la littératie numérique (fracture numérique).
\end{itemize}

\textbf{Exemple 2 : Aliénation — L'orpaillage artisanal et semi-mécanisé à l'Est (Bétaré-Oya, Kambele) et au Nord (Bouré).}
\begin{itemize}
    \item \textbf{Fait :} L'introduction de pompes à eau, broyeurs, détecteurs de métaux, et surtout l'usage du \emph{mercure} et du \emph{cyanure} pour l'extraction, transforment l'activité. Le creuseur devient opérateur de machine, exposé à des toxiques, dépendant du cours de l'or et des acheteurs (souvent étrangers), détruisant son milieu (rivières, forêts).
    \item \textbf{Articulation (Thèse aliénante) :} Cela illustre \textbf{Heidegger} (\emph{La question de la technique}) : l'\emph{engeancement} (\emph{Gestell}). La nature (rivière, sol) devient un « stock » (réserve d'or) à exploiter sans mesure. L'homme lui-même devient une « ressource humaine » jetable. Cela rejoint \textbf{Jonas} (\emph{Le principe responsabilité}) : la puissance technique (chimie de l'extraction) dépasse la responsabilité éthique (santé, écosystème).
    \item \textbf{Lien pro :} Le technicien en mines/géologie/environnement a ici un \emph{devoir} de régulation (normes, réhabilitation, formation aux méthodes sans mercure).
\end{itemize}

\textbf{Bilan :} La technique \emph{ambivalente} (liberté/aliénation) exige du technicien une \emph{phronèsis} (sagesse pratique, Aristote) pour orienter l'innovation vers l'\emph{autonomie} réelle, non la seule performance technique.
\end{exoresolu}

\begin{exercice}[Entraînement : Mobilisation croisée]
\textbf{Consigne :} Choisissez \textbf{trois} notions parmi la liste : \emph{Liberté, Devoir, État, Technique, Art, Vérité}. Pour chacune, proposez \textbf{un} exemple camerounais précis (fait, loi, événement, pratique pro) et rédigez \textbf{la phrase d'articulation conceptuelle} (2 lignes max) qui le relie à un auteur ou une thèse du programme.
\textit{Ex. de réponse attendue : Notion : Devoir. Exemple : L'obligation vaccinale (Covid-19, fièvre jaune) pour les agents de santé/public. Articulation : « Cet exemple illustre le \emph{devoir civique} (Rousseau/Contrat social) où la liberté individuelle s'efface devant la \emph{volonté générale} (santé publique), mais pose la limite du \emph{consentement éclairé} (Kant/Bioéthique). »}
\end{exercice}

\begin{corrige}[Exercice n°1 - Pistes de correction]
\textit{Les réponses varient. Vérifier : 1. Précision de l'exemple (nom, date, contexte). 2. Justesse de la notion visée. 3. Pertinence de l'auteur/thèse cité. 4. Clarté du lien logique (illustre / contredit / nuance).}
\begin{itemize}
    \item \textbf{Liberté} : Ex. Débat loi sur les \textbf{ONG} (2024) / Liberté d'association. Articulation : \emph{Constant} (Liberté des anciens vs modernes) / \emph{Berlin} (Liberté négative/positive).
    \item \textbf{État} : Ex. Création des Régions (2019-2020), élection conseillers régionaux. Articulation : \emph{Tocqueville} (Décentralisation/liberté locale) / \emph{Weber} (Monopole violence légitime partagé).
    \item \textbf{Art} : Ex. Classement « Ngondo » ou « Ngiri » au patrimoine UNESCO (PCI). Articulation : \emph{Hegel} (Art comme manifestation de l'Esprit/Histoire d'un peuple) / \emph{Benjamin} (Aura vs reproductibilité / muséification).
    \item \textbf{Vérité} : Ex. Polémique sur les résultats d'examens (Probatoire / Baccalauréat technique) / notes tronquées. Articulation : \emph{Descartes} (Preuve/évidence) / \emph{Nietzsche} (Volonté de puissance / « Il n'y a pas de faits, que des interprétations » - critique de l'institution).
\end{itemize}
\end{corrige}

\coursec{Rédaction, justification de la démarche et gestion de l'épreuve}

Dans la section précédente, nous avons appris à mobiliser les notions philosophiques de l'unité pour éclairer des situations concrètes de la vie courante. Mais savoir penser ne suffit pas au Baccalauréat technique : il faut encore \emph{faire voir} au correcteur que l'on sait penser. Or, entre le candidat et le correcteur, il n'y a qu'un médiateur : la copie. Une pensée juste mais mal présentée, mal rédigée ou mal organisée sera lue comme une pensée confuse. Cette dernière section est donc consacrée à l'art de transformer un travail de réflexion en une copie claire, structurée et convaincante : rédaction de la présentation du texte, justification de la démarche, qualité de l'écriture, gestion du temps et relecture finale.

\courssub{La structure globale de la copie}

Commençons par une intuition simple. Imaginez que vous receviez une lettre administrative sans objet, sans date, sans signature : vous ne sauriez ni qui écrit, ni de quoi il s'agit, ni ce que l'on attend de vous. La copie de philosophie obéit à la même logique : avant de répondre aux questions, le candidat doit \cle{situer} le texte qu'on lui propose. La structure d'une copie réussie d'exercice sur texte comporte donc deux grands moments :

\begin{enumerate}
\item une \cle{brève présentation du texte} (auteur, œuvre, thème, thèse), rédigée en quelques lignes en tête de copie ;
\item les \cle{réponses numérotées} aux questions, dans l'ordre du sujet, chacune rédigée et justifiée.
\end{enumerate}

\begin{definition}[Présentation du texte]
La présentation du texte est un court paragraphe introductif (quatre à six lignes) qui identifie l'auteur du texte, l'œuvre dont il est extrait, le thème traité et la thèse défendue. Elle n'est ni une paraphrase ni un résumé : elle situe le texte avant de l'exploiter.
\end{definition}

\begin{methode}[Rédiger la présentation du texte en quatre lignes]
\begin{enumerate}
\item \textbf{L'auteur} : nommer l'auteur et, si possible, le situer (époque, courant). Exemple : « Le texte proposé est de Kant, philosophe allemand du XVIII\textsuperscript{e} siècle. »
\item \textbf{L'œuvre} : préciser l'ouvrage dont le passage est extrait, si l'énoncé l'indique. Exemple : « Il est extrait de la \emph{Critique de la raison pure}. »
\item \textbf{Le thème} : formuler en une phrase le problème dont traite le passage. Exemple : « L'auteur y traite du thème de la connaissance, et plus précisément des limites de l'expérience. »
\item \textbf{La thèse} : énoncer en une phrase l'idée directrice défendue. Exemple : « Il y soutient que toute connaissance commence avec l'expérience, mais ne dérive pas entièrement d'elle. »
\end{enumerate}
\end{methode}

\begin{attention}[Erreur fréquente : négliger la relecture et la numérotation]
Beaucoup de candidats, pressés par le temps, rendent leur copie sans la relire et sans numéroter leurs réponses. C'est une faute stratégique : les fautes de langue (orthographe, accords, syntaxe) donnent une mauvaise impression générale, et les réponses non numérotées obligent le correcteur à chercher ce qui correspond à quelle question — au risque qu'il ne trouve pas une réponse pourtant présente. Une réponse introuvable est une réponse non notée. Prévoyez toujours du temps pour la relecture et indiquez clairement le numéro de chaque question avant votre réponse.
\end{attention}

\courssub{Justifier sa démarche : annoncer ce que l'on va faire}

Une copie de philosophie n'est pas seulement un contenu : c'est un \cle{parcours} que le candidat guide. Le correcteur doit pouvoir suivre ce parcours sans effort. Pour cela, deux outils sont indispensables.

\textbf{1. L'annonce de démarche.} Au début de la copie ou au début d'une réponse longue (notamment la question de discussion), il est excellent d'annoncer ce que l'on va faire, à l'aide de formules comme : « Nous montrerons d'abord que..., ensuite nous verrons que..., enfin nous discuterons... ». Cette annonce remplit une double fonction : elle prouve au correcteur que le candidat a un plan, et elle oblige le candidat lui-même à s'y tenir.

\textbf{2. Les connecteurs logiques.} Les articulations du raisonnement doivent être visibles : \emph{d'abord, ensuite, enfin} pour l'ordre ; \emph{en effet, car, puisque} pour la justification ; \emph{mais, cependant, néanmoins} pour l'objection ; \emph{donc, par conséquent, c'est pourquoi} pour la conclusion. Un texte sans connecteurs est une suite d'affirmations ; un texte avec connecteurs est un raisonnement.

\begin{exemplebox}[Annonce de démarche pour une question de discussion]
Sujet : « La liberté se réduit-elle à faire ce que l'on veut ? »

\emph{Annonce rédigée} : « Nous montrerons d'abord, avec l'auteur du texte, que la liberté ne saurait se confondre avec le caprice ; nous verrons ensuite qu'elle suppose la connaissance des motifs qui nous déterminent ; nous conclurons enfin qu'être libre, ce n'est pas faire ce que l'on veut, mais vouloir ce que l'on a soi-même examiné. »
\end{exemplebox}

\courssub{Rédiger des réponses complètes et justifiées}

Chaque réponse à une question de compréhension ou d'analyse doit obéir à une règle simple : \textbf{une réponse = une idée + une justification}. Concrètement, une bonne réponse se construit en trois temps :

\begin{enumerate}
\item \textbf{Répondre directement} à la question posée, en phrases construites (jamais en style télégraphique ni en simple recopie du texte) ;
\item \textbf{Justifier par le texte} à l'aide d'une courte citation entre guillemets, suivie de sa référence (« l'auteur écrit : "..." ») ;
\item \textbf{Expliquer ou reformuler} la citation pour montrer qu'on la comprend (« autrement dit... », « ce qui signifie que... »).
\end{enumerate}

Le vocabulaire employé doit être \cle{philosophiquement précis} : on ne dit pas « l'auteur pense que la technique c'est bien », mais « l'auteur soutient que la technique est un moyen ambivalent, capable d'émanciper comme d'aliéner l'homme ». La précision du vocabulaire est un indice majeur de maîtrise pour le correcteur.

\begin{exoresolu}[Rédiger une réponse de compréhension justifiée]
Énoncé : texte de Bergson sur la conscience. Question : « Selon l'auteur, qu'est-ce qui distingue l'homme conscient de l'animal ? »

\tcblower

\textbf{Solution détaillée.}

\textbf{Question 1.} Selon Bergson, ce qui distingue l'homme conscient de l'animal, c'est la capacité de se représenter le passé et l'avenir, et non seulement de vivre dans l'instant présent. L'auteur écrit en effet : « L'animal vit rivé au présent, tandis que la conscience humaine élargit le présent en y retenant le passé ». Autrement dit, la conscience est pour Bergson une mémoire active : elle permet à l'homme de ne pas subir le moment, mais de le situer dans une durée, ce qui rend possible la prévision et le choix.

\emph{Commentaire méthodologique} : la réponse commence par répondre directement à la question (l'idée), cite le texte entre guillemets (la justification), puis reformule avec « autrement dit » (l'explication). Trois phrases suffisent.
\end{exoresolu}

\begin{attention}[Pièges de rédaction]
\begin{itemize}
\item Ne jamais recopier de longs passages du texte : une citation de plus de deux lignes sans commentaire est sanctionnée comme de la paraphrase.
\item Ne jamais répondre par « oui » ou « non » sans développement : même une question fermée exige une justification.
\item Ne pas confondre l'opinion personnelle et la thèse de l'auteur : dans les questions de compréhension, c'est la pensée \emph{de l'auteur} qui est demandée (« selon l'auteur... ») ; l'opinion du candidat n'intervient que dans la question de discussion.
\end{itemize}
\end{attention}

\courssub{La gestion du temps pendant l'épreuve}

L'épreuve de philosophie au Baccalauréat technique dure quatre heures. Quatre heures, c'est long pour qui n'a pas de plan de bataille, et terriblement court pour qui s'y prend mal. L'expérience des jurys montre que les échecs tiennent moins souvent à l'ignorance qu'à une mauvaise répartition du temps : candidats qui rédigent au propre dès la première heure et s'essoufflent, ou qui peaufinent leur brouillon pendant trois heures et recopient dans la panique.

\begin{propriete}[Répartition conseillée du temps sur 4 heures]
Pour un exercice sur texte, la répartition indicative suivante est recommandée :
\begin{itemize}
\item \textbf{Lecture et analyse du texte} : 30 minutes (lecture active, repérage de la thèse, plan du texte au brouillon) ;
\item \textbf{Brouillon} : 45 minutes (réponse à chaque question sous forme de notes, choix des citations, plan de la discussion) ;
\item \textbf{Rédaction au propre} : 2 h 15 ;
\item \textbf{Relecture} : 30 minutes (orthographe, numérotation, cohérence).
\end{itemize}
\end{propriete}

\begin{popfigure}
\begin{center}
\begin{tikzpicture}
\begin{axis}[
    ybar,
    width=0.95\linewidth,
    height=7cm,
    bar width=1.1cm,
    ymin=0, ymax=160,
    ylabel={Durée (minutes)},
    symbolic x coords={Lecture, Brouillon, Rédaction, Relecture},
    xtick=data,
    nodes near coords,
    nodes near coords align={vertical},
    every node near coord/.append style={font=\small\bfseries},
    ymajorgrids=true,
    grid style={popline!40},
    axis line style={popdark},
]
\addplot[fill=popblue, draw=popdark] coordinates {(Lecture,30) (Brouillon,45) (Rédaction,135) (Relecture,30)};
\end{axis}
\end{tikzpicture}
\end{center}
\legende{Répartition conseillée du temps pour une épreuve d'exercice sur texte de 4 heures (240 minutes) : lecture et analyse 30 min, travail au brouillon 45 min, rédaction au propre 2 h 15 (135 min), relecture 30 min. La rédaction occupe un peu plus de la moitié de l'épreuve, mais elle n'est efficace que si les trois quarts d'heure de préparation l'ont précédée.}
\end{popfigure}

\begin{savaistu}[Le brouillon, votre meilleur assurance contre le hors-sujet]
Les rapports de jury du Baccalauréat le répètent chaque année : les candidats qui écrivent au brouillon le plan du texte (thèse, étapes de l'argumentation, concepts clés) \emph{avant} de répondre aux questions divisent quasiment par deux le risque de hors-sujet. Pourquoi ? Parce que le hors-sujet naît presque toujours d'une lecture trop rapide : le candidat répond à la question qu'il imaginait, non à celle qui était posée. Dix minutes de plan au brouillon valent mieux qu'une heure de rédaction à côté de la plaque.
\end{savaistu}

\courssub{La relecture finale et la présentation matérielle}

La relecture n'est pas un luxe : c'est la dernière étape du travail intellectuel. Elle porte sur quatre points, à vérifier dans l'ordre :

\begin{enumerate}
\item \textbf{La numérotation} : chaque réponse est-elle précédée du numéro de la question correspondante ? Aucune question n'a-t-elle été oubliée ?
\item \textbf{La cohérence} : les réponses se contredisent-elles entre elles ? La thèse attribuée à l'auteur dans la présentation est-elle la même que celle utilisée dans les réponses ?
\item \textbf{La langue} : orthographe des noms d'auteurs (Kant, Descartes, Bergson...), accords, conjugaison, ponctuation.
\item \textbf{Le respect des consignes} : sujet traité intégralement, nombre de questions respecté.
\end{enumerate}

La présentation matérielle joue également un rôle réel, quoique indirect : une copie lisible, aérée (paragraphes séparés, alinéas, pas de ratures envahissantes), écrite à l'encre foncée, met le correcteur dans de bonnes dispositions. À l'inverse, une copie illisible fatigue le lecteur et lui fait manquer des éléments corrects. Il ne s'agit pas d'avoir une belle écriture, mais une écriture \emph{déchiffrable} et une page \emph{organisée}.

\begin{exemplebox}[L'effet chiffré d'une copie bien présentée]
Sur 20 points, les correcteurs constatent régulièrement qu'une copie moyenne sur le fond mais clairement structurée (présentation du texte, réponses numérotées, citations, démarche annoncée) gagne de 2 à 3 points par rapport à une copie de contenu équivalent mais désorganisée. Autrement dit, la forme ne remplace pas le fond, mais elle le fait valoir : à contenu égal, c'est la clarté de la démarche qui fait la différence entre une note de 10 et une note de 13.
\end{exemplebox}

\begin{popfigure}
\begin{center}
\begin{tikzpicture}[
    box/.style={draw=popdark, fill=popblueL, rounded corners=3pt, text width=0.82\linewidth, align=left, inner sep=6pt, font=\small},
    disc/.style={draw=popdark, fill=poporangeL, rounded corners=3pt, text width=0.82\linewidth, align=left, inner sep=6pt, font=\small},
    concl/.style={draw=popdark, fill=popgreenL, rounded corners=3pt, text width=0.82\linewidth, align=left, inner sep=6pt, font=\small},
    arr/.style={-{Stealth[length=3mm]}, thick, popdark}
]
\node[box, align=center] (p) at (0,0){\textbf{Présentation du texte} (4 à 6 lignes)\\ Auteur — œuvre — thème — thèse};
\node[box, below=0.45cm of p] (q1) {\textbf{Question 1} : réponse directe + courte citation + explication};
\node[box, below=0.45cm of q1] (q2) {\textbf{Question 2} : réponse directe + courte citation + explication};
\node[box, below=0.45cm of q2] (q3) {\textbf{Question 3 (analyse)} : idée dégagée + justification par le texte + notion du programme};
\node[disc, below=0.45cm of q3] (qd) {\textbf{Question de discussion} : annonce de la démarche (« nous montrerons d'abord... ensuite... »), thèse de l'auteur, objection, position argumentée du candidat};
\node[concl, below=0.45cm of qd] (c) {\textbf{Conclusion de la discussion} : réponse claire à la question posée + ouverture éventuelle};
\draw[arr] (p) -- (q1);
\draw[arr] (q1) -- (q2);
\draw[arr] (q2) -- (q3);
\draw[arr] (q3) -- (qd);
\draw[arr] (qd) -- (c);
\end{tikzpicture}
\end{center}
\legende{Modèle de structure d'une copie réussie d'exercice sur texte : une brève présentation du texte, puis les réponses numérotées dans l'ordre du sujet (chacune justifiée par une courte citation), la question de discussion construite en étapes annoncées, et une conclusion qui répond explicitement à la question.}
\end{popfigure}

\begin{exoresolu}[Transformer une réponse faible en réponse réussie]
Énoncé : voici la réponse d'un candidat à la question « Quelle est la thèse de l'auteur ? » : \emph{« L'auteur parle de la liberté et il dit que c'est important pour l'homme, tout le monde veut être libre, c'est vrai. »} Identifiez les défauts de cette réponse et rédigez-en une version correcte (le texte est extrait de l'\emph{Éthique} de Spinoza, qui soutient que les hommes se croient libres parce qu'ils ont conscience de leurs désirs mais ignorent les causes qui les déterminent).

\tcblower

\textbf{Solution détaillée.}

\textbf{1. Défauts de la réponse.} (a) Elle ne répond pas à la question : elle parle de la liberté en général au lieu d'énoncer la thèse précise de l'auteur. (b) Elle substitue l'opinion du candidat (« c'est vrai ») à l'analyse du texte. (c) Elle ne contient aucune citation ni aucune justification. (d) Son vocabulaire est vague (« c'est important »).

\textbf{2. Version correcte.} « La thèse de Spinoza, dans ce passage extrait de l'\emph{Éthique}, est que la liberté dont les hommes se croient dotés est une illusion fondée sur l'ignorance. L'auteur écrit en effet que les hommes "ont conscience de leurs désirs, mais ignorent les causes par lesquelles ils sont déterminés à désirer". Autrement dit, nous nous sentons libres non parce que nous le sommes réellement, mais parce que nous percevons nos actes sans en connaître les causes : la véritable liberté suppose donc la connaissance de ce qui nous détermine. »

\emph{Commentaire} : la réponse corrigée nomme l'auteur et l'œuvre, énonce la thèse en une phrase, la justifie par une courte citation, puis l'explique avec le vocabulaire philosophique adéquat (illusion, détermination, connaissance des causes).
\end{exoresolu}

\begin{aretenir}[L'essentiel de la section]
\begin{itemize}
\item La copie s'ouvre sur une \cle{présentation du texte} en quatre lignes : auteur, œuvre, thème, thèse.
\item Chaque réponse est \cle{numérotée}, rédigée en phrases construites, et suit le schéma : idée + courte citation + explication.
\item La démarche est \cle{annoncée} (« nous montrerons d'abord... ensuite... ») et articulée par des connecteurs logiques.
\item Le temps se répartit ainsi : 30 min de lecture, 45 min de brouillon, 2 h 15 de rédaction, 30 min de relecture.
\item La relecture porte sur la numérotation, la cohérence, la langue et le respect des consignes ; une copie claire peut gagner 2 à 3 points à contenu égal.
\end{itemize}
\end{aretenir}

\begin{exercice}[Entraînement à la gestion de l'épreuve]
Un sujet d'exercice sur texte comporte quatre questions de compréhension et une question de discussion. (1) Établissez un emploi du temps détaillé des quatre heures de l'épreuve, en indiquant l'heure à laquelle chaque étape doit être achevée si l'épreuve commence à 8 h. (2) Rédigez la présentation du texte suivant en quatre lignes : extrait du \emph{Discours de la méthode} de Descartes, où l'auteur soutient que le bon sens est la chose du monde la mieux partagée, mais que l'essentiel est de bien conduire sa raison.
\end{exercice}

\begin{corrige}[Exercice : entraînement à la gestion de l'épreuve]
\textbf{1. Emploi du temps.} 8 h – 8 h 30 : lecture active du texte et repérage de la thèse ; 8 h 30 – 9 h 15 : brouillon (plan du texte, réponses en notes, choix des citations) ; 9 h 15 – 11 h 30 : rédaction au propre (environ 20 minutes par question de compréhension, 45 minutes pour la discussion) ; 11 h 30 – 12 h : relecture complète.

\textbf{2. Présentation du texte.} « Le texte proposé est de Descartes, philosophe français du XVII\textsuperscript{e} siècle. Il est extrait du \emph{Discours de la méthode}. L'auteur y traite du thème de la raison et de son usage. Il y soutient que le bon sens est également réparti entre tous les hommes, et que ce qui fait la différence entre eux n'est pas la quantité de raison, mais la manière de la conduire. »
\end{corrige}

\coursec{Travaux pratiques de consolidation guidés}

Après avoir rappelé les règles de rédaction, de justification de la démarche et de gestion du temps d'épreuve, il reste à passer de la théorie à la pratique. Une méthode ne se possède vraiment que lorsqu'on l'a exercée plusieurs fois, corrigée, réexercée. C'est le sens de cette dernière étape : transformer les exigences de l'exercice sur texte en réflexes, grâce à quatre travaux pratiques (TP) progressifs, une grille d'autoévaluation et une correction collective de copies anonymisées. Le principe directeur est simple : \emph{entraîner, corriger, s'autoévaluer, recommencer}.

\begin{popfigure}
\begin{center}
\begin{tikzpicture}[node distance=1.1cm,
  etape/.style={rectangle, rounded corners=4pt, draw=popblue, fill=popblueL, text width=3.4cm, align=center, minimum height=1.1cm, font=\small\bfseries},
  fleche/.style={-{Stealth[length=3mm]}, thick, poporange}]
\node[etape] (e1) {1. Entraînement (TP et sujets)};
\node[etape, right=of e1] (e2) {2. Correction (croisée ou collective)};
\node[etape, below=of e2] (e3) {3. Autoévaluation (grille critériée)};
\node[etape, left=of e3] (e4) {4. Nouvel entraînement (progression)};
\draw[fleche] (e1) -- (e2);
\draw[fleche] (e2) -- (e3);
\draw[fleche] (e3) -- (e4);
\draw[fleche] (e4) -- (e1);
\node[font=\small\itshape, text=popdark] at ($(e1)!0.5!(e3)$) {Cycle de consolidation};
\end{tikzpicture}
\end{center}
\legende{Le cycle de consolidation : chaque passage dans la boucle renforce les réflexes méthodologiques. La progression vient de la répétition corrigée, non du seul entraînement.}
\end{popfigure}

\courssub{TP 1 : lecture guidée d'un texte et repérage collectif de la thèse et des arguments}

\textbf{Objectif.} Vérifier que chaque élève sait dégager la \cle{thèse} d'un texte et la structure de son \cle{argumentation}, première compétence de l'exercice sur texte.

\textbf{Déroulement.} Le professeur distribue un extrait court (10 à 15 lignes) d'un auteur au programme, par exemple Kant, Platon ou Descartes. La lecture se fait en trois temps :
\begin{enumerate}
\item \textbf{Lecture silencieuse individuelle} (5 minutes) : chaque élève surligne d'une couleur la phrase qui exprime la thèse, d'une autre couleur les arguments.
\item \textbf{Mise en commun au tableau} : le professeur recueille les propositions, confronte les désaccords et fait justifier chaque choix par des indices textuels précis (connecteurs logiques comme « donc », « car », « en effet » ; formules d'annonce ; exemples).
\item \textbf{Validation collective} : la classe formule en une phrase la thèse de l'auteur, puis numérote les arguments dans l'ordre du texte.
\end{enumerate}

\begin{exemplebox}[Sujet type : extrait de Kant, \emph{Qu'est-ce que les Lumières ?} (1784)]
Texte proposé : « Les Lumières, c'est la sortie de l'homme hors de l'état de minorité dont il est lui-même responsable. L'état de minorité est l'incapacité de se servir de son entendement sans la conduite d'un autre. [...] \emph{Sapere aude !} Aie le courage de te servir de ton propre entendement ! Voilà la devise des Lumières. »
\par Trois questions type Baccalauréat technique :
\begin{enumerate}
\item Expliquez la thèse de l'auteur.
\item Dégagez l'argumentation qui la soutient.
\item Discutez cette thèse : l'homme est-il toujours responsable de sa minorité ?
\end{enumerate}
Repérage attendu : la \textbf{thèse} est dans la première phrase (sortir de la minorité par l'usage de sa propre raison) ; les \textbf{arguments} sont la définition de la minorité (incapacité de penser par soi-même) et l'appel au courage de penser (\emph{Sapere aude !}).
\end{exemplebox}

\begin{attention}[Piège du repérage]
Ne confonds pas la thèse avec un exemple ou une citation. La thèse est l'idée générale que l'auteur défend ; les exemples et citations ne font que l'illustrer ou l'appuyer. Dans l'extrait de Kant, la devise latine n'est pas la thèse : elle la résume sous forme d'appel.
\end{attention}

\courssub{TP 2 : rédaction individuelle et correction croisée entre pairs}

\textbf{Objectif.} Produire une réponse écrite à une question de compréhension, puis apprendre à évaluer le travail d'autrui avec des critères objectifs.

\textbf{Déroulement.}
\begin{enumerate}
\item \textbf{Rédaction individuelle} (15 minutes) : chaque élève répond par écrit à la question 1 du sujet type (« Expliquez la thèse de l'auteur »), en 10 à 15 lignes, sans paraphrase.
\item \textbf{Échange des copies} : chacun corrige la copie d'un camarade à l'aide de la grille critériée ci-dessous.
\item \textbf{Restitution} : le correcteur explique oralement sa notation à l'auteur de la copie, qui peut défendre ses choix.
\end{enumerate}

\begin{attention}[Erreur fréquente en TP]
Corriger son camarade uniquement sur l'orthographe ou la présentation, et non sur la \textbf{pertinence philosophique} des réponses. Une copie bien écrite mais hors sujet ou paraphrasée reste une mauvaise copie de philosophie. La correction croisée porte d'abord sur la fidélité au texte, la justesse de l'explication et la qualité de la justification.
\end{attention}

\begin{popfigure}
\begin{center}
\begin{tikzpicture}
\node[anchor=north west, align=center] at (0,0){%
\begin{tabular}{|p{3.2cm}|c|p{6.2cm}|}
\hline
\rowcolor{popblueL} \textbf{Critère} & \textbf{Points} & \textbf{Indicateurs de réussite} \\
\hline
Fidélité au texte & /4 & La thèse rapportée est bien celle de l'auteur ; aucune idée étrangère au texte ; aucune paraphrase. \\
\hline
Pertinence de la réponse & /4 & La réponse répond exactement à la question posée ; le sujet n'est pas déplacé. \\
\hline
Qualité de l'argumentation & /6 & Chaque affirmation est justifiée ; les arguments du texte sont correctement dégagés et ordonnés ; les exemples sont bien choisis. \\
\hline
Présentation et expression & /6 & Plan visible, phrases claires, connecteurs logiques, orthographe et vocabulaire philosophique corrects. \\
\hline
\rowcolor{poporangeL} \textbf{Total} & \textbf{/20} & \\
\hline
\end{tabular}};
\end{tikzpicture}
\end{center}
\legende{Grille critériée d'évaluation utilisée pour la correction croisée et l'autoévaluation. Les points peuvent être adaptés au barème réel de l'épreuve.}
\end{popfigure}

\courssub{TP 3 : construction collective d'une discussion (pour / contre)}

\textbf{Objectif.} Apprendre à répondre à la question de \cle{discussion}, la plus difficile de l'exercice, en construisant ensemble un débat argumenté sur la thèse du texte.

\textbf{Déroulement.} Le tableau est divisé en deux colonnes : « Arguments pour la thèse de Kant » et « Arguments contre (ou limites) ». La classe, guidée par le professeur, remplit le tableau à partir de la question 3 : « L'homme est-il toujours responsable de sa minorité ? »

\begin{exemplebox}[Exemple de tableau pour / contre construit en classe]
\begin{itemize}
\item \textbf{Pour} : chacun peut lire, réfléchir, comparer les opinions ; au Cameroun, l'accès à l'école et à l'information permet de juger par soi-même ; celui qui refuse de penser choisit la facilité.
\item \textbf{Contre / limites} : la pauvreté, l'analphabétisme ou la pression sociale peuvent empêcher l'usage libre de la raison ; la responsabilité n'est donc pas toujours entière ; il faut d'abord des conditions (éducation, liberté d'expression) pour que l'appel de Kant soit réaliste.
\end{itemize}
Conclusion collective possible : la thèse de Kant garde sa valeur comme \emph{idéal} (oser penser par soi-même), mais doit être nuancée par les conditions sociales concrètes de son exercice.
\end{exemplebox}

Ce TP montre qu'une bonne discussion n'est pas un rejet de l'auteur : elle \textbf{examine} sa thèse, en reconnaît la force, puis en teste les limites à l'aide d'arguments et d'exemples.

\courssub{TP 4 : simulation d'épreuve chronométrée}

\textbf{Objectif.} Mettre l'élève en conditions réelles de Baccalauréat technique : un texte, trois questions, un temps limité, aucune aide.

\begin{methode}[Protocole de la simulation]
\begin{enumerate}
\item Distribution du sujet (texte de 10 à 15 lignes + trois questions : expliquer, dégager l'argumentation, discuter).
\item Temps imparti : \textbf{1 h 30} (proportionnel à la place de l'exercice dans l'épreuve réelle).
\item Répartition conseillée du temps : 10 min de lecture et repérage, 15 min de plan brouillon, 55 min de rédaction, 10 min de relecture.
\item Ramassage des copies, anonymisation, correction collective à la séance suivante.
\end{enumerate}
\end{methode}

\begin{exoresolu}[Question de compréhension sur le texte de Kant]
Énoncé : Expliquez la thèse de l'auteur dans l'extrait de \emph{Qu'est-ce que les Lumières ?}.
\tcblower
\textbf{Solution détaillée.} Kant affirme que les Lumières consistent pour l'homme à sortir de sa « minorité », c'est-à-dire de l'état où il est incapable d'utiliser sa propre raison sans être guidé par un autre (un prêtre, un livre, un maître). Cette minorité, précise Kant, n'est pas une fatalité : l'homme en est « lui-même responsable », car elle vient de la paresse et de la lâcheté, non d'un manque d'intelligence. La thèse est donc un appel à l'autonomie de la pensée, résumé par la devise \emph{Sapere aude !} : « Ose penser par toi-même. » En une phrase : pour Kant, être éclairé, c'est assumer le courage d'utiliser son propre entendement.
\par \textbf{Ce qui fait la qualité de cette réponse} : elle définit les termes-clés (minorité, entendement), elle reformule sans paraphraser, elle s'appuie sur le texte (« lui-même responsable ») et elle conclut par une synthèse.
\end{exoresolu}

\courssub{Grille d'autoévaluation}

Avant de rendre toute copie (en TP comme à l'examen), l'élève doit pouvoir s'évaluer lui-même. L'autoévaluation porte sur quatre dimensions : la \cle{fidélité au texte}, la \cle{pertinence des réponses}, la \cle{qualité de l'argumentation} et la \cle{présentation}.

\begin{methode}[S'autoévaluer : quatre questions à se poser]
\begin{enumerate}
\item \textbf{Ai-je répondu à la question posée ?} (et non à une question voisine que j'aurais préférée)
\item \textbf{Ai-je évité la paraphrase ?} (ai-je expliqué avec mes propres mots et des idées, au lieu de recopier le texte en le diluant ?)
\item \textbf{Ai-je justifié ?} (chaque affirmation est-elle appuyée par un argument, une référence au texte ou un exemple ?)
\item \textbf{Ai-je conclu ?} (ma réponse se termine-t-elle par une synthèse claire qui répond explicitement à la question ?)
\end{enumerate}
Si l'une des quatre réponses est « non », la copie doit être retravaillée avant d'être rendue.
\end{methode}

\courssub{Correction collective : analyse de deux copies anonymisées}

Le professeur projette ou distribue deux copies réelles anonymisées issues de la simulation : une copie réussie et une copie faible. La classe les analyse avec la grille critériée.

\begin{exemplebox}[Copie réussie (extraits commentés)]
La candidate commence par définir la minorité chez Kant, cite brièvement le texte, puis organise sa réponse en deux parties (ce que dit Kant ; pourquoi il le dit). La discussion reconnaît la force de l'appel à l'autonomie, puis le nuance par l'exemple des conditions sociales d'accès à l'éducation. \textbf{Points forts} : fidélité au texte, plan visible, justification constante, conclusion explicite. \textbf{Note indicative} : 15/20.
\end{exemplebox}

\begin{exemplebox}[Copie faible (extraits commentés)]
Le candidat recopie de longs passages du texte, ajoute « Kant a raison car il faut être intelligent », puis parle de la liberté en général sans lien avec la question. \textbf{Faiblesses repérées collectivement} : paraphrase (critère 1 non rempli), hors sujet partiel (critère 2), absence de justification (critère 3), pas de conclusion (critère 4). \textbf{Note indicative} : 6/20. \textbf{Pistes de progrès} : reformuler au lieu de copier, dégager les arguments du texte, conclure en répondant à la question.
\end{exemplebox}

\begin{aretenir}[L'essentiel des travaux pratiques]
La consolidation repose sur un cycle : entraînement, correction, autoévaluation, nouvel entraînement. La correction — croisée ou collective — porte d'abord sur la pertinence philosophique (fidélité au texte, justification, conclusion), pas seulement sur la forme. S'autoévaluer, c'est se poser quatre questions : ai-je répondu à la question ? ai-je évité la paraphrase ? ai-je justifié ? ai-je conclu ?
\end{aretenir}

\begin{savaistu}[La méthode des annales]
S'entraîner régulièrement sur des sujets des sessions antérieures du Baccalauréat camerounais est la méthode de préparation la plus efficace : les questions d'explication, d'analyse et de discussion suivent des attentes stables d'une session à l'autre. Refaire un sujet d'annale en temps limité, puis comparer sa copie aux critères de correction, fait progresser plus vite que la simple relecture du cours.
\end{savaistu}

\begin{exercice}[Pour s'entraîner à la maison]
Choisis un sujet d'exercice sur texte d'une session antérieure du Baccalauréat technique. Traite-le en 1 h 30 sans aucun document, puis autoévalue ta copie avec la grille critériée (fidélité au texte /4, pertinence /4, argumentation /6, présentation /6). Note tes deux points forts et tes deux axes de progrès, et refais le même exercice une semaine plus tard pour mesurer ta progression.
\end{exercice}

\newpage

\coursec{Activité d'intégration}

\coursec{Consolidation des acquis de la méthodologie de l'exercice sur texte}

\courssub{Objectifs de l'activité}

Cette activité d'intégration est une \cle{simulation complète} d'exercice sur texte en conditions réelles d'examen. Elle vise à faire travailler ensemble l'ensemble des savoir-faire construits dans la leçon. À l'issue de l'activité, l'élève doit être capable de :

\begin{itemize}
\item conduire une \cle{lecture méthodique} d'un texte philosophique (repérage du thème, de la thèse, de la structure argumentative) ;
\item répondre correctement aux trois types de questions du Baccalauréat technique : question de \cle{compréhension}, question d'\cle{analyse} et question de \cle{discussion} ;
\item mobiliser les notions du programme (la liberté, le devoir, l'État, la conscience) pour éclairer des \cle{situations concrètes} de la vie courante au Cameroun ;
\item \cle{rédiger et justifier} une démarche, une réponse et une conclusion personnelle ;
\item gérer son temps (1 h 30), relire sa copie et évaluer sa propre production à l'aide d'une \cle{grille critériée} afin de rédiger une fiche de progrès.
\end{itemize}

\courssub{Situation}

Au lycée technique de Nkongsamba, la classe de Terminale F4 prépare l'épreuve de Philosophie du Baccalauréat technique. Le professeur, M. Etoundi, distribue le sujet suivant, conforme au format officiel de l'examen.

\textbf{Texte (extrait adapté d'Emmanuel Kant, \emph{Réponse à la question : qu'est-ce que les Lumières ?}, 1784).}

\begin{quote}
« Les Lumières, c'est la sortie de l'homme hors de l'état de tutelle dont il est lui-même responsable. L'état de tutelle est l'incapacité à se servir de son entendement sans la conduite d'un autre. On est soi-même responsable de cet état de tutelle quand la cause tient non pas à une insuffisance de l'entendement, mais à une insuffisance de la résolution et du courage de s'en servir sans la conduite d'un autre. \emph{Sapere aude !} Aie le courage de te servir de ton propre entendement ! Voilà la devise des Lumières. La paresse et la lâcheté sont les causes qui expliquent pourquoi une si grande partie des hommes demeure volontiers toute sa vie dans l'état de tutelle, et pourquoi il est si facile à d'autres de se poser en tuteurs. Il est si commode d'être mineur ! Si j'ai un livre qui tient lieu d'entendement, un pasteur qui tient lieu de conscience, un médecin qui juge de mon régime à ma place, je n'ai pas besoin de me donner moi-même du mal. »
\end{quote}

\textbf{Questions (barème indicatif sur 20 points) :}
\begin{enumerate}
\item \textbf{Compréhension (5 points).} Expliquez ce que Kant entend par « état de tutelle » et pourquoi l'homme en est « lui-même responsable ».
\item \textbf{Analyse (7 points).} Dégagez la thèse de l'auteur et montrez comment il la justifie (arguments et exemples).
\item \textbf{Discussion (8 points).} Ce point de vue est-il encore valable aujourd'hui ? Discutez-le en vous appuyant sur des situations concrètes de la vie courante au Cameroun.
\end{enumerate}

\textbf{Données de gestion de l'épreuve :} durée totale 1 h 30, soit 90 minutes ; le professeur recommande la répartition suivante : 15 minutes de lecture et de brouillon, 60 minutes de rédaction, 15 minutes de relecture. La séance suivante (2 heures) est consacrée à la correction croisée avec grille critériée, à la correction collective au tableau et à la rédaction de la fiche de progrès personnelle.

\courssub{Consignes}

\begin{enumerate}
\item Lis le texte deux fois : une première lecture globale, puis une lecture analytique en soulignant les mots-clés. Note au brouillon le thème, le problème posé et la thèse en une phrase chacun.
\item Réponds à la question 1 en définissant « tutelle » à partir du texte, puis en expliquant la distinction entre « insuffisance de l'entendement » et « insuffisance de la résolution et du courage ». Cite brièvement le texte pour justifier chaque élément de ta réponse.
\item Réponds à la question 2 en formulant la thèse en une phrase, puis en reconstituant l'argumentation : quelles causes Kant donne-t-il à la tutelle ? Quels exemples utilise-t-il ? Que signifie la devise \emph{Sapere aude} ?
\item Réponds à la question 3 en construisant une discussion en trois moments : accorde d'abord la pertinence du point de vue (avec deux exemples camerounais précis), puis expose une difficulté ou une limite, enfin propose une conclusion personnelle nuancée et justifiée.
\item Gère ton temps selon la répartition donnée, rédige au propre, puis relis-toi en vérifiant : orthographe, cohérence, présence des citations, réponse effective à chaque question posée.
\item Lors de la séance de correction, évalue la copie de ton voisin avec la grille critériée, compare au tableau une réponse paraphrasée et une réponse expliquée, puis rédige ta fiche de progrès : trois erreurs récurrentes et trois stratégies de correction.
\end{enumerate}

\courssub{Ressources à mobiliser}

\begin{aretenir}[Les trois types de questions]
\begin{itemize}
\item \textbf{Compréhension (« Expliquez »)} : définir les termes, reformuler fidèlement la pensée de l'auteur, justifier par de courtes citations. Interdit : la paraphrase et l'opinion personnelle.
\item \textbf{Analyse (« Dégagez la thèse et l'argumentation »)} : formuler la thèse en une phrase, identifier les arguments, les exemples et la structure du raisonnement.
\item \textbf{Discussion (« Discutez »)} : dire d'abord ce qui est vrai dans la thèse (accorder), puis ce qui pose problème (critiquer avec arguments), enfin conclure par une position personnelle justifiée. Appui obligatoire sur des exemples concrets.
\end{itemize}
\end{aretenir}

\begin{aretenir}[Notions du programme utiles]
La \cle{liberté} (penser par soi-même, autonomie), la \cle{conscience} (faculté de juger du bien et du mal), le \cle{devoir} (agir par raison et non par paresse), la \cle{raison} (entendement, capacité de comprendre). Chez Kant, être majeur, c'est s'arracher à la dépendance intellectuelle par le courage de raisonner seul.
\end{aretenir}

\begin{attention}[Pièges fréquents]
Recopier le texte sans l'expliquer (paraphrase) ; donner son avis dès la question 1 ; oublier de citer le texte ; discuter sans exemples concrets ; négliger la relecture ; rédiger une discussion sans position personnelle finale.
\end{attention}

\courssub{Démarche et corrigé commenté}

\begin{exoresolu}[Simulation corrigée : le texte de Kant sur les Lumières]
\textbf{Énoncé.} Voir la situation ci-dessus : texte de Kant, trois questions (compréhension 5 pts, analyse 7 pts, discussion 8 pts), durée 1 h 30.

\tcblower

\textbf{Étape 1 — Lecture méthodique (15 min).} Au brouillon : \emph{thème} : les Lumières, c'est-à-dire l'émancipation de la pensée. \emph{Problème} : pourquoi les hommes restent-ils mineurs alors qu'ils peuvent penser par eux-mêmes ? \emph{Thèse} : la tutelle est auto-imposée ; elle vient de la paresse et de la lâcheté, non de l'incapacité ; il faut donc oser penser par soi-même. \emph{Mots-clés soulignés} : tutelle, entendement, responsable, courage, paresse, lâcheté, \emph{Sapere aude}.

\textbf{Étape 2 — Question de compréhension (réponse attendue).} Par « état de tutelle », Kant désigne « l'incapacité à se servir de son entendement sans la conduite d'un autre » : l'homme mineur laisse d'autres penser, croire et décider à sa place. Cet état est « auto-imposé » : sa cause n'est pas « une insuffisance de l'entendement » (l'homme possède la raison), mais « une insuffisance de la résolution et du courage ». L'homme est donc \emph{lui-même responsable} de sa minorité, car il pourrait en sortir et choisit, par paresse et par lâcheté, le confort de la dépendance : « Il est si commode d'être mineur ! ». \emph{Commentaire de correction : la réponse définit, distingue les deux causes et cite le texte trois fois : c'est une explication, pas une paraphrase.}

\textbf{Étape 3 — Question d'analyse (réponse attendue).} \emph{Thèse} : la sortie de la tutelle exige le courage de se servir de son propre entendement — « \emph{Sapere aude !} ». \emph{Argumentation} : Kant procède en trois temps. (a) Il définit la tutelle et montre qu'elle est auto-imposée, ce qui rend l'homme responsable de sa situation. (b) Il en donne les causes : « la paresse et la lâcheté », qui font préférer la sécurité de la dépendance à l'effort de penser. (c) Il illustre par trois exemples ironiques : le livre qui « tient lieu d'entendement », le pasteur qui « tient lieu de conscience », le médecin qui décide du régime à ma place — autant de délégations commodes de la pensée. La devise latine conclut le raisonnement comme un impératif : ose savoir. \emph{Commentaire : la thèse est formulée en une phrase, la structure est reconstituée, les exemples sont identifiés comme des arguments.}

\textbf{Étape 4 — Question de discussion (réponse attendue).} \emph{Accord} : le diagnostic de Kant reste d'une actualité frappante. Au Cameroun, combien de citoyens répètent sans vérifier les rumeurs reçues sur WhatsApp ou Facebook, laissant le « livre » numérique tenir lieu d'entendement ? Combien d'élèves attendent le corrigé du professeur sans chercher par eux-mêmes, ou de jeunes qui suivent aveuglément l'avis d'un influenceur, d'un chef traditionnel ou d'un responsable politique sans examiner les arguments ? La paresse intellectuelle décrite par Kant est bien réelle : penser coûte un effort, et la dépendance est commode. \emph{Critique} : cependant, on ne peut pas toujours penser seul. Faire confiance au médecin pour le régime ou à l'ingénieur pour la construction d'un pont n'est pas de la lâcheté : nul ne peut être expert en tout, et la vie en société exige une délégation raisonnable du savoir. De plus, sortir de la tutelle suppose des conditions : instruction, accès à l'information, liberté d'expression ; or tous n'y ont pas également accès. \emph{Conclusion personnelle} : la devise de Kant demeure valable comme exigence d'esprit critique — vérifier avant de croire, raisonner avant d'obéir — mais elle doit être tempérée : penser par soi-même ne signifie pas tout rejeter, c'est exercer son jugement sur ce que disent les autres. \emph{Commentaire : la discussion suit le plan accorder–critiquer–conclure, avec deux exemples camerounais précis et une position nuancée justifiée.}

\textbf{Étape 5 — Relecture (15 min).} Vérifier : chaque question a reçu une réponse ; les citations sont courtes et exactes ; l'orthographe des termes-clés (tutelle, entendement, \emph{Sapere aude}) ; la discussion contient bien une conclusion personnelle.

\textbf{Étape 6 — Correction croisée et fiche de progrès (séance suivante).} Grille critériée : exactitude de la compréhension (5 pts), qualité de l'analyse (thèse + structure, 7 pts), pertinence de la discussion (exemples, argumentation, conclusion, 8 pts). Au tableau, on compare : réponse paraphrasée (« Kant dit que l'homme est mineur et qu'il faut avoir du courage ») et réponse expliquée (définition + distinction des causes + citation). Chaque élève rédige sa fiche de progrès : par exemple, « erreur 1 : je paraphrase → stratégie : définir chaque terme avant de reformuler ; erreur 2 : pas d'exemples en discussion → stratégie : préparer deux exemples camerounais par notion ; erreur 3 : pas de relecture → stratégie : réserver 15 minutes systématiquement ».
\end{exoresolu}

\courssub{Bilan de l'activité}

Cette simulation a permis de mettre en évidence trois acquis essentiels. D'abord, la \cle{lecture méthodique} est la condition de tout le reste : dix minutes d'analyse au brouillon (thème, problème, thèse, structure) évitent la paraphrase et le hors-sujet. Ensuite, chaque type de question exige un travail spécifique : expliquer n'est ni recopier ni discuter ; analyser suppose de reconstituer un raisonnement ; discuter impose un plan en trois moments appuyé sur des exemples concrets — ici, la rumeur sur les réseaux sociaux et la dépendance intellectuelle des élèves, situations bien connues des candidats camerounais. Enfin, la réussite à l'exercice sur texte tient autant à la \cle{gestion de l'épreuve} (répartition du temps, relecture) qu'à la connaissance des notions. La correction croisée et la fiche de progrès transforment l'erreur en outil d'apprentissage : chaque élève repart avec un diagnostic personnel et des stratégies concrètes, prêt pour l'épreuve du Baccalauréat technique.

\newpage

\coursec{Méthodes et exercices résolus}

\coursec{Consolidation des acquis de la méthodologie de l'exercice sur texte}

\courssub{Méthode 1 : Conduire la lecture méthodique du texte}

\begin{methode}[Lire méthodiquement un texte philosophique]
\begin{enumerate}
\item \textbf{Identifier le texte} : relever l'auteur, l'œuvre, le courant de pensée, le contexte historique. Cela éclaire le propos sans jamais s'y substituer.
\item \textbf{Première lecture globale} : lire le texte en entier, sans stylo, pour saisir le sens général et la thèse principale.
\item \textbf{Deuxième lecture analytique} : souligner les mots-clés, repérer la thèse (ce que l'auteur affirme), les arguments (ce qui la justifie) et les exemples (ce qui l'illustre).
\item \textbf{Dégager la structure} : découper le texte en moments (généralement 2 ou 3) et formuler l'idée directrice de chaque moment en une phrase complète.
\item \textbf{Formuler la problématique} : à quelle question philosophique le texte répond-il ? L'écrire sous forme interrogative précise.
\item \textbf{Définir les termes difficiles} : tout terme technique doit être expliqué avec ses propres mots, jamais par simple répétition du dictionnaire.
\end{enumerate}
\textbf{Réflexe à retenir} : Thèse + Arguments + Exemples = le squelette de tout texte. On répond toujours \emph{à partir du texte}, jamais à partir de ses seuls souvenirs de cours.
\end{methode}

\begin{exoresolu}[Lecture méthodique d'un extrait de Kant]
\textbf{Énoncé.} On donne l'extrait suivant : « L'opinion est un savoir qui a conscience d'être insuffisant. La croyance est un savoir qui se tient pour suffisant, sans l'être objectivement. La science seule est un savoir certain. » (Kant, \emph{Critique de la raison pure}, adapté). Après une lecture méthodique, dégagez la thèse, la structure et la problématique de ce texte.
\tcblower
\textbf{Solution détaillée.}
\begin{enumerate}
\item \textbf{Identification} : il s'agit d'un texte d'Emmanuel Kant, philosophe des Lumières allemand, extrait de la \emph{Critique de la raison pure} (Doctrine de la méthode), portant sur la théorie de la connaissance et la critique du dogmatisme.
\item \textbf{Mots-clés soulignés} : \cle{opinion}, \cle{croyance}, \cle{science}, \cle{savoir certain}, \cle{subjectif}/\cle{objectif}.
\item \textbf{Thèse} : il existe une hiérarchie des formes de savoir selon leur degré de certitude objective ; seule la science, par sa méthode rigoureuse, atteint le savoir certain.
\item \textbf{Structure} : le texte comporte trois moments correspondant aux trois phrases : (a) définition de l'opinion comme savoir conscient de son insuffisance (subjectivité assumée) ; (b) définition de la croyance comme savoir qui se croit suffisant à tort (subjectivité non assumée, illusion) ; (c) affirmation de la science comme seul savoir certain (validité objective).
\item \textbf{Problématique} : qu'est-ce qui distingue le vrai savoir (science) de ses simples apparences (opinion, croyance) ? Ou : comment hiérarchiser les degrés de la certitude ?
\item \textbf{Définitions} : l'opinion est un jugement subjectif que l'on sait pouvoir être faux ; la croyance est un jugement subjectif que l'on tient à tort pour universellement valide ; la science est un savoir démontré, vérifiable et objectivement certain.
\end{enumerate}
\textbf{Conclusion} : le texte oppose trois degrés du prétendu savoir pour réserver la certitude à la seule science ; toute réponse aux questions devra s'appuyer sur cette structure ternaire et la distinction subjectif/objectif.
\end{exoresolu}

\courssub{Méthode 2 : Répondre aux questions de compréhension}

\begin{methode}[Traiter les questions de compréhension]
\begin{enumerate}
\item \textbf{Lire la question deux fois} et repérer le verbe d'injonction : « définissez », « relevez », « expliquez », « quelle est la thèse de l'auteur ? », « en quoi consiste... ? ».
\item \textbf{Localiser la réponse dans le texte} : la réponse à une question de compréhension se trouve toujours dans le texte, explicitement ou implicitement (par inférence immédiate).
\item \textbf{Citer le texte} : appuyer chaque réponse sur une citation courte entre guillemets, suivie d'une explication en ses propres termes (reformulation).
\item \textbf{Rédiger en phrases complètes} : jamais de réponse télégraphique ni de simple recopie du texte sans commentaire personnel.
\item \textbf{Rester fidèle à l'auteur} : dire ce que le texte dit, pas ce qu'on voudrait qu'il dise ; aucune opinion personnelle ni connaissance extérieure à ce niveau.
\end{enumerate}
\textbf{Formule à retenir} : réponse = citation précise + reformulation explicative + lien avec la thèse générale.
\end{methode}

\begin{exoresolu}[Question de compréhension type Bac]
\textbf{Énoncé.} Texte : « La liberté ne consiste pas à faire ce que l'on veut, mais à vouloir ce que l'on doit. L'homme qui obéit à ses passions est un esclave ; celui qui obéit à la raison seule est libre. » Question : quelle conception de la liberté l'auteur défend-il ? Appuyez votre réponse sur le texte.
\tcblower
\textbf{Solution détaillée.}
\begin{enumerate}
\item \textbf{Analyse de la consigne} : le verbe « défend-il » exige de dégager la thèse de l'auteur sur la liberté, sans jugement personnel.
\item \textbf{Localisation} : la thèse est dans la première phrase (définition par opposition) ; la justification dans la seconde (argumentation par l'opposition esclave/libre).
\item \textbf{Rédaction} : L'auteur défend une conception \cle{rationaliste} (ou \cle{autonome}) de la liberté. Il affirme que « la liberté ne consiste pas à faire ce que l'on veut, mais à vouloir ce que l'on doit » : la vraie liberté n'est donc pas l'absence de contrainte (indépendance), mais l'obéissance à la raison et au devoir (autonomie). Il le justifie par une opposition : « l'homme qui obéit à ses passions est un esclave », car celui qui suit ses désirs est déterminé par des causes extérieures à sa volonté ; en revanche, « celui qui obéit à la raison seule est libre », car il se donne à lui-même sa propre loi (autonomie).
\item \textbf{Vérification} : la réponse cite le texte, le reformule et ne contient aucune opinion personnelle.
\end{enumerate}
\textbf{Conclusion} : pour l'auteur, être libre, c'est s'affranchir de la tyrannie des passions en suivant la raison : la liberté est autonomie, non licence.
\end{exoresolu}

\courssub{Méthode 3 : Répondre aux questions d'analyse}

\begin{methode}[Traiter les questions d'analyse]
\begin{enumerate}
\item \textbf{Repérer l'objet de l'analyse} : structure logique, procédé d'argumentation, sens précis d'un terme dans le texte, présupposé implicite.
\item \textbf{Découper} : pour la structure, diviser le texte en moments et justifier chaque coupure par un indice textuel (connecteur logique : « mais », « donc », « cependant », « or », « car » ; changement de thème ou de ton).
\item \textbf{Nommer les procédés} : définition, exemple, comparaison, analogie, opposition, concession, argument d'autorité, raisonnement par l'absurde, réduction à l'absurde, dilemme.
\item \textbf{Expliquer la fonction} : à quoi sert ce procédé dans la démonstration de la thèse ? (ex: la concession désarme l'adversaire, l'exemple rend concret, l'opposition clarifie).
\item \textbf{Synthétiser} : montrer comment l'analyse confirme la cohérence d'ensemble du texte et la progression de la pensée.
\end{enumerate}
\textbf{Réflexe} : une analyse sans justification textuelle (citation du connecteur, mot-clé repéré) est une affirmation gratuite qui ne rapporte pas de points.
\end{methode}

\begin{exoresolu}[Analyse de la structure d'un texte]
\textbf{Énoncé.} Texte : « Tout le monde reconnaît que le travail est pénible. Cependant, c'est par le travail que l'homme transforme la nature et se transforme lui-même. Le travail est donc, malgré sa peine, la condition de l'humanisation de l'homme. » Dégagez la structure logique de ce texte et montrez son fonctionnement argumentatif.
\tcblower
\textbf{Solution détaillée.}
\begin{enumerate}
\item \textbf{Repérage des connecteurs} : « cependant » (opposition/concession) et « donc » (conséquence/conclusion) découpent le texte en trois moments logiques.
\item \textbf{Moment 1} (phrase 1) : \cle{concession} à l'opinion commune — l'auteur accorde que le travail est pénible (« tout le monde reconnaît que... »). Fonction : rejoindre le lecteur sur un fait évident pour mieux le faire basculer.
\item \textbf{Moment 2} (phrase 2) : \cle{renversement} / argument central — le connecteur « cependant » introduit la thèse positive : le travail a une double fonction ontologique (transformer la nature) et anthropologique (transformer l'homme lui-même).
\item \textbf{Moment 3} (phrase 3) : \cle{conclusion} — « donc » marque la thèse finale : le travail est la condition de l'humanisation. La reprise « malgré sa peine » montre que la conclusion intègre la concession sans l'annuler.
\item \textbf{Fonctionnement} : le texte procède par \cle{concession} puis \cle{réfutation implicite} de l'opinion commune réduite à son aspect immédiat (la peine). La structure concession–argument–conclusion rend la démonstration progressive et convaincante.
\item \textbf{Synthèse} : la structure ternaire conduit du constat commun (peine du travail) à la thèse philosophique (valeur humanisante du travail) par un pivot argumentatif fort (« cependant »).
\end{enumerate}
\textbf{Conclusion} : le texte est construit en trois moments articulés par « cependant » et « donc », conduisant du sensible (peine) à l'essentiel (humanisation).
\end{exoresolu}

\courssub{Méthode 4 : Répondre aux questions de discussion}

\begin{methode}[Discuter une thèse avec méthode]
\begin{enumerate}
\item \textbf{Reformuler la thèse à discuter} en une phrase claire, pour montrer qu'on l'a comprise et qu'on en a saisi la portée.
\item \textbf{Dégager l'enjeu} : pourquoi cette question est-elle importante pour la vie, la société ou la pensée ? (ex: fondement du droit, de la morale, de la connaissance).
\item \textbf{Examiner la pertinence (arguments \emph{pour})} : présenter d'abord ce qui rend la thèse défendable, ses forces, ses arguments internes, en citant brièvement un auteur ou une référence de cours si pertinent.
\item \textbf{Confronter (arguments \emph{contre} / nuances)} : opposer une objection solide, un contre-exemple, une limite, ou une thèse rivale (en citant brièvement un auteur adverse si possible).
\item \textbf{Prendre position de manière justifiée} : dire clairement si l'on adhère, si l'on nuance (distinction, condition) ou si l'on rejette, et \emph{pourquoi} (argument décisif).
\item \textbf{Conclure} par une réponse tranchée à la question posée, qui résume la position adoptée.
\end{enumerate}
\textbf{Formule à retenir} : discussion = thèse reformulée + enjeu + arguments pour + arguments contre/nuances + position personnelle justifiée + conclusion. Une discussion sans position personnelle argumentée est inachevée.
\end{methode}

\begin{exoresolu}[Question de discussion type Bac]
\textbf{Énoncé.} « La conscience fait la grandeur de l'homme. » Discutez cette thèse en vous appuyant sur des arguments philosophiques et des exemples concrets.
\tcblower
\textbf{Solution détaillée.}
\begin{enumerate}
\item \textbf{Reformulation} : cette thèse affirme que ce qui élève l'homme au-dessus des autres êtres vivants, c'est sa conscience, entendue comme capacité de réflexivité (se connaître soi-même), de délibération morale et de jugement sur ses propres actes.
\item \textbf{Enjeu} : il s'agit de définir la dignité spécifique de l'homme, fondement de la morale, du droit et de la responsabilité juridique.
\item \textbf{Pertinence de la thèse} : la conscience permet la délibération morale et la responsabilité. Grâce à elle, l'homme peut distinguer le bien du mal et se tenir pour auteur de ses actes : un élève de Yaoundé qui avoue sa faute au lieu d'accuser un camarade agit moralement parce qu'il s'examine lui-même (conscience morale). Descartes fait de la pensée (« je pense, donc je suis ») le fondement indubitable de l'existence humaine ; Pascal dit que « l'homme n'est qu'un roseau, le plus faible de la nature, mais c'est un roseau pensant ».
\item \textbf{Objection / Nuance} : la conscience peut aussi être source de faiblesse, de malheur et de perversion : elle rend l'homme capable de remords, d'angoisse existentielle, et même de \cle{mauvaise foi} (Sartre) ou de lucidité cynique. Certains actes criminels sont commis en pleine conscience (prémeditation) : la conscience ne garantit donc pas la grandeur \emph{morale}, elle n'est que la condition de possibilité du bien comme du mal.
\item \textbf{Position personnelle} : nous admettons la thèse, mais nuancée : la conscience fait la grandeur de l'homme non pas automatiquement (comme un attribut statique), mais \emph{lorsqu'elle est exercée} comme conscience morale, éclairée par la raison et tournée vers le bien. C'est l'usage droit de la conscience (l'examen de soi, le repentir, le choix du bien) qui fonde la dignité, non sa simple possession psychologique.
\item \textbf{Conclusion} : la conscience est bien ce qui distingue ontologiquement l'homme (grandeur métaphysique), à condition qu'elle soit actualisée comme conscience morale et responsable (grandeur éthique).
\end{enumerate}
\end{exoresolu}

\courssub{Méthode 5 : Appliquer les notions philosophiques à des situations concrètes}

\begin{methode}[Mobiliser une notion dans une situation de la vie courante]
\begin{enumerate}
\item \textbf{Rappeler la définition rigoureuse} de la notion concernée (liberté, devoir, État, travail, vérité, justice, conscience, technique, religion, art, langage, histoire, nature, culture, politique...).
\item \textbf{Décrire la situation concrète} avec précision : qui ? quoi ? où ? dans quelles conditions ? (contexte camerounais privilégié : école, marché, quartier, élections, réseaux sociaux, entreprise, famille, administration).
\item \textbf{Établir le lien conceptuel} : montrer comment la notion éclaire la situation (analyse) ou comment la situation illustre la notion (exemplification). Utiliser les distinctions clés du cours (ex: liberté/licence, légalité/légitimité, travail/emploi, droit/devoir).
\item \textbf{Tirer la leçon philosophique} : quelle conduite, quel jugement, quelle critique ou quelle décision la notion permet-elle de fonder dans cette situation ?
\end{enumerate}
\textbf{Réflexe} : privilégier des exemples camerounais précis et datés si possible (ex: « lors des élections municipales de 2020 à Douala... », « dans un lycée technique de Bafoussam... ») plutôt que des exemples vagues (« un homme », « dans une société »).
\end{methode}

\begin{exoresolu}[Application concrète : la liberté et la règle]
\textbf{Énoncé.} À l'aide de la notion de liberté définie comme « obéissance à la loi qu'on s'est prescrite » (Rousseau, \emph{Du contrat social}), analysez la situation suivante : dans un lycée technique de Douala, le règlement intérieur, discuté et adopté avec les délégués de classe, interdit l'usage du téléphone en classe. Un élève respecte cette règle tout en souhaitant utiliser son téléphone. Est-il libre ?
\tcblower
\textbf{Solution détaillée.}
\begin{enumerate}
\item \textbf{Définition} : pour Rousseau, la liberté civile n'est pas l'indépendance naturelle (faire ce qu'on veut), mais l'obéissance à la loi qu'on s'est prescrite collectivement : « l'impulsion du seul appétit est esclavage, et l'obéissance à la loi qu'on s'est prescrite est liberté ». Obéir à une règle à laquelle on a participé comme citoyen, c'est obéir à soi-même (autonomie).
\item \textbf{Description de la situation} : la règle n'a pas été imposée arbitrairement par l'administration seule ; elle a été discutée et adoptée avec les représentants légitimes des élèves (délégués). L'élève, par le biais de ses délégués, a participé à l'élaboration de la norme.
\item \textbf{Lien notion–situation} : en respectant le règlement, l'élève n'obéit pas à une volonté étrangère (hétéronomie), mais à une norme qu'il a contribué à établir comme membre du corps citoyen de l'établissement. Sa frustration immédiate (envie du téléphone) relève du \cle{désir} ou de l'inclination sensible, non de la liberté : céder au désir serait, au sens de Kant, une \cle{servitude} (hétéronomie) envers ses inclinations.
\item \textbf{Leçon philosophique} : l'élève est donc libre \emph{par son respect même de la règle}, car la contrainte qu'il subit est une \cle{autocontrainte} rationnelle, consentie collectivement. La situation montre que la règle légitime (adoptée démocratiquement) ne détruit pas la liberté : elle la réalise et la protège contre l'arbitraire des désirs immédiats.
\end{enumerate}
\textbf{Conclusion} : l'élève qui respecte une règle à l'élaboration de laquelle il a participé exerce sa liberté au sens fort (autonomie politique et morale) : la liberté véritable est autonomie, non caprice ; la loi commune est la condition de la liberté de chacun.
\end{exoresolu}

\courssub{Méthode 6 : Rédiger, justifier sa démarche et gérer l'épreuve}

\begin{methode}[Organiser la copie et le temps de l'épreuve]
\begin{enumerate}
\item \textbf{Analyser le sujet et le barème} (5 min) : lire le texte et toutes les questions ; repérer leur nature (compréhension, analyse, discussion) et leur pondération (points).
\item \textbf{Répartir le temps} : pour une épreuve de 2 h, environ 20 min de lecture méthodique, prise de notes et brouillon structuré ; 80 min de rédaction soignée ; 20 min de relecture active.
\item \textbf{Répondre dans l'ordre des questions} et numéroter clairement chaque réponse (1., 2., 3.) en reprenant les termes de la question.
\item \textbf{Justifier chaque réponse} : toute affirmation doit être suivie d'une preuve (citation du texte entre guillemets, argument logique, référence d'auteur, exemple concret).
\item \textbf{Soigner la langue} : phrases complètes, syntaxe correcte, orthographe soignée, connecteurs logiques variés (« d'abord », « ensuite », « cependant », « par conséquent », « en définitive »).
\item \textbf{Relire activement} : vérifier qu'aucune question n'est oubliée, que chaque réponse correspond bien au verbe de la consigne, que les citations sont exactes, que l'écriture est lisible.
\end{enumerate}
\textbf{Réflexe} : une réponse juste mais non justifiée ne vaut que la moitié des points ; une copie non relue en perd toujours (erreurs de lecture, oublis, contresens).
\end{methode}

\begin{exoresolu}[Sujet complet d'entraînement corrigé]
\textbf{Énoncé.} Texte : « L'État ne se résume pas à la force. Certes, il dispose de la contrainte ; mais une autorité qui ne repose que sur la force ne dure pas. L'État légitime est celui qui assure la justice et la sécurité de tous. » Questions : (1) Quelle est la thèse du texte ? (2) Dégagez sa structure. (3) La force suffit-elle à fonder l'État ? Discutez.
\tcblower
\textbf{Solution détaillée.}
\begin{enumerate}
\item \textbf{Question 1 (compréhension)} : la thèse est que la légitimité de l'État ne repose pas sur la seule force, mais sur la justice et la sécurité assurées à tous les citoyens. L'auteur affirme explicitement : « L'État légitime est celui qui assure la justice et la sécurité de tous. » La force (contrainte) est un attribut, pas le fondement.
\item \textbf{Question 2 (analyse)} : le texte comporte trois moments. Moment 1 (phrase 1) : thèse annoncée négativement — « L'État ne se résume pas à la force » (négation de l'identification État/force). Moment 2 (phrases 2 et 3) : argumentation par \cle{concession} (« certes, il dispose de la contrainte ») suivie d'une \cle{réfutation} par l'expérience historique (« une autorité qui ne repose que sur la force ne dure pas »). Moment 3 (phrase 4) : conclusion positive définissant l'État légitime par sa finalité (justice, sécurité). Les connecteurs « certes... mais » organisent le passage de la concession à la réfutation ; « mais » marque l'opposition centrale.
\item \textbf{Question 3 (discussion)} :
\begin{itemize}
\item \emph{Pour la thèse (force insuffisante)} : un pouvoir purement fondé sur la force physique provoque la révolte et l'instabilité ; l'histoire montre que les dictatures fondées sur la terreur finissent par s'effondrer (ex: chute de régimes autoritaires). Rousseau écrit dans \emph{Du contrat social} : « La force ne fait pas le droit, et on n'obéit qu'aux pouvoirs légitimes » : la force crée la soumission (contrainte extérieure), non l'obéissance (reconnaissance intérieure).
\item \emph{Nuance (force nécessaire)} : la force (monopole de la violence légitime, police, justice contraignante) reste un \emph{moyen} indispensable. Sans contrainte étatique, les lois seraient lettre morte ; un État sans moyens de coercition ne peut garantir la sécurité (ex: zones de non-droit, insécurité dans certaines régions en crise).
\item \emph{Position} : la force ne suffit pas à \emph{fonder} l'État (fondement = légitimité = justice/consentement), mais elle est un \emph{instrument} nécessaire au service de la justice ; c'est la finalité (justice, sécurité de tous) qui légitime l'usage de la force, et non l'inverse. L'État de droit conjugue \cle{droit} (fondement) et \cle{force} (moyen).
\end{itemize}
\end{enumerate}
\textbf{Conclusion} : la force est une condition de l'efficacité de l'État, mais seule la justice en fonde la légitimité ; l'État véritable conjugue le droit et la contrainte au service de tous.
\end{exoresolu}

\begin{attention}[Les 5 erreurs qui coûtent des points au Bac]
\begin{enumerate}
\item \textbf{Répondre hors texte} : réciter son cours ou donner son opinion personnelle alors que la question porte sur ce que dit l'auteur (compréhension, analyse). Toute réponse de compréhension doit s'appuyer sur une citation textuelle explicite.
\item \textbf{Recopier le texte sans l'expliquer} : une citation non commentée, non reformulée en ses propres termes, ne rapporte aucun point ; il faut toujours expliciter le sens de la citation.
\item \textbf{Confondre compréhension et discussion} : donner son avis personnel dans une question de compréhension (« Expliquez ce que veut dire l'auteur... »), ou oublier de prendre position claire et justifiée dans une question de discussion (« Discutez... », « Qu'en pensez-vous ? »). Repérer le verbe de la consigne avant de rédiger.
\item \textbf{Découper le texte sans le justifier} : annoncer des « moments » ou des « parties » sans citer les connecteurs logiques (« mais », « donc », « cependant », « or », « car ») ou les indices thématiques qui fondent objectivement la coupure.
\item \textbf{Négliger la présentation et la gestion du temps} : réponses non numérotées, phrases incomplètes, orthographe négligée, absence de relecture, dernière question bâclée faute de temps. Prévoir impérativement 15 à 20 minutes de relecture finale.
\end{enumerate}
\end{attention}

\newpage

\coursec{Exercices}

\courssub{Je vérifie mes connaissances}

\begin{exercice}[QCM : les fondamentaux de l'exercice sur texte]
Pour chaque question, une seule réponse est correcte. Entoure la lettre correspondante.

\textbf{1.} La première étape du travail sur un texte philosophique est :
\begin{itemize}
\item[a)] rédiger immédiatement la discussion personnelle ;
\item[b)] lire attentivement le texte et identifier l'auteur, le thème et la thèse ;
\item[c)] recopier le texte pour bien le mémoriser ;
\item[d)] chercher les citations d'autres philosophes.
\end{itemize}

\textbf{2.} La thèse d'un texte est :
\begin{itemize}
\item[a)] le résumé mot à mot du texte ;
\item[b)] la position défendue par l'auteur, la réponse qu'il apporte au problème posé ;
\item[c)] l'opinion personnelle du candidat ;
\item[d)] le titre de l'ouvrage dont le texte est extrait.
\end{itemize}

\textbf{3.} Expliquer un texte, c'est :
\begin{itemize}
\item[a)] le paraphraser en le répétant avec d'autres mots ;
\item[b)] dégager le sens des expressions clés et montrer comment l'auteur construit sa démonstration ;
\item[c)] donner son avis sur le sujet ;
\item[d)] critiquer l'auteur dès le début de la copie.
\end{itemize}

\textbf{4.} Dans la question de discussion, le candidat doit :
\begin{itemize}
\item[a)] répéter la thèse de l'auteur sans la discuter ;
\item[b)] prendre clairement position, argumenter et illustrer par des exemples pertinents ;
\item[c)] refuser systématiquement la thèse de l'auteur ;
\item[d)] éviter de conclure pour ne pas se tromper.
\end{itemize}
\end{exercice}

\begin{exercice}[Vrai ou faux — justifie ta réponse]
Indique si chaque affirmation est vraie ou fausse, puis justifie en une ou deux phrases.

\textbf{1.} Il est interdit de numéroter ses réponses dans un exercice sur texte.

\textbf{2.} La paraphrase du texte est une explication philosophique acceptable à l'examen.

\textbf{3.} Un exemple tiré de la vie courante camerounaise peut renforcer une argumentation philosophique.

\textbf{4.} La présentation du texte (auteur, œuvre, thème, problème, thèse) doit précéder les réponses aux questions.

\textbf{5.} Dans une discussion, il suffit d'annoncer son avis sans le justifier.
\end{exercice}

\begin{exercice}[Définitions express]
Définis en deux ou trois phrases chacune des notions suivantes, telles qu'elles s'emploient dans la méthodologie de l'exercice sur texte :
\begin{enumerate}
\item la \cle{thèse} ;
\item l'\cle{argumentation} ;
\item l'\cle{explication de texte} ;
\item la \cle{discussion} ;
\item la \cle{problématique}.
\end{enumerate}
\end{exercice}

\begin{exercice}[Repérage dans un texte]
On donne l'extrait suivant de Kant, \emph{Qu'est-ce que les Lumières ?} : « Les Lumières, c'est la sortie de l'homme hors de l'état de minorité dont il est lui-même responsable. [...] \emph{Sapere aude !} Aie le courage de te servir de ton propre entendement ! Voilà la devise des Lumières. »

Réponds en deux phrases maximum à chaque question :
\begin{enumerate}
\item Qui est l'auteur du texte et de quelle œuvre est-il extrait ?
\item Quel est le thème du texte ?
\item Quelle est la thèse défendue par l'auteur ?
\item Relevez dans le texte un argument qui soutient cette thèse.
\end{enumerate}
\end{exercice}

\courssub{Je m'entraîne}

\begin{exercice}[De la paraphrase à l'explication]
On donne la phrase suivante, extraite d'un texte de Platon : « Les prisonniers de la caverne prennent les ombres pour la réalité. »

\textbf{1.} On propose trois reformulations. Indique laquelle est une simple paraphrase, laquelle est un hors-sujet et laquelle est une véritable explication philosophique :
\begin{itemize}
\item[a)] « Les hommes enchaînés dans la grotte croient que les reflets qu'ils voient sont les vraies choses. »
\item[b)] « Platon montre ici que l'opinion commune est une connaissance trompeuse : ce que nous percevons immédiatement n'est que l'apparence du réel, et non le réel lui-même. »
\item[c)] « Les grottes du Cameroun, comme celles de l'Ouest, abritent de nombreuses chauves-souris. »
\end{itemize}

\textbf{2.} Rédige toi-même une explication philosophique (trois à cinq lignes) de l'expression : « sortir de la caverne ».
\end{exercice}

\begin{exercice}[Construire une réponse de compréhension]
Texte support : « La technique n'est pas la science appliquée : elle est aussi ancienne que l'homme, car fabriquer un outil, c'est déjà transformer la nature selon un projet. »

Rédige, en cinq à huit lignes, la réponse à la question suivante, en respectant la méthode (reformulation de la question, réponse argumentée, appui sur le texte) : « Expliquez la distinction que l'auteur établit entre technique et science. »
\end{exercice}

\begin{exercice}[Choisir des exemples pertinents]
Pour chacune des thèses suivantes, propose un exemple concret tiré de la vie courante camerounaise qui l'illustre correctement, et explique en une phrase pourquoi cet exemple est pertinent.

\textbf{1.} « L'éducation libère l'homme de l'ignorance. »

\textbf{2.} « La technique peut aliéner le travailleur en le réduisant à un geste répétitif. »

\textbf{3.} « Penser par soi-même exige du courage. »
\end{exercice}

\begin{exercice}[La copie fautive : repérage d'erreurs méthodologiques]
Voici des extraits de la copie d'un élève répondant à un exercice sur un texte de Kant. Cette copie contient \textbf{cinq erreurs méthodologiques}. Identifie chaque erreur, nomme-la (paraphrase, hors-sujet, absence de position personnelle, exemple non pertinent, réponses non numérotées) et propose une correction.

« Kant dit que l'homme doit sortir de la minorité dont il est lui-même responsable, c'est-à-dire que l'homme est responsable de sa minorité dont il doit sortir. [...] D'ailleurs, mon quartier à Yaoundé est très animé le soir et les jeunes aiment le football, ce qui prouve que la liberté est importante. [...] Kant a raison, c'est évident, tout le monde le sait. [...] (l'élève enchaîne ensuite toutes ses réponses à la suite, sans numérotation ni présentation du texte, et ne répond jamais à la question de discussion posée par le sujet, préférant raconter la vie de Kant). »
\end{exercice}

\courssub{Je me prépare au Bac}

\begin{exercice}[Sujet type Baccalauréat technique n\textsuperscript{o} 1 — Kant, \emph{Qu'est-ce que les Lumières ?}]
\textbf{Texte :} « Les Lumières, c'est la sortie de l'homme hors de l'état de minorité dont il est lui-même responsable. La minorité, c'est l'incapacité de se servir de son entendement sans la direction d'autrui ; minorité dont l'homme est lui-même responsable, puisque la cause en réside non dans un défaut de l'entendement, mais dans un défaut de décision et de courage de s'en servir sans la direction d'autrui. \emph{Sapere aude !} Aie le courage de te servir de ton propre entendement ! Voilà donc la devise des Lumières. La paresse et la lâcheté sont les causes qui expliquent qu'un si grand nombre d'hommes [...] demeurent volontiers mineurs leur vie durant. » (Kant, \emph{Réponse à la question : qu'est-ce que les Lumières ?}, 1784.)

\textbf{Questions :}
\begin{enumerate}
\item Expliquez l'expression : « minorité dont l'homme est lui-même responsable ». (4 pts)
\item Dégagez la thèse du texte et montrez les étapes de son argumentation. (6 pts)
\item L'éducation suffit-elle à rendre l'homme autonome ? Discutez cette question en vous appuyant sur le texte et en illustrant votre réponse par le fonctionnement du système éducatif camerounais (école, famille, médias, réseaux sociaux). (10 pts)
\end{enumerate}
\emph{Durée conseillée : 2 h. N'oubliez pas de présenter le texte (auteur, œuvre, thème, problème, thèse) avant de répondre aux questions.}
\end{exercice}

\begin{exercice}[Sujet type Baccalauréat technique n\textsuperscript{o} 2 — Platon, l'allégorie de la caverne]
\textbf{Texte :} « Figure-toi des hommes dans une demeure souterraine, en forme de caverne. Depuis l'enfance, ils ont les jambes et le cou enchaînés, de sorte qu'ils ne peuvent voir que devant eux. Derrière eux brûle un feu, et entre le feu et les prisonniers passe un chemin où des hommes transportent des objets dont les ombres se projettent sur la paroi. Les prisonniers prennent ces ombres pour la réalité. Si l'un d'eux était délivré et contraint de monter vers la lumière, il souffrirait d'abord ; puis, habitué peu à peu, il verrait les choses elles-mêmes, et enfin le soleil. » (Adapté de Platon, \emph{La République}, livre VII.)

\textbf{Questions :}
\begin{enumerate}
\item Expliquez le sens philosophique de l'allégorie : que représentent la caverne, les chaînes, les ombres et la sortie vers la lumière ? (5 pts)
\item Dégagez la thèse du texte et les étapes de l'argumentation de Platon. (5 pts)
\item Peut-on sortir seul de l'ignorance, ou avons-nous besoin d'un maître ? Discutez en prenant clairement position et en illustrant votre propos par des situations concrètes de la vie scolaire et sociale au Cameroun. (10 pts)
\end{enumerate}
\emph{Durée conseillée : 2 h.}
\end{exercice}

\begin{exercice}[Sujet type Baccalauréat technique n\textsuperscript{o} 3 — Technique et travail + entraînement chronométré]
\textbf{Texte :} « La machine ne supprime pas le travail, elle le transforme. L'ouvrier qui répète mille fois le même geste ne fabrique plus un objet : il sert un rythme. Mais la même technique qui enchaîne le geste peut aussi délivrer l'homme des tâches pénibles et lui ouvrir le temps de penser. Tout dépend de l'usage que la société fait de ses outils. » (Texte d'inspiration philosophique sur la technique.)

\textbf{Partie A — Exercice sur texte (1 h 15) :}
\begin{enumerate}
\item Expliquez la distinction que le texte établit entre la technique qui aliène et la technique qui libère. (4 pts)
\item Dégagez la thèse du texte et son argument principal. (4 pts)
\item La technique libère-t-elle ou aliène-t-elle le travailleur ? Discutez en illustrant votre réponse par la situation des entreprises industrielles de Douala (usines, brasseries, port autonome, ateliers de fabrication). (8 pts)
\end{enumerate}

\textbf{Partie B — Entraînement chronométré (45 min) :}
Sur le même texte, rédigez \textbf{uniquement} :
\begin{enumerate}
\item la présentation complète du texte (auteur supposé, thème, problème, thèse) en cinq lignes maximum ;
\item le plan détaillé de votre réponse à la question 3 (introduction de la discussion, deux arguments avec exemples, objection, réponse à l'objection, conclusion).
\end{enumerate}
\emph{Objectif : apprendre à gérer le temps de l'épreuve et à rédiger une démarche justifiée.}
\end{exercice}

\newpage

\coursec{Corrigés des exercices}

\coursec{Consolidation des acquis de la méthodologie de l'exercice sur texte}

\courssub{Exercices de révision méthodologique}

\begin{corrige}[Exercice 1 — QCM : les fondamentaux de l'exercice sur texte]
\textbf{1. Réponse b).} Avant toute rédaction, le candidat doit lire attentivement le texte (au moins deux lectures : une lecture globale, puis une lecture analytique) et identifier l'auteur, l'œuvre, le thème, le problème et la thèse. Rédiger la discussion d'emblée (a) revient à ignorer le texte ; recopier (c) est une perte de temps ; chercher des citations (d) ne vient qu'après la compréhension du texte.

\textbf{2. Réponse b).} La thèse est la position que l'auteur défend : c'est sa réponse au problème posé. Un résumé mot à mot (a) serait une paraphrase ; l'opinion du candidat (c) relève de la discussion, pas de l'explication ; le titre de l'ouvrage (d) n'est qu'une indication matérielle.

\textbf{3. Réponse b).} Expliquer, c'est dégager le sens des expressions clés, définir les concepts et montrer l'enchaînement logique de la démonstration. La paraphrase (a) est la faute méthodologique la plus sanctionnée ; l'avis personnel (c) et la critique (d) appartiennent à la discussion, qui vient après l'explication.

\textbf{4. Réponse b).} La question de discussion exige une prise de position claire, soutenue par des arguments et illustrée par des exemples pertinents. Répéter la thèse (a) n'est pas discuter ; refuser systématiquement (c) est un dogmatisme inverse ; une discussion sans conclusion (d) est une copie inachevée.
\end{corrige}

\begin{corrige}[Exercice 2 — Vrai ou faux]
\textbf{1. FAUX.} Il est au contraire \emph{obligatoire} de numéroter ses réponses en reprenant la numérotation du sujet : le correcteur doit pouvoir vérifier immédiatement que chaque question a été traitée.

\textbf{2. FAUX.} La paraphrase se contente de répéter le texte avec d'autres mots, sans en dégager le sens philosophique. L'explication exige de définir les concepts, de reformuler la thèse et de montrer la structure de l'argumentation.

\textbf{3. VRAI.} Un exemple concret (vie scolaire, travail, vie sociale au Cameroun) rend l'argumentation vivante et prouve que le candidat sait appliquer les notions philosophiques à des situations réelles, ce qui est explicitement exigé par le programme.

\textbf{4. VRAI.} La présentation du texte (auteur, œuvre, thème, problème, thèse) est l'introduction obligatoire de l'exercice : elle situe le texte et montre que le candidat l'a compris avant de répondre aux questions.

\textbf{5. FAUX.} Une opinion non justifiée n'a aucune valeur philosophique. Toute prise de position doit être soutenue par au moins deux arguments, des exemples, et idéalement la discussion d'une objection.
\end{corrige}

\begin{corrige}[Exercice 3 — Définitions express]
\textbf{1. La thèse.} La thèse est la position défendue par l'auteur du texte : c'est la réponse qu'il apporte au problème philosophique posé. Elle s'énonce en une phrase claire, du type : « Pour l'auteur, ... ».

\textbf{2. L'argumentation.} L'argumentation est l'ensemble des raisons et des preuves qu'un auteur (ou le candidat) enchaîne logiquement pour établir sa thèse et convaincre le lecteur. Elle repose sur des arguments (raisons) et des exemples (illustrations).

\textbf{3. L'explication de texte.} L'explication de texte est le travail qui consiste à dégager le sens d'un texte philosophique : définir les expressions clés, reformuler la thèse sans paraphraser, et montrer les étapes de la démonstration de l'auteur.

\textbf{4. La discussion.} La discussion est la partie de l'exercice où le candidat prend personnellement position sur le problème posé : il évalue la thèse de l'auteur, avance ses propres arguments, examine une objection et conclut de manière justifiée.

\textbf{5. La problématique.} La problématique est la question philosophique à laquelle le texte répond. Elle se formule sous forme interrogative (par exemple : « L'éducation suffit-elle à rendre l'homme libre ? ») et permet de comprendre l'enjeu de la thèse.
\end{corrige}

\begin{corrige}[Exercice 4 — Repérage dans un texte]
\textbf{1. Auteur et œuvre.} L'auteur est Emmanuel Kant (1724--1804), philosophe allemand, et le texte est extrait de l'opuscule \emph{Réponse à la question : qu'est-ce que les Lumières ?} (1784).

\textbf{2. Thème.} Le thème du texte est les Lumières, c'est-à-dire l'émancipation de l'homme par l'usage de sa propre raison.

\textbf{3. Thèse.} Pour Kant, les Lumières consistent en la sortie de l'homme hors de la minorité (la dépendance intellectuelle) dont il est lui-même responsable : il faut oser penser par soi-même.

\textbf{4. Argument.} L'argument est contenu dans la devise \emph{Sapere aude !} (« Aie le courage de te servir de ton propre entendement ») : si la minorité vient non d'un manque d'intelligence mais d'un manque de courage, alors seul un effort personnel de décision permet d'en sortir.
\end{corrige}

\begin{corrige}[Exercice 5 — De la paraphrase à l'explication]
\textbf{1. Classement des reformulations.}
\begin{itemize}
\item a) est une \cle{paraphrase} : elle répète la phrase de Platon avec d'autres mots (« enchaînés dans la grotte » pour « prisonniers de la caverne », « reflets » pour « ombres ») sans en dégager aucun sens philosophique.
\item b) est une véritable \cle{explication philosophique} : elle nomme les concepts en jeu (opinion, apparence, réalité) et formule la thèse de Platon (la connaissance sensible est trompeuse).
\item c) est un \cle{hors-sujet} : elle prend le mot « caverne » au sens propre et zoologique, sans aucun rapport avec la philosophie de Platon.
\end{itemize}

\textbf{2. Explication de « sortir de la caverne » (corrigé type).} Sortir de la caverne, c'est accomplir le difficile passage de l'opinion à la connaissance véritable. Le prisonnier délivré symbolise l'homme qui renonce aux apparences sensibles et aux préjugés reçus pour s'élever, par l'éducation et l'effort de la raison, vers la vérité. Cette ascension est d'abord douloureuse (la lumière éblouit), car remettre en cause ses certitudes exige du courage. Mais elle aboutit à la contemplation du « soleil », c'est-à-dire du Bien, principe de toute vérité : philosopher, c'est donc accepter de se détacher du monde des ombres.
\end{corrige}

\begin{corrige}[Exercice 6 — Construire une réponse de compréhension]
\textbf{Corrigé type (5 à 8 lignes).}

La question demande de dégager la distinction que l'auteur établit entre la technique et la science. Or, l'auteur refuse la conception courante selon laquelle la technique ne serait que l'application de la science : il affirme au contraire que la technique est « aussi ancienne que l'homme ». Son argument repose sur une définition : fabriquer un outil, c'est déjà « transformer la nature selon un projet », c'est-à-dire agir de manière finalisée, ce qui ne suppose aucune théorie scientifique préalable. Ainsi, l'homme préhistorique qui taille un \textbf{silex} pratique une technique sans posséder aucune science. La technique est donc une activité originelle et autonome, et non un produit dérivé de la science.

\textbf{Points de vigilance.} La réponse doit : (i) reformuler la question en introduction ; (ii) citer ou s'appuyer sur des expressions du texte (« aussi ancienne que l'homme », « transformer la nature selon un projet ») ; (iii) ne pas donner d'avis personnel, car il s'agit d'une question de compréhension, non de discussion.
\end{corrige}

\begin{corrige}[Exercice 7 — Choisir des exemples pertinents]
\textbf{1. « L'éducation libère l'homme de l'ignorance. »} Exemple : un jeune de Maroua qui apprend à lire et à calculer au lycée technique peut ensuite vérifier lui-même un contrat ou gérer un budget, au lieu de dépendre de l'avis des autres. Pertinence : l'exemple montre concrètement le passage de la dépendance (ignorance) à l'autonomie (savoir), qui est exactement le contenu de la thèse.

\textbf{2. « La technique peut aliéner le travailleur en le réduisant à un geste répétitif. »} Exemple : l'ouvrier d'une chaîne de mise en bouteille dans une brasserie de Douala, qui répète le même geste des milliers de fois par jour sans voir le produit fini. Pertinence : le travailleur ne fabrique plus un objet dont il serait l'auteur, il « sert un rythme » : c'est précisément l'aliénation décrite par la thèse.

\textbf{3. « Penser par soi-même exige du courage. »} Exemple : un élève de Terminale qui, lors d'un débat de classe ou sur les réseaux sociaux, ose défendre une opinion argumentée contraire à celle de la majorité de ses camarades. Pertinence : l'exemple illustre le \emph{Sapere aude} kantien : l'usage de son propre entendement suppose d'affronter la pression du groupe, donc le courage.
\end{corrige}

\begin{corrige}[Exercice 8 — La copie fautive : repérage d'erreurs méthodologiques]
\begin{enumerate}
\item \textbf{Paraphrase.} « L'homme est responsable de sa minorité dont il doit sortir » répète la phrase de Kant sans l'expliquer. \emph{Correction :} définir les concepts — la minorité est l'incapacité de se servir de son entendement sans la direction d'autrui ; elle est « responsable » parce que sa cause est le manque de courage, non le manque d'intelligence.
\item \textbf{Exemple non pertinent.} Le quartier animé et le football à Yaoundé n'ont aucun rapport avec la thèse de Kant. \emph{Correction :} choisir un exemple qui illustre la dépendance intellectuelle (un citoyen qui répète une rumeur sans la vérifier) ou l'émancipation (un élève qui vérifie une information par lui-même).
\item \textbf{Absence d'argumentation (pétition de principe).} « Kant a raison, c'est évident, tout le monde le sait » affirme sans justifier. \emph{Correction :} prendre position et argumenter : « Kant a raison, car... » suivi d'au moins deux arguments et d'exemples.
\item \textbf{Réponses non numérotées et absence de présentation du texte.} \emph{Correction :} ouvrir la copie par la présentation (auteur, œuvre, thème, problème, thèse) et numéroter chaque réponse selon la numérotation du sujet.
\item \textbf{Hors-sujet et question de discussion non traitée.} Raconter la vie de Kant ne répond à aucune question, et la discussion est escamotée. \emph{Correction :} supprimer la biographie et rédiger une discussion complète : position personnelle, arguments, exemples, objection, réponse, conclusion.
\end{enumerate}
\end{corrige}

\courssub{Sujets type Baccalauréat — Entraînements complets}

\begin{corrige}[Exercice 9 — Sujet type Bac n\textsuperscript{o} 1 : Kant, \emph{Qu'est-ce que les Lumières ?}]
\textbf{Présentation du texte (attendue avant les réponses).} Le texte est extrait de \emph{Réponse à la question : qu'est-ce que les Lumières ?} (1784) d'Emmanuel Kant. Thème : les Lumières. Problème : comment l'homme peut-il sortir de sa minorité intellectuelle ? Thèse : les Lumières sont la sortie de la minorité dont l'homme est lui-même responsable, par le courage de penser par soi-même.

\textbf{Question 1 (4 pts) — Explication de « minorité dont l'homme est lui-même responsable ».} La minorité est l'incapacité de se servir de son entendement sans la direction d'autrui : le mineur intellectuel laisse d'autres penser à sa place (un livre, un guide, l'opinion commune). Elle est « responsable » parce que sa cause ne réside pas dans un défaut d'intelligence, mais dans un défaut de décision et de courage : la paresse (il est plus commode de suivre) et la lâcheté (il est risqué de penser seul). L'homme peut donc en sortir par un acte de volonté.

\textbf{Question 2 (6 pts) — Thèse et étapes de l'argumentation.} Thèse : les Lumières sont la sortie de la minorité par l'usage autonome de la raison. Étapes : (1) définition des Lumières comme sortie de minorité ; (2) définition de la minorité comme dépendance de l'entendement ; (3) diagnostic de la cause : non l'inintelligence, mais le manque de courage ; (4) formulation de la devise \emph{Sapere aude !} ; (5) confirmation par l'analyse des causes psychologiques (paresse et lâcheté) qui maintiennent les hommes « volontiers » mineurs.

\textbf{Question 3 (10 pts) — Discussion : l'éducation suffit-elle à rendre autonome ?} Plan attendu : prise de position claire (par exemple : l'éducation est nécessaire mais non suffisante). Arguments : (i) l'école camerounaise transmet les savoirs et forme l'esprit critique — sans elle, \textbf{pas de point de départ} ; (ii) mais Kant montre que l'obstacle est le courage, non le savoir : un diplômé peut rester mineur s'il suit passivement la famille, les médias ou les réseaux sociaux ; (iii) objection : sans éducation, le courage seul ne suffit pas ; réponse : c'est bien pourquoi l'éducation est nécessaire, mais elle doit être complétée par la décision personnelle. Conclusion nuancée.

\textbf{Barème indicatif.} Présentation du texte : 2 pts ; explication des concepts (minorité, responsabilité) : 4 pts ; thèse et étapes : 6 pts ; discussion (position 2 pts, arguments 4 pts, exemples camerounais 2 pts, conclusion 2 pts).

\textbf{Points de vigilance.} Ne pas paraphraser ; citer la devise \emph{Sapere aude} ; distinguer nettement explication (questions 1--2) et discussion (question 3) ; illustrer par le système éducatif camerounais comme le demande explicitement le sujet.
\end{corrige}

\begin{corrige}[Exercice 10 — Sujet type Bac n\textsuperscript{o} 2 : Platon, l'allégorie de la caverne]
\textbf{Présentation du texte.} Texte adapté de Platon, \emph{La République}, livre VII. Thème : la connaissance et l'ignorance. Problème : comment l'homme passe-t-il de l'opinion trompeuse à la vérité ? Thèse : l'homme ordinaire prend les apparences pour la réalité, et seule une éducation difficile le conduit à la connaissance véritable.

\textbf{Question 1 (5 pts) — Sens philosophique de l'allégorie.} La \cle{caverne} représente le monde sensible, celui des apparences et de l'opinion. Les \cle{chaînes} symbolisent les préjugés, les habitudes et les passions qui empêchent l'homme de voir au-delà du sensible. Les \cle{ombres} figurent les opinions et les images que l'on prend pour la réalité (rumeurs, apparences). La \cle{sortie vers la lumière} représente l'ascension intellectuelle vers la connaissance vraie, progressive et douloureuse ; le soleil symbolise le Bien, principe de toute vérité.

\textbf{Question 2 (5 pts) — Thèse et argumentation.} Thèse : l'homme vit dans l'illusion et \textbf{n'accède} à la vérité que par un effort libérateur. Étapes : (1) description de la situation des prisonniers (l'ignorance) ; (2) identification des ombres à la pseudo-réalité (l'opinion) ; (3) récit de la délivrance et de la souffrance initiale (le coût de l'éducation) ; (4) progression graduelle vers la lumière puis le soleil (les degrés de la connaissance).

\textbf{Question 3 (10 pts) — Discussion : peut-on sortir seul de l'ignorance ?} Position possible : la sortie exige à la fois un maître et un effort personnel. Arguments : (i) dans l'allégorie, le prisonnier est « délivré » et « contraint » de monter : un éducateur (enseignant, parent) est nécessaire pour briser les chaînes — l'école camerounaise joue ce rôle ; (ii) mais personne ne peut voir à la place d'un autre : l'élève doit fournir l'effort de comprendre, comme Kant l'exige avec le \emph{Sapere aude} ; (iii) objection : certains autodidactes réussissent seuls ; réponse : même l'autodidacte s'appuie sur des livres et des maîtres indirects. Exemples : l'élève qui progresse grâce à un professeur exigeant ; le jeune qui vérifie les informations des réseaux sociaux au lieu de les subir.

\textbf{Barème indicatif.} Présentation : 2 pts ; symbolique des quatre éléments : 5 pts ; thèse et étapes : 5 pts ; discussion : 8 pts (position 2, arguments 3, exemples 2, conclusion 1).

\textbf{Points de vigilance.} Ne pas raconter l'allégorie sans l'interpréter ; nommer les concepts (opinion/\emph{doxa}, connaissance, monde sensible) ; prendre clairement position dans la discussion ; ancrer les exemples dans la vie scolaire et sociale camerounaise.
\end{corrige}

\begin{corrige}[Exercice 11 — Sujet type Bac n\textsuperscript{o} 3 : Technique et travail]
\textbf{Partie A — Exercice sur texte.}

\textbf{Question 1 (4 pts) — Technique qui aliène / technique qui libère.} La technique aliène lorsque la machine impose son rythme à l'ouvrier : réduit à répéter « mille fois le même geste », celui-ci ne fabrique plus un objet dont il serait l'auteur, il « sert un rythme » — son travail perd son sens. Mais la même technique libère lorsqu'elle décharge l'homme des tâches pénibles et lui ouvre « le temps de penser ». La distinction ne tient donc pas à la machine elle-même, mais à l'usage social qui en est fait.

\textbf{Question 2 (4 pts) — Thèse et argument principal.} Thèse : la technique n'est ni bonne ni mauvaise en elle-même ; tout dépend de l'usage que la société fait de ses outils. Argument principal : la machine transforme le travail sans le supprimer, et cette transformation peut prendre deux directions opposées (asservissement au rythme ou délivrance des tâches pénibles) ; c'est donc le choix social, et non la technique, qui décide.

\textbf{Question 3 (8 pts) — Discussion.} Position possible : la technique libère ou aliène selon l'organisation du travail. Arguments : (i) aliénation réelle : l'ouvrier de chaîne dans les usines de Douala (mise en bouteille, conditionnement) répète un geste sans maîtriser le produit ; (ii) libération réelle : les machines du port autonome de Douala ou des ateliers modernes suppriment les efforts épuisants et exigent des qualifications nouvelles qui valorisent le travailleur ; (iii) objection : la machine détruit des emplois ; réponse : elle en crée d'autres (maintenance, conduite), à condition que la société organise la formation. Conclusion : la libération dépend d'un usage humain et réfléchi de la technique.

\textbf{Partie B — Entraînement chronométré (corrigé type).}

\textbf{1. Présentation du texte (5 lignes max).} Texte philosophique sur la technique et le travail. Thème : les effets de la machine sur le travailleur. Problème : la technique libère-t-elle ou aliène-t-elle l'homme au travail ? Thèse : la technique transforme le travail et peut aliéner ou libérer selon l'usage que la société fait de ses outils.

\textbf{2. Plan détaillé de la discussion.}
\begin{itemize}
\item \emph{Introduction :} reformulation du problème et annonce de la position (la technique est ambivalente, tout dépend de l'usage social).
\item \emph{Argument 1 :} la technique aliène quand elle réduit l'ouvrier au geste répétitif — exemple : chaînes de production des usines de Douala.
\item \emph{Argument 2 :} la technique libère quand elle supprime la peine et ouvre le temps de penser — exemple : mécanisation portuaire, machines agricoles allégeant le travail des champs.
\item \emph{Objection :} la machine supprime des emplois et aggrave le chômage.
\item \emph{Réponse :} elle déplace les emplois vers la qualification ; le problème est social (formation, organisation), non technique.
\item \emph{Conclusion :} la technique est un instrument ; c'est la société qui doit en décider l'usage au service de l'homme.
\end{itemize}

\textbf{Barème indicatif.} Partie A : 16 pts (4 + 4 + 8) ; Partie B : présentation 2 pts, plan détaillé 2 pts. \textbf{Points de vigilance.} Respecter le temps imparti (1 h 15 / 45 min) ; en Partie B, rédiger \emph{uniquement} ce qui est demandé ; dans le plan, chaque argument doit être accompagné de son exemple ; ne jamais conclure sans avoir répondu à l'objection.
\end{corrige}

\newpage

\coursec{Fiche bilan}

\begin{popfigure}
\begin{tikzpicture}[
    node distance=3.5cm and 2.5cm,
    central/.style={circle, minimum size=2.2cm, draw=popblue, very thick, fill=popblue!80, text=white, font=\bfseries\large, align=center},
    branch/.style={rounded rectangle, minimum height=1.8cm, text width=2.8cm, align=center, font=\small\bfseries, draw, very thick, fill=#1!20, text=#1},
    connect/.style={very thick, draw=#1, -{Stealth[length=3mm]}},
]
\node[central, align=center] (center){Exercice\\sur Texte\\Philosophique};

% Branch 1: Lecture
\node[branch=popblue, above left=of center, align=center] (b1){1. Lecture\\Méthodique};
\draw[connect=popblue] (center) -- (b1);

% Branch 2: Compréhension/Analyse
\node[branch=poporange, above=of center, align=center] (b2){2. Questions\\Compréhension\\ \& Analyse};
\draw[connect=poporange] (center) -- (b2);

% Branch 3: Discussion
\node[branch=popgreen, above right=of center, align=center] (b3){3. Questions\\de Discussion};
\draw[connect=popgreen] (center) -- (b3);

% Branch 4: Application
\node[branch=poppurple, right=of center, align=center] (b4){4. Application\\Notions $\to$\\Concret};
\draw[connect=poppurple] (center) -- (b4);

% Branch 5: Rédaction/Gestion
\node[branch=poppink, below right=of center, align=center] (b5){5. Rédaction,\\Justification\\ \& Gestion Temps};
\draw[connect=poppink] (center) -- (b5);

% Branch 6: Pièges
\node[branch=popdark, below left=of center, align=center] (b6){6. Pièges\\Fréquents\\ \& Check-list};
\draw[connect=popdark] (center) -- (b6);
\end{tikzpicture}
\legende{Carte mentale : les 6 piliers de la maîtrise de l'exercice sur texte (Terminale Technique).}
\end{popfigure}

\begin{bilan}
\textbf{Définitions clés}
\begin{itemize}
    \item \cle{Exercice sur texte} : Épreuve écrite (coeff. souvent fort) consistant en l'étude ordonnée d'un extrait philosophique : compréhension, analyse, discussion.
    \item \cle{Problématique} : La question philosophique fondamentale à laquelle le texte tente de répondre (souvent implicite).
    \item \cle{Thèse} : La réponse affirmative ou négative que l'auteur défend explicitement.
    \item \cle{Argumentation} : L'enchaînement logique (concepts, exemples, distinctions, réfutations) soutenant la thèse.
    \item \cle{Concept} : Notion opératoire définie dans le texte (ex: \emph{liberté}, \cle{conscience}, \cle{devoir}, \cle{État}, \cle{technique}).
\end{itemize}

\textbf{Méthodes express (à mémoriser)}
\begin{center}
\begin{tabularx}{\linewidth}{|l|X|}\hline
\rowcolor{popblueL} \textbf{Étape} & \textbf{Actions Clés} \\\hline
\textbf{Lecture (15-20 min)} & 1. Lecture rapide (sujet global). 2. Lecture crayon en main : repérer thèse, structure, concepts, connecteurs. 3. Formuler la problématique. \\\hline
\textbf{Questions Compréhension} & \textbf{Citer} $\to$ \textbf{Expliquer} (reformuler) $\to$ \textbf{Préciser} (portée, sens technique). Ne pas paraphraser bêtement. \\\hline
\textbf{Questions Analyse} & Identifier l'\textbf{argument} $\to$ Reconstruire le \textbf{raisonnement} (déduction, induction, analogie) $\to$ Évaluer la \textbf{cohérence} interne. \\\hline
\textbf{Question Discussion} & 1. \textbf{Thèse de l'auteur} (résumé). 2. \textbf{Critique/Objets} (limites, contre-exemples, autres auteurs). 3. \textbf{Nuance/Synthèse} (dépassement, condition de validité). 4. \textbf{Conclusion} personnelle argumentée. \\\hline
\textbf{Application Concrète} & Concept $\to$ Situation camerounaise (ex: \emph{corruption}, \emph{éducation}, \emph{terroir}, \emph{modernité}) $\to$ Démontrer la pertinence/limite du concept. \\\hline
\textbf{Rédaction Finale} & Phrases complètes, connecteurs logiques (\emph{donc, toutefois, en effet, par conséquent}), vocabulaire précis, encre noire, lisibilité. \\\hline
\end{tabularx}
\end{center}

\textbf{Structure type d'une réponse de discussion (Modèle « Dialectique Maîtrisée »)}
\begin{enumerate}
    \item \textbf{Reformulation de la thèse} : « L'auteur soutient que... parce que... »
    \item \textbf{Analyse critique} : « Cette thèse est forte car... / Cependant, elle omet... / Un contre-exemple :... »
    \item \textbf{Apport extérieur} : « Auteur X / Notion Y montre que... » (Programme Terminale : Politique, Morale, Science, Technique, Art).
    \item \textbf{Synthèse personnelle} : « En définitive, la thèse est valable à condition que... / sous réserve de... »
\end{enumerate}
\end{bilan}

\begin{aretenir}[Checklist avant le Bac — Exercice sur Texte]
\begin{itemize}
    \item $\square$ Je connais la \textbf{durée} exacte de l'épreuve et mon planning de gestion du temps (Lecture / Brouillon / Propre).
    \item $\square$ Je sais \textbf{repérer la thèse} (souvent dans l'intro ou la conclusion du texte) et la \textbf{problématique} implicite.
    \item $\square$ Je distingue \textbf{comprendre} (reformuler), \textbf{analyser} (expliquer le mécanisme) et \textbf{discuter} (évaluer critiquement).
    \item $\square$ Je maîtrise le \textbf{vocabulaire technique} du programme (liberté, devoir, justice, travail, technique, État, science, conscience, art, religion).
    \item $\square$ Je sais \textbf{citer} précisément (guillemets, numéro de ligne) pour étayer mes réponses.
    \item $\square$ Je sais construire une \textbf{discussion structurée} (Thèse $\to$ Objections $\to$ Nuances $\to$ Position personnelle) et non une simple liste d'opinions.
    \item $\square$ Je suis capable de \textbf{mobiliser un exemple concret camerounais} (social, politique, culturel, économique) pour illustrer un concept.
    \item $\square$ Je vérifie la \textbf{cohérence} de mon propos : pas de contradiction entre l'analyse et la discussion.
    \item $\square$ Je relis pour corriger la \textbf{syntaxe}, l'\textbf{orthographe} et la \textbf{ponctuation} (critères de notation).
    \item $\square$ J'ai révisé les \textbf{pièges classiques} : hors-sujet, paraphrase, absence de thèse personnelle, exemples vagues, confusion analyse/discussion.
    \item $\square$ J'ai mon \textbf{matériel} prêt (stylos noirs, correcteur, montre, texte officiel si autorisé) et je gère mon \textbf{stress} par la respiration.
\end{itemize}
\end{aretenir}

\findecours

\end{document}
